\documentclass[11pt,oneside,openany]{book}
\usepackage[T1]{fontenc}
\usepackage{lmodern}
\usepackage[letterpaper,margin=1.10in,includeheadfoot]{geometry}
\usepackage{amsmath,amssymb,amsthm,mathtools}
\usepackage{microtype,xcolor,titlesec,fancyhdr,needspace}
\usepackage[titles]{tocloft}
\usepackage{longtable,booktabs,array}
\usepackage{hyperref}
\definecolor{ink}{HTML}{172923}
\definecolor{deepgreen}{HTML}{2B5242}
\definecolor{rulegreen}{HTML}{8A9F92}
\hypersetup{colorlinks=true,linkcolor=deepgreen,citecolor=deepgreen,
urlcolor=deepgreen,pdftitle={Lie Algebras and Finite-Dimensional Representations},
pdfauthor={MathGenome},bookmarksnumbered=true}
\ifdefined\pdftrailerid
  \pdftrailerid{MathGenome Lie Algebras and Finite-Dimensional Representations}
\fi
\setcounter{tocdepth}{1}
\setcounter{secnumdepth}{2}
\setlength{\cftbeforechapskip}{0.45em}
\setlength{\headheight}{15pt}
\setlength{\emergencystretch}{2em}
\raggedbottom
\titleformat{\chapter}[display]{\normalfont\color{ink}}
{\color{deepgreen}\large\scshape Chapter \thechapter}{9pt}{\huge\bfseries}
\titlespacing*{\chapter}{0pt}{4pt}{25pt}
\titleformat{\section}{\normalfont\Large\bfseries\color{ink}}{\thesection}{0.7em}{}
\pagestyle{fancy}
\fancyhf{}
\fancyhead[L]{\small\sffamily\color{deepgreen}MathGenome}
\fancyhead[R]{\small\nouppercase{\leftmark}}
\fancyfoot[C]{\small\thepage}
\renewcommand{\headrulewidth}{0.35pt}
\renewcommand{\chaptermark}[1]{\markboth{\thechapter.\ #1}{}}
\fancypagestyle{plain}{\fancyhf{}\fancyfoot[C]{\small\thepage}\renewcommand{\headrulewidth}{0pt}}
\newtheorem{theorem}{Theorem}[chapter]
\newtheorem{lemma}[theorem]{Lemma}
\newtheorem{proposition}[theorem]{Proposition}
\newtheorem{corollary}[theorem]{Corollary}
\theoremstyle{definition}
\newtheorem{definition}[theorem]{Definition}
\newtheorem{example}[theorem]{Example}
\theoremstyle{remark}
\newtheorem{remark}[theorem]{Remark}
\DeclareMathOperator{\ad}{ad}
\DeclareMathOperator{\tr}{tr}
\DeclareMathOperator{\rad}{rad}
\DeclareMathOperator{\nilrad}{nilrad}
\DeclareMathOperator{\Der}{Der}
\DeclareMathOperator{\End}{End}
\DeclareMathOperator{\Hom}{Hom}
\DeclareMathOperator{\Aut}{Aut}
\DeclareMathOperator{\Sym}{Sym}
\DeclareMathOperator{\gr}{gr}
\DeclareMathOperator{\Span}{span}
\DeclareMathOperator{\GL}{GL}
\DeclareMathOperator{\SL}{SL}
\DeclareMathOperator{\Int}{Int}
\DeclareMathOperator{\ch}{ch}
\allowdisplaybreaks[1]

\begin{document}
\frontmatter
\hypersetup{pageanchor=false}
\begin{titlepage}
\thispagestyle{empty}
\thispdfpagelabel{Cover}
\begin{flushleft}
\vspace*{0.65in}
{\large\scshape\color{deepgreen}MathGenome\par}
\vspace{0.35in}
{\color{rulegreen}\rule{0.94\textwidth}{0.65pt}\par}
\vspace{0.43in}
{\color{ink}\fontsize{34}{40}\selectfont\bfseries
Lie Algebras and\\[5pt]Finite-Dimensional\\[5pt]Representations\par}
\vspace{0.35in}
{\color{deepgreen}\Large Structure, Highest Weights, Cohomology, and Real Forms\par}
\vfill
{\large Mathematical lecture notes\par}
\vspace{0.12in}
{\normalsize MathGenome v0.21.0 dependency baseline\par}
{\normalsize November 2026\par}
\vspace*{0.45in}
\end{flushleft}
\end{titlepage}
\hypersetup{pageanchor=true}
\begingroup\small\tableofcontents\endgroup

\chapter*{Scope and conventions}
\addcontentsline{toc}{chapter}{Scope and conventions}
The bracket records an infinitesimal symmetry. Its ideals determine
which parts of the symmetry can be separated, its roots organize the
semisimple part, and its highest weights classify the finite-dimensional
modules. Enveloping algebras connect these constructions to ordinary
modules over associative algebras. Cohomology measures the failure of
extensions and infinitesimal deformations to split.

The first three chapters develop the algebraic tools, including PBW,
the structural criteria, rank one, and complete reducibility. Chapters
4--6 construct the root system and classify the complex semisimple
algebras, with the existence part of the classification treated
separately from uniqueness. Chapters 7--8 construct the modules and
compute their characters, dimensions, multiplicities, and central
characters. Chapter 9 treats extensions, Levi factors and faithful
representations. Chapter 10 studies descent, real forms and the
boundaries of the characteristic-zero statements.

The coefficient field is stated at the start of each argument.
The root and highest-weight chapters use \(\mathbb C\); arbitrary
characteristic-zero fields and real forms are treated with their
additional hypotheses. An alternating bracket means \([x,x]=0\), also
in characteristic two. Modules are left modules, with
\[
 x(yv)-y(xv)=[x,y]v.
\]
For a chosen positive root system,
\(\mathfrak g=\mathfrak n^-\oplus\mathfrak h\oplus\mathfrak n^+\).
The root vector \(e_i\) raises weight and \(f_i\) lowers it.
The highest vector is killed by \(\mathfrak n^+\).
We distinguish the ordinary Weyl action from
\(w\mathbin{\cdot}\lambda=w(\lambda+\rho)-\rho\).
Here \(\rho\) is half the sum of the positive roots.

The invariant bilinear form used for weights is extended complex
bilinearly; on the real span of the roots it is positive definite.
For a direct sum of simple factors, this convention is imposed on
each factor. The coroot satisfies
\(\langle\lambda,\alpha^\vee\rangle
=2(\lambda,\alpha)/(\alpha,\alpha)\).
The normalization of a root vector paired with its opposite is stated
when the Casimir is computed. The trivial algebra, zero modules and
rank-zero root system obey the usual empty-sum and empty-product
conventions.

Finite-dimensional linear algebra, tensor and exterior powers, the
basic theory of associative modules, and integral extensions are used
from the accepted algebra core \cite{p103}. For CE complexes, Ext and
coefficient long exact sequences we also use the detailed constructions
in \cite[Chapter 6]{homological}; the particular maps needed here are
written out. That volume is a mathematical prerequisite, rather than
a claim that the present results have already been formalized.
The result map accompanying this volume records all 86 planned result
IDs, textbook sections and pages, actual statement locations, proof
dependencies and the fixed-version candidate Lean interfaces.

The finite-dimensional inverse function theorem used in the density
and compact-orbit arguments is the one proved in the accepted
multivariable analysis notes \cite{multivariable}.
The needed facts about closed matrix groups and solvable projective
actions are proved here before they are used; no general Lie-group
integration theorem is a hidden prerequisite.

The principal sources are \cite{humphreys,etingof,weibel}.
The matrix and real-form comparisons also use \cite{hall,knapp}.
Broader subjects, such as algebraic groups over general fields,
integral Kostant forms, category \(\mathcal O\), Kac--Moody algebras,
positive-characteristic classification and quantum groups, require
their own additional statements and are not silently identified with
the finite-dimensional characteristic-zero results proved here.

\mainmatter
\chapter{Constructions and enveloping algebras}
\label{ch:constructions}

\section{Brackets, quotients, and representations}

We work over a field \(k\). Statements involving factorials or exponential
series will explicitly assume characteristic zero. All tensor products
in this chapter are over \(k\). An associative algebra is unital, and a
module over it is a unital left module.

\begin{definition}[Lie algebras and their quotients]\label{res:A01}
A Lie algebra is a vector space \(\mathfrak g\) with a bilinear operation
\([x,y]\) satisfying
\[
 [x,x]=0,\qquad
 [x,[y,z]]+[y,[z,x]]+[z,[x,y]]=0.
\]
A linear subspace is a subalgebra if it is closed under brackets; it is
an ideal \(I\) if \([\mathfrak g,I]\subseteq I\). A homomorphism preserves
brackets. For an ideal, the quotient has the bracket
\([x+I,y+I]=[x,y]+I\). If \(f:\mathfrak g\to\mathfrak h\) is a
homomorphism, it induces an isomorphism
\(\mathfrak g/\ker f\longrightarrow\operatorname{im}f\).
\end{definition}

\begin{proof}
Bilinearity and \([x+y,x+y]=0\) imply \([x,y]=-[y,x]\), including in
characteristic two; the alternating identity itself remains part of the
definition. If \(x\) and \(y\) are replaced by \(x+a,y+b\), with
\(a,b\in I\), the bracket changes by
\([a,y]+[x,b]+[a,b]\in I\). Thus the quotient bracket is well defined.
Its alternating and Jacobi identities are images of the identities
in \(\mathfrak g\). The kernel of a homomorphism is an ideal, its image
is a subalgebra, and \(x+\ker f\mapsto f(x)\) is well defined, injective,
surjective onto the image, and bracket preserving.
\end{proof}

An associative algebra \(A\) becomes a Lie algebra under
\([a,b]=ab-ba\). Expansion gives its Jacobi identity: each of the six
ordered triple products occurs once with each sign. In particular,
\(\mathfrak{gl}(V)=\operatorname{End}_k(V)\) is a Lie algebra.

\begin{proposition}[Central and derived series]\label{res:A02}
Set
\[
 D^0\mathfrak g=\mathfrak g,\quad
 D^{r+1}\mathfrak g=[D^r\mathfrak g,D^r\mathfrak g],\qquad
 \gamma _1\mathfrak g=\mathfrak g,\quad
 \gamma _{r+1}\mathfrak g=[\mathfrak g,\gamma _r\mathfrak g].
\]
Here brackets of subspaces mean their linear span. All these subspaces
are ideals. They define solvability and nilpotence by eventual
vanishing, respectively. Subalgebras and quotients of a solvable
(respectively nilpotent) algebra have the same property. An extension
of solvable algebras is solvable.

The upper central series \(Z_0=0\),
\(Z_{r+1}/Z_r=Z(\mathfrak g/Z_r)\) reaches \(\mathfrak g\) if and only
if \(\mathfrak g\) is nilpotent. Moreover,
\[
 [\gamma_i\mathfrak g,\gamma_j\mathfrak g]
 \subseteq\gamma_{i+j}\mathfrak g.
\]
\end{proposition}

\begin{proof}
Jacobi in the form
\([x,[a,b]]=[[x,a],b]+[a,[x,b]]\) shows that the bracket of two
ideals is an ideal; it proves the assertion about both series by
induction. For a homomorphism \(f\), induction gives
\(f(D^r\mathfrak g)=D^r(f\mathfrak g)\) and similarly for \(\gamma_r\).
The analogous inclusions for a subalgebra suffice. If an ideal \(I\)
and the quotient are solvable, some \(D^r\mathfrak g\) lies in \(I\),
and then \(D^{r+s}\mathfrak g\subseteq D^s I\).

For the last inclusion, induct on \(i\). The case \(i=1\) is the
definition. Jacobi gives
\[
 [[\mathfrak g,\gamma_i],\gamma_j]
 \subseteq[\mathfrak g,[\gamma_i,\gamma_j]]
       +[\gamma_i,[\mathfrak g,\gamma_j]]
 \subseteq\gamma_{i+j+1}.
\]
If \(\gamma_{c+1}=0\), induction shows
\(\gamma_{c+1-r}\subseteq Z_r\), starting with
\(\gamma_c\subseteq Z_1\). Thus \(Z_c=\mathfrak g\).
Conversely, \(Z_c=\mathfrak g\) implies
\(\gamma_{r+1}\subseteq Z_{c-r}\), by induction on \(r\), so
\(\gamma_{c+1}=0\).
\end{proof}

Nilpotence implies solvability, since \(D^r\mathfrak g\subseteq
\gamma_{r+1}\mathfrak g\). It does not follow merely from an extension
of two nilpotent algebras: \([h,e]=e\) gives a two-dimensional
counterexample.

\begin{proposition}[Actions, derivations, and semidirect products]
\label{res:A03}
A representation is a linear map
\(\rho:\mathfrak g\to\mathfrak{gl}(M)\) satisfying
\[
 \rho([x,y])=\rho(x)\rho(y)-\rho(y)\rho(x).
\]
The map \(\operatorname{ad}x:y\mapsto[x,y]\) is a representation and
has kernel \(Z(\mathfrak g)\). A derivation of a Lie algebra
\(\mathfrak n\) is a linear map \(d\) with
\(d[a,b]=[da,b]+[a,db]\). Derivations form a Lie subalgebra of
\(\mathfrak{gl}(\mathfrak n)\).

If \(\theta:\mathfrak g\to\operatorname{Der}(\mathfrak n)\) is a
representation, then
\[
 [(x,a),(y,b)]
   =([x,y],\,\theta(x)b-\theta(y)a+[a,b])
 \tag{1.1}\label{eq:semidirect}
\]
makes \(\mathfrak g\oplus\mathfrak n\) a Lie algebra, denoted
\(\mathfrak g\ltimes\mathfrak n\). For an abelian \(\mathfrak n=M\),
the Jacobi identity in this formula is equivalent to the
representation identity.
\end{proposition}

\begin{proof}
Jacobi is exactly the assertion that \(\operatorname{ad}x\) is a
derivation and that \([\operatorname{ad}x,\operatorname{ad}y]
=\operatorname{ad}[x,y]\). Its kernel is the stated center.
Expanding \(de[a,b]-ed[a,b]\) cancels the two mixed terms and proves
that \([d,e]\) is a derivation.

For \eqref{eq:semidirect}, check Jacobi according to the number of
entries in each summand. Three entries from one summand give that
summand's Jacobi identity. Two from \(\mathfrak g\) and one from
\(\mathfrak n\) give
\(\theta([x,y])a-[\theta(x),\theta(y)]a=0\). One from
\(\mathfrak g\) and two from \(\mathfrak n\) give the derivation
identity for \(\theta(x)\). These cases exhaust the expansion.
When \(\mathfrak n\) is abelian, the last condition is automatic,
and the second case is necessary as well as sufficient.
\end{proof}

\section{The universal enveloping algebra}

\begin{theorem}[Construction and universal property]\label{res:A04}
Let \(T(\mathfrak g)=\bigoplus_{r\geq0}\mathfrak g^{\otimes r}\).
The quotient
\[
 U(\mathfrak g)=T(\mathfrak g)/
  \bigl(x\otimes y-y\otimes x-[x,y]\mid x,y\in\mathfrak g\bigr)
\]
has a canonical Lie map \(\iota:\mathfrak g\to U(\mathfrak g)\).
Every Lie map \(\mathfrak g\to A\), with \(A\) associative, extends
uniquely to an algebra map \(U(\mathfrak g)\to A\).
Representations of \(\mathfrak g\) and left \(U(\mathfrak g)\)-modules
are consequently equivalent, with the same linear morphisms.
\end{theorem}

\begin{proof}
The tensor algebra extends a linear map \(f:\mathfrak g\to A\)
uniquely by
\(x_1\otimes\cdots\otimes x_r\mapsto f(x_1)\cdots f(x_r)\).
It annihilates the displayed ideal exactly when
\(f([x,y])=[f(x),f(y)]\). The quotient therefore has the universal
property. Its generators force uniqueness. Taking
\(A=\operatorname{End}_k(M)\) yields the module correspondence.
A linear map commutes with every \(\iota(x)\) if and only if it
commutes with all products and their sums, so the morphisms also
agree. The two constructions leave the underlying action unchanged
and are inverse.
\end{proof}

\begin{theorem}[Poincare--Birkhoff--Witt]\label{res:A05}
Fix a totally ordered basis \((x_i)\) of \(\mathfrak g\). The ordered
monomials
\[
 1,\qquad x_{i_1}\cdots x_{i_r}\quad
 (r\geq1,\ i_1\leq\cdots\leq i_r)
\]
form a basis of \(U(\mathfrak g)\). In particular, \(\iota\) is
injective. For the word-length filtration \(F_rU\),
\[
 S(\mathfrak g)\ \xrightarrow{\ \sim\ }\operatorname{gr}U(\mathfrak g)
\]
is the natural graded algebra isomorphism. In characteristic zero
the symmetrization map
\[
 \operatorname{sym}(y_1\cdots y_r)
 =\frac1{r!}\sum_{\sigma\in S_r}y_{\sigma(1)}\cdots y_{\sigma(r)}
 \tag{1.2}\label{eq:symmetrization}
\]
is a filtered vector-space isomorphism \(S(\mathfrak g)\to U(\mathfrak g)\);
it is equivariant for the adjoint action, but in general is not an
algebra homomorphism.
\end{theorem}

\begin{proof}
In a word, replace a descending adjacent pair \(yx\) by
\(xy+[y,x]\), expressing the bracket in the chosen basis.
Order terms first by word length and then, at fixed length, by
number of inversions. Each replacement produces one word with
one fewer inversion and terms of smaller length. Thus every
reduction terminates.

We verify independence of the first replacement. Disjoint
replacements commute. An overlapping pair can only occur in
a strictly descending triple \(zyx\). Reducing its left pair or
its right pair and then moving the remaining length-three
letters to \(xyz\) leaves a difference of shorter terms
\[
 [y,x]z-z[y,x]+y[z,x]-[z,x]y+[z,y]x-x[z,y].
\]
After reducing each commutator of two basis-linear expressions,
this becomes
\[
 [[y,x],z]+[y,[z,x]]+[[z,y],x]=0.
\]
Induction on length and inversion number therefore gives a unique
normal form. More explicitly, to compare two first reductions,
commute them when disjoint, use this overlap calculation when
adjacent, and normalize all shorter resulting terms by the
induction hypothesis. Applying this comparison successively
compares any two finite reduction sequences.

The normal-form map kills each defining relation with any prefix
and suffix. It consequently kills their linear combinations,
which are exactly the generated two-sided ideal. Conversely,
each reduction changes a word by a defining relation. Thus the
normal forms identify the quotient with the vector space having
the ordered words as basis. In particular, no linear combination
of the length-one words vanishes.

The relation becomes commutativity in top degree, so it defines
the map from \(S(\mathfrak g)\) to \(\operatorname{gr}U\).
Both sides have the same ordered monomial basis in each degree,
which proves bijectivity. In \eqref{eq:symmetrization}, every
permuted product has the same top symbol. Thus symmetrization
induces the identity on the associated graded spaces.
Injectivity and surjectivity now follow by induction on the
finite filtration degree of an element. The formula is
symmetric and multilinear, so is well defined on \(S^r\).
Applying \(\operatorname{ad}x\) to a product differentiates its
factors; averaging commutes with this operation and proves
equivariance.
\end{proof}

This argument works for a free Lie algebra over a commutative
coefficient ring when an ordered module basis is supplied.
The characteristic-zero assumption is needed for
symmetrization, not for the ordered basis theorem. See
\cite[Sections 17.3--17.4]{humphreys} and
\cite[Section 13]{etingof}.

\begin{proposition}[Tensor, dual, and Hopf structures]\label{res:A06}
The enveloping algebra has algebra maps
\[
 \Delta(x)=x\otimes1+1\otimes x,\qquad \epsilon(x)=0
\]
and a linear anti-automorphism \(S\) with \(S(x)=-x\).
They satisfy coassociativity, the counit identities, and
\(m(S\otimes1)\Delta=m(1\otimes S)\Delta=\eta\epsilon\).
For representations \(M,N\),
\[
 x(m\otimes n)=(xm)\otimes n+m\otimes(xn),\qquad
 (x\phi)(m)=-\phi(xm)
 \tag{1.3}\label{eq:tensor-dual-action}
\]
give the tensor and dual representations. On
\(\operatorname{Hom}_k(M,N)\), the corresponding action is
\((x f)(m)=x f(m)-f(xm)\).
\end{proposition}

\begin{proof}
The primitive assignment for \(\Delta\) preserves brackets:
the two tensor factors commute, and the remaining two
commutators are \([x,y]\otimes1+1\otimes[x,y]\).
The universal property constructs \(\Delta\); it also constructs
\(\epsilon\). The assignment \(x\mapsto-x\) is a Lie map to
\(U^{\mathrm{op}}\), so constructs \(S:U\to U^{\mathrm{op}}\).
The two maps \(U\to U^{\otimes3}\) in coassociativity agree on
every generator; the counit identities are checked the same way.
For the antipode identity, the assertion holds on \(1\) and on
each generator. If it holds for \(a,b\), multiplicativity of
\(\Delta\), anti-multiplicativity of \(S\), and the two preceding
identities give it for \(ab\):
\[
 \sum S(a_{(1)}b_{(1)})a_{(2)}b_{(2)}
 =\sum S(b_{(1)})\Big(\sum S(a_{(1)})a_{(2)}\Big)b_{(2)}
 =\epsilon(a)\epsilon(b)1.
\]
Linearity and generation prove it on \(U\).
The first formula in \eqref{eq:tensor-dual-action} is induced
by \(\Delta\), and the second by \(S\).
Alternatively, direct commutator expansion gives the
representation identities and shows that evaluation and
composition are equivariant. Identifying \(M^*\otimes N\)
with \(\operatorname{Hom}_k(M,N)\) when \(M\) is finite
dimensional gives the last formula.
\end{proof}

\section{Free Lie algebras and the BCH product}

\begin{theorem}[The free Lie algebra]\label{res:A07}
For a vector space \(V\) there is a Lie algebra \(L(V)\) and a
linear map \(V\to L(V)\) through which every linear map from \(V\)
to a Lie algebra extends uniquely to a Lie homomorphism.
Moreover,
\[
 U(L(V))\simeq T(V),
\]
and the resulting inclusion identifies \(L(V)\) with the Lie
subalgebra of \(T(V)\) generated by \(V\).
\end{theorem}

\begin{proof}
Take the free nonassociative algebra on \(V\): its degree \(r\)
part is the sum of copies of \(V^{\otimes r}\) indexed by binary
bracketings. Quotient by the ideal generated by alternation
and Jacobi, including all their substitutions and bracket
contexts. Evaluation of a bracketed expression shows the
universal property. No relations on \(V\) are forced, since a
linear map \(V\to V\) with abelian target extends.

For every associative \(A\), a map \(U(L(V))\to A\) is a Lie map
\(L(V)\to A\), hence a linear map \(V\to A\). Construct the
corresponding maps between \(U(L(V))\) and \(T(V)\); both
composites fix their algebra generators \(V\), so are identities.
PBW embeds \(L(V)\) in its enveloping algebra. Its image is
generated under brackets by \(V\), by the nonassociative
construction, proving the last assertion.
\end{proof}

\begin{theorem}[Primitive elements and formal BCH]\label{res:A08}
Suppose \(\operatorname{char}k=0\), and complete \(T(V)\) by tensor
degree. Its elements of positive degree that are primitive for
\(\Delta\) are exactly the degree-completed free Lie algebra
\(\widehat{L(V)}\). Exponentiation and logarithm give mutually
inverse correspondences between primitive elements and
group-like elements with constant term \(1\). Consequently,
\[
 \operatorname{BCH}(X,Y)=\log(\exp X\exp Y)
 =X+Y+\tfrac12[X,Y]
   +\tfrac1{12}[X,[X,Y]]+\tfrac1{12}[Y,[Y,X]]+\cdots
 \tag{1.4}\label{eq:formal-bch}
\]
is a formal Lie series. It satisfies
\(\operatorname{BCH}(\operatorname{BCH}(X,Y),Z)
=\operatorname{BCH}(X,\operatorname{BCH}(Y,Z))\).
\end{theorem}

\begin{proof}
First consider an ordinary, noncompleted enveloping algebra.
Let \(p\) be primitive, and take its highest PBW symbol
\(\bar p\in S^r(L(V))\). The PBW coproduct says that this
symbol is primitive in the symmetric algebra.
In a polynomial algebra over a characteristic-zero field,
the only primitive polynomials are linear: in
\(q(s+t)=q(s)+q(t)\), the coefficient of a single \(t_i\)
is \(\partial q/\partial s_i\), and therefore is constant.
All second derivatives vanish, so \(q\) has degree at most
one. A polynomial uses only finitely many variables, and the
argument applies to it even when the basis is infinite.
Thus \(r\leq1\), so \(p\in k+L(V)\). Comparing constant terms
in the primitive identity removes the scalar.

The tensor grading and the coproduct are compatible. A
primitive element of the completion has each homogeneous
tensor component primitive. Each such component is an
ordinary polynomial in \(T(V)=U(L(V))\), so the preceding
argument puts it in \(L(V)\). Conversely, generators are
primitive, brackets of primitives are primitive, and degreewise
limits preserve the identity. This proves the assertion
about completed primitives.

The series for \(\exp p\) and \(\log(1+a)\) are well defined:
in each tensor degree only finitely many summands contribute.
Their usual inverse identities are formal identities in one
variable. If \(p\) is primitive, then
\[
 \Delta(\exp p)=\exp(p\otimes1+1\otimes p)
              =\exp p\otimes\exp p
\]
because the two summands commute. If \(g\) is group-like,
\(\Delta g=(g\otimes1)(1\otimes g)\), again with commuting
factors, and hence
\(\Delta\log g=\log g\otimes1+1\otimes\log g\).
These computations prove the correspondence.

The product of two group-like elements is group-like.
Therefore its logarithm is a Lie series. Expanding the
associative series through degree three yields the displayed
coefficients in \eqref{eq:formal-bch}. Exponentiating the two
claimed associativity expressions gives
\(\exp X\exp Y\exp Z\) on both sides. Applying the formal
logarithm proves their equality. Substitution is legitimate
because each substituted argument has positive degree.
\end{proof}

\begin{theorem}[Finite BCH on a nilpotent algebra]\label{res:A09}
Let \(\mathfrak n\) be nilpotent of class at most \(c\) over a
characteristic-zero field, so \(\gamma_{c+1}\mathfrak n=0\).
The evaluation of \eqref{eq:formal-bch} is a finite polynomial
product \(x*y\) on \(\mathfrak n\). It is associative, with unit
\(0\) and inverse \(-x\). Lie homomorphisms preserve this
product. If \(D\) is a derivation, then
\[
 D(x*y)=d(\operatorname{BCH})_{(x,y)}(Dx,Dy).
 \tag{1.5}\label{eq:bch-derivation}
\]
\end{theorem}

\begin{proof}
A homogeneous Lie monomial of degree \(r\) is a linear
combination of brackets of length \(r\), which lie in
\(\gamma_r\mathfrak n\) by Proposition~\ref{res:A02}.
Terms of degree greater than \(c\) therefore vanish.
Evaluate the universal formal associativity identity:
all terms of bracket length at most \(c\) form finite sums,
and the others evaluate to zero. This gives associativity.
The formal identities \(\exp0=1\) and
\(\exp x\exp(-x)=1\) give the unit and inverse.

Every Lie homomorphism preserves every finite nested bracket,
so preserves the polynomial. A derivation differentiates a
nested bracket by applying \(D\) to one argument at a time,
by induction on bracket length. Applying this to each term
of the finite BCH polynomial is exactly its polynomial
directional derivative, which proves \eqref{eq:bch-derivation}.
\end{proof}

The construction here is algebraic: it does not use the
integration of an arbitrary Lie algebra into a global Lie
group. That distinction will matter when enlarging a
solvable algebra in the proof of Ado's theorem.

\chapter{Solvability, nilpotence, and semisimplicity}
\label{ch:structure}

\section{Common invariant vectors}

Throughout this chapter the Lie algebras and the spaces acted on
are finite dimensional. Engel's theorem is valid over any field.
Lie's theorem and the structure theorems that follow it use
characteristic zero.

\begin{theorem}[Engel]\label{res:B01}
Let \(0\ne V\) and let \(\mathfrak a\subseteq\mathfrak{gl}(V)\)
be a Lie subalgebra every element of which is a nilpotent
endomorphism. There exists \(0\ne v\in V\) with \(\mathfrak a v=0\).
There is a basis in which all elements of \(\mathfrak a\) are
strictly upper triangular.

A finite-dimensional Lie algebra \(\mathfrak g\) is nilpotent
if and only if \(\operatorname{ad}x\) is nilpotent for every
\(x\in\mathfrak g\).
\end{theorem}

\begin{proof}
Induct on \(\dim\mathfrak a\) to prove the common-kernel
assertion for every representation by nilpotent operators.
The zero algebra is immediate. Choose a maximal proper
subalgebra \(\mathfrak b\) of a nonzero \(\mathfrak a\).
For \(b\in\mathfrak b\), nilpotence of \(b\) on \(V\)
implies nilpotence of \(\operatorname{ad}b\) on
\(\operatorname{End}(V)\): if \(b^N=0\), expand
\((L_b-R_b)^{2N-1}\) and use commutativity of \(L_b,R_b\).
Its induced action on \(\mathfrak a/\mathfrak b\) is
therefore nilpotent. The acting image has dimension at most
\(\dim\mathfrak b<\dim\mathfrak a\), so induction gives
an \(x\notin\mathfrak b\) with
\([\mathfrak b,x]\subseteq\mathfrak b\).
The subspace \(\mathfrak b+kx\) is a subalgebra.
Maximality implies it is all of \(\mathfrak a\); hence
\(\mathfrak b\) is an ideal of codimension one.

Induction on \(\mathfrak b\) acting on \(V\) gives
\(W=\{v:\mathfrak b v=0\}\ne0\).
For \(v\in W\) and \(b\in\mathfrak b\),
\(b(xv)=x(bv)+[b,x]v=0\), so \(W\) is \(x\)-stable.
The nilpotent operator \(x|_W\) has a nonzero kernel.
A vector in that kernel is killed by all of \(\mathfrak a\).
Applying the assertion on successive quotients by the
resulting one-dimensional invariant subspaces gives a full
flag on which every diagonal action is zero. This is the
strict triangularization.

If \(\gamma_{c+1}\mathfrak g=0\), each
\((\operatorname{ad}x)^c\) is zero. Conversely, apply strict
triangularization to the adjoint image on \(\mathfrak g\).
The product of \(\dim\mathfrak g\) strictly upper triangular
operators is zero. Such products span the iterated brackets
defining the lower central series, so that series terminates.
\end{proof}

\begin{lemma}[The cyclic-space calculation]\label{supp:cyclic-weight}
Assume characteristic zero. Let \(\mathfrak b\) be an ideal
of \(\mathfrak b+kx\), and suppose \(0\ne v\in V\) satisfies
\(bv=\lambda(b)v\) for all \(b\in\mathfrak b\).
The space \(W=k[x]v\) is stable under both \(x\) and
\(\mathfrak b\), and each \(b\) has the constant diagonal
\(\lambda(b)\) on \(W\). Moreover,
\(\lambda([x,\mathfrak b])=0\), and the entire common
\(\lambda\)-eigenspace of \(\mathfrak b\) is \(x\)-stable.
\end{lemma}

\begin{proof}
For every \(b\), induction on \(j\), using
\(bx^jv=x(bx^{j-1}v)+[b,x]x^{j-1}v\), gives
\[
 bx^jv-\lambda(b)x^jv
       \in\operatorname{span}(v,xv,\ldots,x^{j-1}v).
\]
The induction is simultaneous for all \(b\), since
\([b,x]\in\mathfrak b\).
Let \(m\) be the first integer for which \(x^mv\)
depends on earlier powers. The first \(m\) powers form
a basis of \(W\), proving stability and the diagonal claim.
Since the trace of a commutator on the invariant space \(W\)
is zero,
\[
 0=\operatorname{tr}_W[x,b]=m\lambda([x,b]).
\]
Characteristic zero gives the asserted vanishing.
If \(w\) is any common \(\lambda\)-eigenvector, then
\(b(xw)=\lambda(b)xw+[b,x]w=\lambda(b)xw\).
This proves stability of the full eigenspace.
\end{proof}

\begin{theorem}[Lie's theorem]\label{res:B02}
Let \(k\) be algebraically closed of characteristic zero.
A nonzero representation of a solvable Lie algebra
\(\mathfrak a\) has a common eigenvector. Every such
representation admits a common upper triangular basis.
In particular, its irreducible representations are
one dimensional.
\end{theorem}

\begin{proof}
Induct on \(\dim\mathfrak a\). For a nonzero solvable
algebra, \([\mathfrak a,\mathfrak a]\ne\mathfrak a\);
a hyperplane \(\mathfrak b\) containing the derived algebra
is an ideal. Write \(\mathfrak a=\mathfrak b+kx\).
Induction gives a common \(\mathfrak b\)-eigenvector and
its character \(\lambda\).
By Lemma~\ref{supp:cyclic-weight}, the common eigenspace
\(V_\lambda\ne0\) is \(x\)-stable. Algebraic closure
supplies an eigenvector of \(x|_{V_\lambda}\).
It is a common eigenvector of \(\mathfrak a\).
Now induct on \(\dim V\), applying the assertion to the
quotient by its eigenline. Lifting the quotient flag
produces a common triangular basis. An irreducible
module must equal its invariant eigenline.
\end{proof}

One consequence needed later is worth proving now. If
\(\mathfrak r\) is a solvable ideal of an arbitrary
\(\mathfrak g\), and \(V\) is an irreducible
\(\mathfrak g\)-module over an algebraically closed
characteristic-zero field, then \(\mathfrak r\) acts by
scalars and \([\mathfrak g,\mathfrak r]\) acts by zero.
Indeed, choose a common \(\mathfrak r\)-eigenvector.
For every \(x\in\mathfrak g\), the cyclic-space lemma for
\(\mathfrak r+kx\) proves
\(\lambda([x,\mathfrak r])=0\) and stability of the same
nonzero eigenspace. It is consequently a
\(\mathfrak g\)-submodule, and therefore is all of \(V\).

\section{Jordan decomposition and the trace criterion}

\begin{proposition}[Jordan decomposition and derivations]
\label{res:B03}
Over a characteristic-zero field every \(A\in\operatorname{End}(V)\)
has a unique decomposition
\[
 A=A_s+A_n,\qquad [A_s,A_n]=0,
\]
where \(A_s\) is semisimple after scalar extension and
\(A_n\) is nilpotent. Both are polynomials in \(A\).
Every \(A\)-stable subspace is stable under both parts.
If \(A=D\) is a derivation of a Lie algebra, then both
\(D_s,D_n\) are derivations.
\end{proposition}

\begin{proof}
Over a splitting field, decompose \(V=\bigoplus V_\lambda\)
into generalized eigenspaces. Choose, by the Chinese
remainder theorem, \(P(t)\) with
\(P(t)\equiv\lambda\pmod{(t-\lambda)^N}\) for every
eigenvalue, where \(N\) is large enough.
Then \(P(A)\) is \(\lambda\) times the identity on
\(V_\lambda\), and \(A-P(A)\) is nilpotent there.
The polynomial can be uniquely chosen modulo the minimal
polynomial of \(A\); its interpolation data are Galois
invariant. In characteristic zero the splitting field is
separable, so this representative has coefficients in \(k\).

For uniqueness, a second commuting decomposition commutes
with \(A\), hence with its polynomial parts. The difference
of the two semisimple parts then equals the opposite
difference of the nilpotent parts. The former is
semisimple (simultaneously diagonalize the commuting
semisimple operators), while the latter is nilpotent
(expand a sufficiently high power of commuting nilpotents).
Both differences are therefore zero. Polynomial expressions
also prove the assertion about invariant subspaces.

Suppose \(D\) is a derivation. For generalized eigenvectors
\(u\in V_\lambda,v\in V_\mu\), the derivation rule gives
\[
 (D-\lambda-\mu)^N[u,v]
 =\sum_{i+j=N}\binom Ni[(D-\lambda)^iu,(D-\mu)^jv].
\]
For large \(N\) every term is zero, so
\([V_\lambda,V_\mu]\subseteq V_{\lambda+\mu}\).
Consequently
\[
 D_s[u,v]=(\lambda+\mu)[u,v]
          =[D_su,v]+[u,D_sv].
\]
By bilinearity this holds everywhere. Thus \(D_s\) is a
derivation, and so is \(D_n=D-D_s\). The identities descend
from the splitting field.
\end{proof}

\begin{theorem}[Cartan's trace criterion]\label{res:B04}
For a Lie subalgebra \(\mathfrak a\subseteq\mathfrak{gl}(V)\)
over a characteristic-zero field, the following are equivalent:
\[
 \mathfrak a\text{ is solvable},\qquad
 \operatorname{tr}(xy)=0
 \quad(x\in[\mathfrak a,\mathfrak a],\ y\in\mathfrak a).
 \tag{2.1}\label{eq:trace-criterion}
\]
In particular, with
\(K_{\mathfrak g}(x,y)=
\operatorname{tr}_{\mathfrak g}(\operatorname{ad}x\operatorname{ad}y)\),
an abstract algebra is solvable if and only if
\(K_{\mathfrak g}([\mathfrak g,\mathfrak g],\mathfrak g)=0\).
\end{theorem}

\begin{proof}
Extension to an algebraic closure preserves traces,
brackets, and vanishing of a term of the derived series.
We may therefore assume \(k\) algebraically closed.
For the forward implication apply Lie's theorem:
derived matrices are strictly upper triangular and the
other matrices upper triangular, so their product has zero
diagonal.

Conversely, fix \(x\in[\mathfrak a,\mathfrak a]\), with
distinct eigenvalues \(\lambda_1,\ldots,\lambda_m\) and
generalized eigenspace dimensions \(d_1,\ldots,d_m\).
Let \(E\) be their finite-dimensional rational span in \(k\)
and choose any rational linear map \(b:E\to\mathbb Q\).
Let \(B\) act by \(b(\lambda_i)\) on the generalized
\(\lambda_i\)-eigenspace of \(x\). We claim
\[
 \operatorname{ad}B=R(\operatorname{ad}x)
 \quad\text{for a polynomial }R\text{ with }R(0)=0.
 \tag{2.2}\label{eq:cartan-polynomial}
\]
Indeed, \(\operatorname{ad}x_s\) has eigenvalues
\(\lambda_i-\lambda_j\) on the corresponding Hom blocks.
Interpolation gives a polynomial \(Q\) with
\(Q(\lambda_i-\lambda_j)=b(\lambda_i)-b(\lambda_j)\)
and \(Q(0)=0\). Thus
\(\operatorname{ad}B=Q(\operatorname{ad}x_s)\).
The identity
\((\operatorname{ad}x)_s=\operatorname{ad}x_s\)
follows by separating the commuting diagonalizable and
nilpotent operators
\(\operatorname{ad}x_s,\operatorname{ad}x_n\).
By Jordan decomposition, the former is \(P(\operatorname{ad}x)\),
where \(P(0)=0\), since zero is an eigenvalue on the identity
endomorphism. Set \(R=Q\circ P\) to obtain
\eqref{eq:cartan-polynomial}.

Write \(x=\sum_j[y_j,z_j]\) in \(\mathfrak a\).
Because \(R\) has no constant term,
\([B,y_j]=R(\operatorname{ad}x)y_j\) belongs to
\([\mathfrak a,\mathfrak a]\). Cyclicity of trace and
\eqref{eq:trace-criterion} give
\[
 \sum_i d_i b(\lambda_i)\lambda_i
 =\operatorname{tr}(Bx)
 =\sum_j\operatorname{tr}([B,y_j]z_j)=0.
\]
Apply the rational linear functional \(b\) to this
identity in \(E\). It yields
\(\sum_i d_i b(\lambda_i)^2=0\) in \(\mathbb Q\).
Each \(d_i\) is a positive integer, so all
\(b(\lambda_i)=0\). As every rational linear functional
has this property, \(E=0\). Hence every eigenvalue of
\(x\) is zero and \(x\) is nilpotent.

Engel's theorem now makes \([\mathfrak a,\mathfrak a]\)
nilpotent, hence solvable. Its abelian quotient proves
solvability of \(\mathfrak a\).
For the abstract assertion apply the matrix criterion
to the adjoint image \(\mathfrak g/Z(\mathfrak g)\);
the missing kernel is central and therefore solvable.
\end{proof}

\section{The radical and semisimple structure}

\begin{theorem}[The radical and the Killing criterion]\label{res:B05}
The sum of all solvable ideals is a solvable ideal
\(\operatorname{rad}\mathfrak g\), and
\(\mathfrak g/\operatorname{rad}\mathfrak g\) has no nonzero
solvable ideal. In characteristic zero,
\(\operatorname{rad}\mathfrak g=0\) if and only if the Killing
form is nondegenerate. Such an algebra is called semisimple.
\end{theorem}

\begin{proof}
For solvable ideals \(I,J\), the quotient
\((I+J)/I\simeq J/(I\cap J)\) is solvable, so \(I+J\)
is solvable. A finite-dimensional sum is already a finite
sum, proving the first assertion. A solvable ideal in the
quotient lifts to a solvable extension of the radical,
so must be zero.

Trace cyclicity gives symmetry and invariance of the Killing
form:
\[
 K([x,y],z)=K(x,[y,z]).
 \tag{2.3}\label{eq:killing-invariant}
\]
Its kernel \(I\) is an ideal. For any ideal \(J\) and
\(x,y\in J\), the action of \(x,y\) on \(\mathfrak g/J\)
is zero, so a block triangular trace calculation gives
\[
 K_J(x,y)=K_{\mathfrak g}(x,y).
 \tag{2.4}\label{eq:killing-ideal}
\]
If \(\mathfrak g\) is semisimple, \(K_I=0\).
The trace criterion makes \(I\) solvable; hence \(I=0\).

Conversely, suppose \(K_{\mathfrak g}\) is nondegenerate
and the radical is nonzero. Its last nonzero derived term
is a nonzero abelian ideal \(A\) of \(\mathfrak g\).
For \(a\in A,y\in\mathfrak g\),
\(\operatorname{ad}a\operatorname{ad}y\) sends
\(\mathfrak g\) into \(A\) and kills \(A\):
\(\operatorname{ad}y\) preserves \(A\), and
\([a,A]=0\). Its square is zero, so its trace is zero.
Thus \(A\) lies in the Killing kernel, a contradiction.
\end{proof}

This also proves that scalar extension preserves
semisimplicity: its Killing matrix is the same matrix
over the larger field and remains invertible. We shall
use this consequence when passing from real algebras
to their complexifications.

\begin{theorem}[Simple-ideal decomposition]\label{res:B06}
A finite-dimensional characteristic-zero semisimple algebra
is a direct sum of nonabelian simple ideals. Every ideal
has an ideal complement, namely its Killing orthogonal
complement. Its center is zero and
\([\mathfrak g,\mathfrak g]=\mathfrak g\).
Every ideal is a sum of a subset of its simple ideals.
\end{theorem}

\begin{proof}
If \(I\) is an ideal, invariance shows that \(I^\perp\)
is an ideal. Set \(J=I\cap I^\perp\). By
\eqref{eq:killing-ideal}, its Killing form is zero,
since \(J\subseteq I^\perp\) and \(J\subseteq I\).
Thus \(J\) is solvable, and semisimplicity forces \(J=0\).
Nondegeneracy and dimensions then give
\(\mathfrak g=I\oplus I^\perp\).
The two ideals commute because their bracket lies in
their zero intersection.

Take a minimal nonzero ideal and split it off.
An ideal of this summand is an ideal of the full direct
sum, so minimality makes the summand simple; it cannot
be abelian. Iterating decreases dimension and gives the
decomposition. A simple nonabelian algebra is perfect,
since its derived algebra is a nonzero ideal, and its
center is zero. These assertions pass to direct sums.

Write \(\mathfrak g=\bigoplus_i\mathfrak g_i\) and let
\(I\) be any ideal. Its projection to \(\mathfrak g_i\)
is either zero or all of that simple algebra. In the
latter case
\([\mathfrak g_i,I]=
[\mathfrak g_i,\operatorname{pr}_i I]=\mathfrak g_i\),
so \(\mathfrak g_i\subseteq I\). Removing these
summands leaves zero. This proves the final assertion
and uniqueness of the complementary sum.
\end{proof}

\begin{theorem}[Inner derivations and abstract Jordan decomposition]
\label{res:B07}
Every derivation of a characteristic-zero semisimple
\(\mathfrak g\) is inner. Each \(x\in\mathfrak g\) has a
unique expression
\[
 x=x_s+x_n,\qquad [x_s,x_n]=0,
\]
where \(\operatorname{ad}x_s\) is semisimple and
\(\operatorname{ad}x_n\) is nilpotent. If \([x,y]=0\),
then \([x_s,y]=[x_n,y]=0\).
\end{theorem}

\begin{proof}
Inside \(\operatorname{Der}\mathfrak g\), the adjoint
image is an ideal, because
\([D,\operatorname{ad}x]=\operatorname{ad}(Dx)\).
On derivations use the invariant trace form
\(B(D,E)=\operatorname{tr}_{\mathfrak g}(DE)\).
Its restriction to \(\operatorname{ad}\mathfrak g\)
is the nondegenerate Killing form. Hence
\[
 \operatorname{Der}\mathfrak g
 =\operatorname{ad}\mathfrak g\oplus
   (\operatorname{ad}\mathfrak g)^\perp.
\]
The second summand is an ideal by invariance.
It commutes with the first, since their intersection
is zero. For \(D\) in this complement,
\(\operatorname{ad}(Dx)=[D,\operatorname{ad}x]=0\).
The center is zero, so \(Dx=0\) for all \(x\), and
the complement vanishes.

Apply Proposition~\ref{res:B03} to
\(\operatorname{ad}x\) on \(\mathfrak g\).
Both matrix Jordan parts are derivations, so are
\(\operatorname{ad}x_s,\operatorname{ad}x_n\)
for uniquely determined elements. Their sum and
commutator identify \(x=x_s+x_n\) and
\([x_s,x_n]=0\), using injectivity of the adjoint
map. Matrix uniqueness gives uniqueness here.
If \([x,y]=0\), then \(\operatorname{ad}y\)
commutes with \(\operatorname{ad}x\) and its
polynomial parts. Hence the two asserted brackets
are central, and therefore zero.
\end{proof}

\begin{theorem}[Characteristic radical and nilradical]\label{res:B08}
In characteristic zero the radical and the largest nilpotent
ideal \(\operatorname{nilrad}\mathfrak g\) are preserved by
every derivation. More precisely,
\[
 D(\operatorname{rad}\mathfrak g)
       \subseteq\operatorname{nilrad}\mathfrak g
 \quad(D\in\operatorname{Der}\mathfrak g),\qquad
 [\mathfrak g,\operatorname{rad}\mathfrak g]
       \subseteq\operatorname{nilrad}\mathfrak g.
 \tag{2.5}\label{eq:nilrad-characteristic}
\]
The nilradical consists exactly of those \(x\) in the
radical for which \(\operatorname{ad}_{\mathfrak g}x\)
is nilpotent.
\end{theorem}

\begin{proof}
We first establish the useful equality
\[
 \operatorname{rad}\mathfrak g
     =[\mathfrak g,\mathfrak g]^\perp
 \tag{2.6}\label{eq:rad-perp}
\]
under the Killing form. Over an algebraic closure,
choose a composition series of the adjoint module.
The consequence of Lie's theorem proved after
Theorem~\ref{res:B02} says that the radical acts
by a scalar on each simple quotient. A commutator
has zero trace on each quotient. Thus
\(K(\operatorname{rad}\mathfrak g,
[\mathfrak g,\mathfrak g])=0\).
For the reverse containment let
\(J=[\mathfrak g,\mathfrak g]^\perp\).
It is an ideal by invariance. If \(a,b,c\in J\),
then \([a,b]\in[\mathfrak g,\mathfrak g]\), so
\(K([a,b],c)=0\). By \eqref{eq:killing-ideal},
the trace criterion makes \(J\) solvable.
It therefore lies in the radical.
The same argument proves the equality over \(k\):
one inclusion is checked after extension, and for
the other \(J\) is already defined over \(k\).
In particular this equality shows that the radical
commutes with scalar extension.

For a derivation \(D\), the identity
\([D,\operatorname{ad}x]=\operatorname{ad}(Dx)\)
and trace cyclicity give
\[
 K(Dx,y)+K(x,Dy)=
 \operatorname{tr}[D,\operatorname{ad}x\operatorname{ad}y]=0.
\]
Since \(D\) preserves the derived algebra,
\eqref{eq:rad-perp} proves \(D\mathfrak r\subseteq\mathfrak r\),
where \(\mathfrak r=\operatorname{rad}\mathfrak g\).

Over an algebraic closure, triangularize the action of
\(\mathfrak r\) on \(\mathfrak g\). Let \(\mathfrak n\)
be the subspace on which all its diagonal characters
vanish. Equivalently,
\[
 \mathfrak n=\{x\in\mathfrak r:
             \operatorname{ad}_{\mathfrak g}x
             \text{ is nilpotent}\}.
\]
For any derivation \(D\), the semidirect algebra
\(kD\ltimes\mathfrak r\) is solvable and acts on
\(\mathfrak g\) through \(D\) and
\(\operatorname{ad}\mathfrak r\). Lie's theorem
triangularizes this whole action. Its commutators,
in particular
\([D,\operatorname{ad}x]=\operatorname{ad}(Dx)\)
for \(x\in\mathfrak r\), are strictly upper
triangular. Thus \(D\mathfrak r\subseteq\mathfrak n\).
Taking \(D=\operatorname{ad}y\) shows that
\(\mathfrak n\) is an ideal. Engel applied to its
adjoint action on itself makes it nilpotent.

If \(A\) is any nilpotent ideal, it is solvable,
so lies in \(\mathfrak r\). For \(a\in A\),
\(\operatorname{ad}a\) sends \(\mathfrak g\) into
\(A\), and its further powers send \(A\) down
the lower central series of \(A\).
It is nilpotent on \(\mathfrak g\), so \(A\subseteq
\mathfrak n\). This proves maximality.
Finally \(D\mathfrak n\subseteq D\mathfrak r
\subseteq\mathfrak n\).

Over a nonclosed field, define \(\mathfrak n\) by
the same nilpotent-adjoint condition. The preceding
triangular calculation after scalar extension proves
that its \(k\)-rational elements form a subspace and
an ideal, and proves all the stated derivation
inclusions there. Engel and the maximality argument
are already valid over \(k\). Hence the conclusion
descends without choosing eigenvalues in \(k\).
\end{proof}

The trace criterion, rather than a representation
classification, supplied every semisimple fact used
in this chapter. The complete reducibility of arbitrary
finite-dimensional representations is a separate
theorem, proved next.

\chapter{Rank one, Casimir, and complete reducibility}
\label{ch:rank-one}

\section{The rank-one calculation}

Write \(\mathfrak{sl}_2(k)\) with basis
\[
 e=\begin{pmatrix}0&1\\0&0\end{pmatrix},\qquad
 f=\begin{pmatrix}0&0\\1&0\end{pmatrix},\qquad
 h=\begin{pmatrix}1&0\\0&-1\end{pmatrix}.
\]
Thus \([h,e]=2e,[h,f]=-2f,[e,f]=h\). In characteristic
zero this algebra is simple: an ideal is stable under
\(\operatorname{ad}h\), whose three distinct eigenspaces
are \(ke,kf,kh\). Polynomial projections onto these
eigenspaces show that a nonzero ideal contains one of
the basis vectors. Its brackets with the others then
give all three.

\begin{theorem}[Simple \(\mathfrak{sl}_2\)-modules]\label{res:C01}
Over an algebraically closed characteristic-zero field
the nonzero finite-dimensional simple
\(\mathfrak{sl}_2\)-modules are exactly \(L_n\),
\(n\in\mathbb Z_{\geq0}\). They have basis
\(v_0,\ldots,v_n\) on which
\[
 hv_j=(n-2j)v_j,\quad
 fv_j=v_{j+1},\quad
 ev_j=j(n-j+1)v_{j-1},
 \tag{3.1}\label{eq:sl2-basis}
\]
where \(v_{-1}=v_{n+1}=0\). They can be realized on
the degree-\(n\) homogeneous polynomials in \(X,Y\) by
\[
 e=X\partial_Y,\qquad f=Y\partial_X,\qquad
 h=X\partial_X-Y\partial_Y.
\]
\end{theorem}

\begin{proof}
Choose an \(h\)-eigenvector in a simple module. The
relation with \(e\) shows that every nonzero successive
\(e\)-image has its eigenvalue increased by \(2\).
There are only finitely many eigenvalues, so the last
nonzero image is a vector \(v\) with
\(ev=0,hv=\lambda v\).
Commuting \(h\) and \(e\) past powers of \(f\) gives
\[
 hf^jv=(\lambda-2j)f^jv,\qquad
 ef^jv=j(\lambda-j+1)f^{j-1}v.
 \tag{3.2}\label{eq:sl2-ladder}
\]
For the second formula, use
\([e,f^j]=\sum_{a=0}^{j-1}f^a h f^{j-1-a}\);
on \(v\) its coefficient is
\(\sum_{a=0}^{j-1}(\lambda-2(j-1-a))
=j(\lambda-j+1)\).

PBW with the order \(f,h,e\) says that the vectors
\(f^jv\) span the generated submodule, hence the whole
simple module. Their distinct \(h\)-eigenvalues make
the nonzero ones independent. Let \(n\) be the final
index with \(f^nv\ne0\). Applying \(e\) to
\(f^{n+1}v=0\) gives
\((n+1)(\lambda-n)f^nv=0\), so \(\lambda=n\).
This proves \eqref{eq:sl2-basis} and uniqueness.

Conversely the three polynomial operators satisfy
the displayed brackets. Set
\[
 v_j=f^jX^n=\frac{n!}{(n-j)!}X^{n-j}Y^j.
\]
Direct differentiation gives
\eqref{eq:sl2-basis}. A nonzero submodule contains an
\(h\)-eigenvector: the distinct-eigenvalue projections
are polynomials in \(h\). Each eigenline is \(kv_j\);
applying \(e^j\), whose coefficient is nonzero,
gives \(v_0\), and applying \(f\) gives every \(v_j\).
Thus the module is simple. This argument includes
\(n=0\), the trivial one-dimensional module.
\end{proof}

\section{Central operators and cohomology}

We now return to an arbitrary characteristic-zero field
and a finite-dimensional semisimple Lie algebra
\(\mathfrak g\). Its modules in this section are
finite dimensional. No root decomposition or highest
weight classification is needed for complete
reducibility.

\begin{theorem}[The module-dependent Casimir]\label{res:C02}
Let \(M\) be a nontrivial simple \(\mathfrak g\)-module.
There is a central element \(C_M\in U(\mathfrak g)\)
which annihilates the trivial module and acts
invertibly on \(M\).

More explicitly, let
\(\mathfrak h=\rho(\mathfrak g)\subseteq\operatorname{End}_k(M)\).
The trace form
\(B_M(x,y)=\operatorname{tr}_M(xy)\) is nondegenerate
on \(\mathfrak h\). Lift \(\mathfrak h\) to the
complementary ideal of \(\ker\rho\) in \(\mathfrak g\).
For a basis \(a_i\) and its \(B_M\)-dual basis \(a^i\)
of that ideal, one may take
\[
 C_M=\sum_i a_i a^i.
 \tag{3.3}\label{eq:module-casimir}
\]
Its trace on \(M\) is \(\dim\mathfrak h\).
If \(k\) is algebraically closed, it acts as the
nonzero scalar \(\dim\mathfrak h/\dim M\).
\end{theorem}

\begin{proof}
The image \(\mathfrak h\) is semisimple by
Theorem~\ref{res:B06}. The radical \(I\) of its
invariant trace form is an ideal. On \(I\), the
trace pairing in the faithful representation on
\(M\) is zero. In particular it vanishes on
\([I,I]\times I\), so the matrix trace criterion
shows that \(I\) is solvable. Semisimplicity of
\(\mathfrak h\) gives \(I=0\).

The element \(\sum_i a_i\otimes a^i\) is independent
of the dual basis: under the identification of one
factor with its dual by \(B_M\), it corresponds to
the identity endomorphism. Invariance of \(B_M\)
makes this tensor invariant under the diagonal
adjoint action. Concretely, if
\([x,a_i]=\sum_j b_{ji}a_j\), then
\([x,a^j]=-\sum_i b_{ji}a^i\), and the two sums
in its derivative cancel. Multiplication of the
tensor gives
\([x,C_M]=0\).
The other ideal \(\ker\rho\) commutes with its
complement, so \(C_M\) is central in the entire
\(U(\mathfrak g)\).

Each \(a_i\) kills the trivial module, hence so
does \(C_M\). On \(M\),
\[
 \operatorname{tr}_M C_M
 =\sum_i B_M(a_i,a^i)=\dim\mathfrak h\ne0.
 \tag{3.4}\label{eq:casimir-trace}
\]
Thus its action is a nonzero module endomorphism.
Its kernel and image are submodules of the simple
module \(M\), so it is invertible. Over an
algebraically closed field Schur's lemma makes
it scalar, and \eqref{eq:casimir-trace} gives the
value. Over a general field invertibility, rather
than a scalar formula, is all that is asserted.
\end{proof}

\begin{theorem}[Casimir contraction of nontrivial simple coefficients]
\label{res:C03}
For a nontrivial finite-dimensional simple
\(\mathfrak g\)-module \(M\),
\[
 H^q(\mathfrak g,M)=0\qquad(q\geq0).
\]
The vanishing is induced by a nullhomotopy for the
central Casimir action on a free resolution of the
trivial module.
\end{theorem}

\begin{proof}
Use the Chevalley--Eilenberg free resolution
\[
 P_q=U(\mathfrak g)\otimes\bigwedge^q\mathfrak g
       \longrightarrow k.
 \tag{3.5}\label{eq:ce-resolution}
\]
Its differential, freeness, and exactness are
proved in \cite[Chapter 6, A06--A07]{homological}:
the PBW filtration identifies the associated graded
resolution with the Koszul resolution, and the
finite-degree lifting argument proves exactness.
Thus its use here has a previously established
proof dependency.

Multiplication by \(C_M\) is a \(U(\mathfrak g)\)-linear
chain endomorphism of \(P_\bullet\), because it is
central; it induces zero on \(k\).
We give the comparison homotopy explicitly at the
lifting level. Put \(h_{-1}=0\). Suppose \(h_i\)
has been chosen below degree \(q\) so that
\[
 C_M=d h_i+h_{i-1}d
\]
there. The map
\(C_M-h_{q-1}d:P_q\to P_q\) has image in
\(\ker d\): applying \(d\), using the already
constructed equation in degree \(q-1\), and using
\(d^2=0\), gives zero. At degree zero the same
statement uses the augmentation. Exactness of
\eqref{eq:ce-resolution} and projectivity of the
free \(P_q\) lift this map through
\(d:P_{q+1}\to\ker d\). Choose that lift as \(h_q\).
Induction constructs a homotopy in all degrees.

Apply \(\operatorname{Hom}_{U(\mathfrak g)}(-,M)\).
For a cochain \(f:P_q\to M\), define
\(Hf=f h_{q-1}\).
Then
\[
 (\delta H+H\delta)f=f C_M=C_M f,
\]
the final equality following from \(U\)-linearity.
Thus the cochain map induced by the invertible
endomorphism \(C_M:M\to M\) is nullhomotopic.
It is an isomorphism on cohomology and also zero
on cohomology; hence every cohomology group is zero.
This proof is the central-action argument of
\cite[Theorem 7.8.9]{weibel}, with its comparison
homotopy made explicit.
\end{proof}

\begin{theorem}[Whitehead's first lemma]\label{res:C04}
For every finite-dimensional \(\mathfrak g\)-module \(M\),
\[
 H^1(\mathfrak g,M)=0.
 \tag{3.6}\label{eq:whitehead-first}
\]
Equivalently, every linear map \(a:\mathfrak g\to M\)
satisfying
\[
 a([x,y])=x\,a(y)-y\,a(x)
\]
has the form \(a(x)=xm\) for some \(m\in M\).
\end{theorem}

\begin{proof}
The low-degree Chevalley--Eilenberg formulas identify
one-cocycles and one-coboundaries exactly as in the
statement; see \cite[Chapter 6, A08]{homological}.
If \(M\) is simple and nontrivial, apply
Theorem~\ref{res:C03}. A simple trivial module is \(k\),
and
\[
 H^1(\mathfrak g,k)=
 \operatorname{Hom}_k(\mathfrak g/[\mathfrak g,\mathfrak g],k)=0
\]
by perfectness from Theorem~\ref{res:B06}.
For arbitrary \(M\), induct along a composition
series. For \(0\to U\to M\to Q\to0\), the cochain
complexes form a degreewise short exact sequence,
since \(\operatorname{Hom}_k(\bigwedge^q\mathfrak g,-)\)
is exact over a field. The cohomology exact sequence
contains
\(H^1(\mathfrak g,U)\to H^1(\mathfrak g,M)
\to H^1(\mathfrak g,Q)\).
Both outer terms vanish by induction, proving the
assertion.
\end{proof}

\begin{theorem}[Weyl complete reducibility]\label{res:C05}
Every submodule of a finite-dimensional
\(\mathfrak g\)-module has a \(\mathfrak g\)-invariant
complement. Consequently every such module is a
finite direct sum of simple modules.
\end{theorem}

\begin{proof}
Let \(0\to W\to V\xrightarrow{p}Q\to0\), and choose
a \(k\)-linear section \(s:Q\to V\).
Its equivariance defect is
\[
 a(x)=\rho_V(x)s-s\rho_Q(x)\in\operatorname{Hom}_k(Q,W),
 \tag{3.7}\label{eq:splitting-cocycle}
\]
since applying \(p\) gives zero. The Hom module has
the action \(x\cdot t=\rho_W(x)t-t\rho_Q(x)\).
Expanding \eqref{eq:splitting-cocycle} at \([x,y]\)
and using the representation identities gives
\(a([x,y])=x\cdot a(y)-y\cdot a(x)\);
the two mixed products cancel.
Whitehead's first lemma supplies \(t:Q\to W\) with
\(a(x)=x\cdot t\).
The adjusted section \(s-t\) has zero defect and
still satisfies \(p(s-t)=1\). Its image is the
desired invariant complement. Split off a simple
submodule, and repeat on the smaller complement;
finite dimension terminates the process.
\end{proof}

This proof distinguishes two statements often both
called semisimplicity: Theorem~\ref{res:B06} concerns
simple ideals of the algebra, whereas
Theorem~\ref{res:C05} concerns arbitrary modules over
it. Nor does the theorem apply to an arbitrary
reductive algebra with no condition on the action of
its center. A central generator can act by a
nonzero nilpotent Jordan block.

\section{Rank-one consequences and functorial Jordan decomposition}

\begin{theorem}[Clebsch--Gordan, duals, and invariant forms]
\label{res:C06}
Over an algebraically closed characteristic-zero field,
every finite-dimensional \(\mathfrak{sl}_2\)-module is
a direct sum of the \(L_n\), and
\[
 L_m\otimes L_n\simeq
 \bigoplus_{j=0}^{\min(m,n)}L_{m+n-2j}.
 \tag{3.8}\label{eq:clebsch-gordan}
\]
Each \(L_n\) is self-dual and has, up to a scalar,
a unique nonzero invariant bilinear form. It is
nondegenerate, symmetric for even \(n\), and
alternating for odd \(n\).
\end{theorem}

\begin{proof}
Simplicity of \(\mathfrak{sl}_2\), the Killing
criterion, and Weyl's theorem give complete
reducibility; Theorem~\ref{res:C01} identifies its
simple summands.
Record the \(h\)-weights by
\(\chi_n(t)=t^n+t^{n-2}+\cdots+t^{-n}\).
The coefficient of \(t^{m+n-2s}\) in
\(\chi_m\chi_n\) counts pairs
\((a,b)\) with \(0\leq a\leq m,0\leq b\leq n,a+b=s\).
If \(m\leq n\), this number is
\[
 \min(m,s)-\max(0,s-n)+1
\]
when the interval is nonempty, and zero otherwise.
The same count results from
\(\sum_{j=0}^{m}\chi_{m+n-2j}\): its \(j\)-th
term contributes exactly when
\(j\leq s\leq m+n-j\). Hence
\(\chi_m\chi_n=\sum_{j=0}^{\min(m,n)}
\chi_{m+n-2j}\).
The characters \(\chi_r\) are linearly independent,
because a nontrivial finite linear relation has a
largest exponent with a nonzero coefficient.
Complete reducibility therefore gives
\eqref{eq:clebsch-gordan}.

In the basis \eqref{eq:sl2-basis}, define
\[
 B(v_i,v_j)=
 \begin{cases}(-1)^i,&i+j=n,\\0,&i+j\ne n.\end{cases}
 \tag{3.9}\label{eq:sl2-form}
\]
For \(h\), invariance
\(B(hv_i,v_j)+B(v_i,hv_j)=0\) follows from the
weights. For \(f\), its only possible nonzero
case is \(i+j=n-1\), and the two terms are
\((-1)^{i+1}+(-1)^i=0\).
For \(e\), its only possible case is \(i+j=n+1\);
the coefficients \(i(n-i+1)\) and
\(j(n-j+1)\) are then equal, and the signs
are opposite. Thus \(B\) is invariant.
Its antidiagonal entries are nonzero, so it
identifies \(L_n\) with \(L_n^*\).
Any invariant form must be antidiagonal by \(h\),
and the \(f\)-equations force successive entries
to change sign. Thus it is a scalar multiple
of \eqref{eq:sl2-form}.
Finally
\(B(v_j,v_i)=(-1)^n B(v_i,v_j)\).
For odd \(n\) its diagonal is zero, giving the
alternating assertion.
\end{proof}

\begin{theorem}[Representations preserve Jordan decomposition]
\label{res:C07}
Let \(k\) have characteristic zero and let
\(\rho:\mathfrak g\to\operatorname{End}_k(V)\)
be finite dimensional, with \(\mathfrak g\) semisimple.
For the abstract Jordan decomposition of
Theorem~\ref{res:B07},
\[
 \rho(x_s)=\rho(x)_s,\qquad
 \rho(x_n)=\rho(x)_n.
 \tag{3.10}\label{eq:representation-jordan}
\]
In particular, each toral subalgebra of \(\mathfrak g\) acts semisimply on \(V\).
Over an algebraically closed field, the restriction of \(\rho\) to such a
subalgebra gives a decomposition of \(V\) into simultaneous weight spaces.
\end{theorem}

\begin{proof}
First take an algebraically closed field and a
simple nonzero \(V\), and put
\(\mathfrak a=\rho(\mathfrak g)\).
It is semisimple. Its normalizer in
\(\operatorname{End}(V)\) acts by derivations on
\(\mathfrak a\). Since all such derivations are
inner, every normalizing operator is \(z+c\),
with \(z\in\mathfrak a\) and \(c\) centralizing
\(\mathfrak a\). Schur's lemma makes \(c\)
a scalar.

Let \(y\in\mathfrak a\). Its matrix semisimple
part \(y_s\) normalizes \(\mathfrak a\):
\(\operatorname{ad}y\) preserves \(\mathfrak a\),
so its polynomial semisimple part does also,
and that part is \(\operatorname{ad}y_s\), by
the commuting semisimple/nilpotent calculation
used in the trace criterion. Write \(y_s=z+c\)
as above. Since \(\mathfrak a\) is perfect,
\(\operatorname{tr}y=\operatorname{tr}z=0\).
Nilpotents have zero trace, so
\(\operatorname{tr}y_s=\operatorname{tr}y=0\).
It follows that \((\dim V)c=0\), hence \(c=0\).
Thus \(y_s\) and \(y_n=y-y_s\) actually belong
to \(\mathfrak a\). Their adjoint operators
on \(\mathfrak a\) are semisimple and nilpotent,
respectively: on the ambient endomorphism
space they have these properties, and
restriction to an invariant subspace preserves
them. They are therefore the abstract Jordan
parts in \(\mathfrak a\).

The kernel of \(\rho\) is a direct ideal
summand of \(\mathfrak g\), by
Theorem~\ref{res:B06}. The adjoint action is
block diagonal for this splitting, so abstract
Jordan decomposition commutes with projection
onto the complementary summand. The preceding
paragraph proves \eqref{eq:representation-jordan}
on \(V\).
For an arbitrary \(V\), Weyl's theorem writes
it as a sum of simple modules. Jordan parts
of block diagonal matrices are the sums of
the parts on their blocks, proving the same
formula.

For a nonclosed field, extend to an algebraic
closure. The two sides of
\eqref{eq:representation-jordan} are scalar
extensions of their \(k\)-linear operators,
by the polynomial and uniqueness construction
of Jordan decomposition. Their equality
therefore descends.

Finally a toral algebra is abelian with
semisimple adjoint action for each element.
The formula just proved makes its representing
operators semisimple. Commuting semisimple
operators over an algebraically closed field
are simultaneously diagonalizable: diagonalize
one generator, observe that its eigenspaces
are stable under all the others, and repeat
inside those eigenspaces. The simultaneous
eigenvalues are linear functionals on the
toral algebra and give its weight spaces.
\end{proof}

\chapter{Cartan Subalgebras and Roots}
\label{ch:roots}

Throughout this chapter \(\mathfrak g\) is a finite-dimensional
semisimple complex Lie algebra and \(B\) is its Killing form.
The words semisimple and nilpotent, when applied to an element, refer
to its adjoint endomorphism. Inner automorphisms are generated by
\(\exp(\operatorname{ad}x)\), \(x\in\mathfrak g\); the exponential here
is an ordinary convergent matrix exponential. No integration theorem
for abstract Lie algebras is needed.

\section{Tori and their centralizers}

\begin{proposition}[Maximal toral subalgebras]\label{res:D01}
A subalgebra consisting of commuting semisimple elements is called
toral. A maximal toral subalgebra \(\mathfrak h\) exists and satisfies
\[
 C_{\mathfrak g}(\mathfrak h)=\mathfrak h.
\]
In particular \(\mathfrak h\) consists of simultaneously diagonalizable
adjoint endomorphisms.
\end{proposition}
\begin{proof}
Commuting diagonalizable matrices are simultaneously diagonalizable:
an eigenspace of one is preserved by all the others, and induction on
dimension finishes the assertion. Consequently the sum of two commuting
toral subalgebras is toral. Maximal ones exist by finite dimensionality.
If \(\mathfrak g\ne0\), one is nonzero: otherwise every element would
be adjoint-nilpotent by the Jordan decomposition
of Theorem~\ref{res:B07}, and Engel would make \(\mathfrak g\) nilpotent.

Write \(\mathfrak c=C_{\mathfrak g}(\mathfrak h)\). Simultaneous
diagonalization gives
\(\mathfrak g=\bigoplus_\lambda\mathfrak g_\lambda\), where
\([h,x]=\lambda(h)x\) on \(\mathfrak g_\lambda\).
Invariance says
\[
 (\lambda(h)+\mu(h))B(x,y)=0
 \quad(x\in\mathfrak g_\lambda,\ y\in\mathfrak g_\mu).
\]
Thus \(B|_{\mathfrak c}\) is nondegenerate: a vector in its radical
is also orthogonal to every nonzero weight space and hence to all
of \(\mathfrak g\).

For \(x\in\mathfrak c\), both Jordan components \(x_s,x_n\) centralize
\(\mathfrak h\). This follows either from uniqueness of Jordan
decomposition or from the polynomial expressions for the two adjoint
components. Maximality implies \(x_s\in\mathfrak h\).
Therefore \(\operatorname{ad}_{\mathfrak c}x=
\operatorname{ad}_{\mathfrak c}x_n\) is nilpotent.
Engel implies that \(\mathfrak c\) is nilpotent, in particular solvable.
Lie's theorem triangularizes its action on \(\mathfrak g\).
The operators from \([\mathfrak c,\mathfrak c]\) then have zero
diagonal, so
\(B([\mathfrak c,\mathfrak c],\mathfrak c)=0\).
Nondegeneracy gives \([\mathfrak c,\mathfrak c]=0\).
Finally \(\operatorname{ad}x_n\) commutes with every
\(\operatorname{ad}y\), \(y\in\mathfrak c\). Their product is nilpotent,
so \(B(x_n,y)=0\) for every such \(y\). Hence \(x_n=0\).
Every element of \(\mathfrak c\) is semisimple and
\(\mathfrak c=\mathfrak h\).
\end{proof}

\begin{proposition}[Root decomposition and opposite pairings]\label{res:D02}
For a maximal toral \(\mathfrak h\) there is a decomposition
\[
 \mathfrak g=\mathfrak h\oplus
       \bigoplus_{\alpha\in R}\mathfrak g_\alpha,\qquad
 \mathfrak g_\alpha=\{x:[h,x]=\alpha(h)x\text{ for all }h\in\mathfrak h\},
\]
where \(R\subset\mathfrak h^*\setminus\{0\}\) is finite.
One has
\([\,\mathfrak g_\alpha,\mathfrak g_\beta\,]\subset
\mathfrak g_{\alpha+\beta}\), with the zero weight interpreted as
\(\mathfrak h\). Moreover \(B(\mathfrak g_\alpha,\mathfrak g_\beta)=0\)
unless \(\alpha+\beta=0\), and
\(\mathfrak g_\alpha\times\mathfrak g_{-\alpha}\to\mathbb C\)
is a perfect pairing.
\end{proposition}
\begin{proof}
The decomposition and its zero space follow from
Proposition~\ref{res:D01}. Jacobi gives
\([h,[x,y]]=(\alpha+\beta)(h)[x,y]\).
The orthogonality equation in that proof gives the asserted
orthogonality. If a vector in \(\mathfrak g_\alpha\) pairs to zero
with \(\mathfrak g_{-\alpha}\), it pairs to zero with every summand,
so is zero. Apply this on both sides of the finite-dimensional pairing.
The same argument establishes nondegeneracy on \(\mathfrak h\).
\end{proof}

\section{Root triples and root strings}

Use \(B|_{\mathfrak h}\) to define \(t_\alpha\in\mathfrak h\) by
\(B(t_\alpha,h)=\alpha(h)\), and put
\((\alpha,\beta)=B(t_\alpha,t_\beta)\). Invariance immediately gives
\[
 [e,f]=B(e,f)t_\alpha
 \quad(e\in\mathfrak g_\alpha,\ f\in\mathfrak g_{-\alpha}).
 \tag{4.1}\label{eq:opposite-bracket}
\]

\begin{theorem}[Root triples and reducedness]\label{res:D03}
For every root \(\alpha\), \((\alpha,\alpha)\ne0\).
There are \(e_\alpha\in\mathfrak g_\alpha\),
\(f_\alpha\in\mathfrak g_{-\alpha}\), and
\[
 h_\alpha=\frac{2t_\alpha}{(\alpha,\alpha)}
\]
with \([e_\alpha,f_\alpha]=h_\alpha\),
\([h_\alpha,e_\alpha]=2e_\alpha\), and
\([h_\alpha,f_\alpha]=-2f_\alpha\).
Every root space is one dimensional, and the only roots proportional
to \(\alpha\) are \(\alpha\) and \(-\alpha\).
\end{theorem}
\begin{proof}
Choose \(e,f\) with \(B(e,f)\ne0\), and let \(h=[e,f]\).
If \(\alpha(t_\alpha)=0\), the three-dimensional span of \(e,f,h\)
is solvable: \(h\) commutes with \(e,f\). Lie's theorem makes
\(\operatorname{ad}h=[\operatorname{ad}e,\operatorname{ad}f]\)
strictly upper triangular, hence nilpotent. Since \(h\in\mathfrak h\),
it is also semisimple. The zero center of \(\mathfrak g\) would
give \(h=0\), contradicting (4.1). Normalize \(B(e,f)\) to
\(2/(\alpha,\alpha)\); (4.1) gives the three relations.

First consider the root spaces whose roots are integral multiples
of \(\alpha\), together with \(\mathbb C h_\alpha\).
Their sum is a subalgebra stable under this \(\mathfrak{sl}_2\).
All its \(h_\alpha\)-weights are even, and its zero space is
one dimensional. Complete reducibility for \(\mathfrak{sl}_2\)
implies that it is one simple even-highest-weight module: every
such simple summand has a one-dimensional zero space.
It contains the adjoint three-dimensional module
\(\langle e_\alpha,h_\alpha,f_\alpha\rangle\), so equals that module.
Thus \(\dim\mathfrak g_\alpha=1\) and no integral multiple
\(m\alpha\), \(|m|\ge2\), is a root.

For an arbitrary proportional root \(c\alpha\), integrality of
\(\mathfrak{sl}_2\)-weights gives \(2c\in\mathbb Z\).
Interchanging the two roots gives \(2/c\in\mathbb Z\).
Consequently \(c\in\{\pm1,\pm2,\pm\tfrac12\}\).
The last four cases contradict the already established exclusion
of integral multiples, applied to one of the two roots.
\end{proof}

\begin{theorem}[Strings, integers, and reflections]\label{res:D04}
If \(\beta\ne\pm\alpha\), the roots in \(\beta+\mathbb Z\alpha\)
are exactly
\[
 \beta-p\alpha,\ldots,\beta+q\alpha,\qquad
 p,q\ge0,\qquad p-q=\beta(h_\alpha).
\]
All successive nontrivial brackets along this string are nonzero.
In particular
\[
 \beta(h_\alpha)=\frac{2(\beta,\alpha)}{(\alpha,\alpha)}\in\mathbb Z,
 \qquad
 s_\alpha(\beta)=\beta-\beta(h_\alpha)\alpha\in R.
\]
The reflection is implemented on \(\mathfrak g\) by
\[
 n_\alpha=
 \exp(\operatorname{ad}e_\alpha)
 \exp(-\operatorname{ad}f_\alpha)
 \exp(\operatorname{ad}e_\alpha).
 \tag{4.2}\label{eq:root-reflection-lift}
\]
\end{theorem}
\begin{proof}
The finite sum
\(V=\bigoplus_{k\in\mathbb Z}\mathfrak g_{\beta+k\alpha}\)
is a module for the root triple. Each weight space for
\(h_\alpha\) is one dimensional, and all its weights differ by
even integers. Two nonzero simple summands would share the
weight \(0\) if their highest weights are even, or the weights
\(\pm1\) if odd. Hence \(V\) is simple.
Its \(\mathfrak{sl}_2\) weights form one interval, with endpoints
\(\beta(h_\alpha)-2p\) and \(\beta(h_\alpha)+2q\), opposite to each
other. This proves the equation and the nonzero bracket assertion.
Symmetry of that interval sends the root \(\beta\) to
\(\beta-\beta(h_\alpha)\alpha\).

On each finite simple \(\mathfrak{sl}_2\) module the three
exponentials in (4.2) interchange opposite weights. On the triple
itself direct multiplication gives \(h_\alpha\mapsto-h_\alpha\).
For \(h\in\ker\alpha\), all three factors fix \(h\), so
\(n_\alpha(h)=h-\alpha(h)h_\alpha\).
It follows by applying brackets that \(n_\alpha\) maps the entire
\(\beta\)-space to the \(s_\alpha\beta\)-space.
The exponentials are finite on these modules, and an exponential
of a derivation preserves brackets; thus they are automorphisms.
\end{proof}

\begin{proposition}[The real root space]\label{res:D05}
The roots span \(\mathfrak h^*\). Let
\(\mathfrak h_{\mathbb R}=\operatorname{span}_{\mathbb R}
\{h_\alpha:\alpha\in R\}\). Then
\[
 \mathfrak h=\mathfrak h_{\mathbb R}\oplus
 i\mathfrak h_{\mathbb R},
\]
and \(B\) is real positive definite on \(\mathfrak h_{\mathbb R}\).
The roots, viewed in its real dual \(E\), form a finite reduced
crystallographic root system. Their inner products for the
Killing-induced form are rational, and their rational span is
a rational form of \(E\).
\end{proposition}
\begin{proof}
If all roots vanish on \(h\), then \(\operatorname{ad}h=0\);
therefore \(h=0\). Nondegeneracy identifies the spanning property
with the fact that the \(h_\alpha\) span \(\mathfrak h\) over
\(\mathbb C\).
For \(h\) in their real span all eigenvalues of \(\operatorname{ad}h\)
are real, by Theorem~\ref{res:D04}. If also \(h\in
i\mathfrak h_{\mathbb R}\), they are purely imaginary and hence
zero, giving \(h=0\). Moreover
\[
 B(h,h)=\sum_{\beta\in R}\beta(h)^2>0\quad(h\ne0),
\]
and polarization makes \(B\) real valued on this real space.
Dimension now gives the displayed complexification.
The dual Euclidean form is the restriction of the form already
used for roots. Reducedness, the integral coroot pairings,
reflection invariance, and spanning have all been proved.
These are precisely the root system axioms.

The rationality assertion follows without choosing a Dynkin type.
By the one-dimensional root spaces and integral coroot pairings,
\[
 B(h_\alpha,h_\beta)
     =\sum_{\gamma\in R}\gamma(h_\alpha)\gamma(h_\beta)
       \in\mathbb Z.
\]
Since \(B(h_\alpha,h_\alpha)=4/(\alpha,\alpha)\ne0\),
the square length \((\alpha,\alpha)\) is rational. The identity
\((\alpha,\beta)=\beta(h_\alpha)(\alpha,\alpha)/2\)
then gives all rational inner products. Choose any real basis
from the roots. Its nonsingular Gram matrix has rational entries,
as do the pairings of any other root with that basis; solving
the Gram system expresses that root with rational coefficients.
Thus the rational span has dimension \(\dim_{\mathbb R}E\),
and its scalar extension to \(\mathbb R\) is \(E\).
\end{proof}

\section{Regular elements and conjugacy}

The following argument separates the ordinary kernel dimension
from the multiplicity of zero in a characteristic polynomial.
For \(x\in\mathfrak g\), let \(n(x)\) be the latter multiplicity,
and set \(r=\min_x n(x)\).

\begin{theorem}[Regular semisimple density]\label{res:D06}
The set of elements with \(n(x)=r\) is a connected dense open
subset of \(\mathfrak g\); every one of these elements is semisimple.
For every maximal toral \(\mathfrak h\),
\(\dim\mathfrak h=r\), and
\[
 \mathfrak h_{\rm reg}
 =\{h\in\mathfrak h:\alpha(h)\ne0\text{ for all }\alpha\in R\}.
\]
A regular semisimple element has a unique maximal toral
subalgebra containing it, namely its centralizer.
For every \(x\in\mathfrak g\),
\(\dim C_{\mathfrak g}(x)\ge r\), with equality on this open set.
\end{theorem}
\begin{proof}
For \(\mathfrak g=0\) every assertion is immediate, with
\(r=0\); hence assume \(\mathfrak g\ne0\).
Write
\(\det(t-\operatorname{ad}x)=t^r
\bigl(t^{m}+a_{m-1}(x)t^{m-1}+\cdots+a_0(x)\bigr)\),
where \(m=\dim\mathfrak g-r\).
By the definition of \(r\), \(a_0\) is a nonzero polynomial,
and its nonvanishing is exactly \(n(x)=r\).
The complement of a nonzero complex polynomial's zero set is
open and dense: vanishing on a ball forces all Taylor
coefficients to vanish. It is path connected because the complex
line between two of its points meets the zero set at only
finitely many parameters, and a path in that line can avoid them.

Fix a maximal toral \(\mathfrak h\), and put
\(V=\mathfrak h_{\rm reg}\) as defined in the statement.
For \(h\in V\),
\(\mathfrak g=[\mathfrak g,h]\oplus\mathfrak h\).
Choose a linear complement \(M\) to \(\mathfrak h\).
The analytic map
\[
 M\times\mathfrak h\longrightarrow\mathfrak g,\qquad
 (y,z)\longmapsto\exp(\operatorname{ad}y)z
\]
has differential \((y,z)\mapsto[y,h]+z\) at \((0,h)\),
an isomorphism. The inverse function theorem makes the inner
saturation of \(V\) open. It therefore meets the dense set
\(n(x)=r\). Nullity is invariant under automorphisms, and equals
\(\dim\mathfrak h\) on \(V\). Thus \(\dim\mathfrak h=r\).

Let first \(x\) be semisimple with \(n(x)=r\).
Its centralizer \(\mathfrak c\) has dimension \(r\), and
\(B|_{\mathfrak c}\) is nondegenerate by the zero weight argument.
For \(y\in\mathfrak c\), \(\operatorname{ad}(x+ty)\)
is invertible on the sum of nonzero \(x\)-eigenspaces for small
\(t\). On \(\mathfrak c\) it equals \(t\operatorname{ad}y\).
Since its zero multiplicity cannot be less than \(r=\dim\mathfrak c\),
\(\operatorname{ad}_{\mathfrak c}y\) is nilpotent.
Engel and Lie's theorem, followed by nondegeneracy of
\(B|_{\mathfrak c}\), show that \(\mathfrak c\) is abelian,
exactly as in Proposition~\ref{res:D01}.
For the same reason every nilpotent component of an element of
\(\mathfrak c\) is zero. Thus \(\mathfrak c\) is toral.

For arbitrary \(x\) of minimal nullity, the Jordan components
commute, and \(\operatorname{ad}x\) and \(\operatorname{ad}x_s\)
have the same generalized zero space. Thus \(n(x_s)=r\).
The preceding paragraph makes \(C_{\mathfrak g}(x_s)\) toral.
It contains \(x\), so \(x\) is semisimple and \(x_n=0\).
Any toral subalgebra containing \(x\) is contained in this
centralizer; maximality then gives uniqueness.

On the dense open set just described,
\(\operatorname{ad}x\) has rank \(\dim\mathfrak g-r\).
Every larger matrix minor therefore vanishes on a dense set
and hence identically. This proves the claimed lower bound
for the kernel dimension at all \(x\), not just semisimple ones.
\end{proof}

\begin{theorem}[Cartan subalgebras and their conjugacy]\label{res:D07}
A Cartan subalgebra, defined as a nilpotent self-normalizing
subalgebra, is the same as a maximal toral subalgebra of
\(\mathfrak g\). All of them are conjugate by inner automorphisms.
\end{theorem}
\begin{proof}
A maximal toral subalgebra is abelian and self-normalizing:
in the root decomposition, a nonzero root component cannot
normalize its zero space. It is therefore Cartan.

Conversely let \(\mathfrak a\) be nilpotent and
\(N_{\mathfrak g}(\mathfrak a)=\mathfrak a\).
We shall use the following elementary consequence of Engel:
if \(\mathfrak a\) acts nilpotently on a space \(V\supset\mathfrak a\)
and preserves \(\mathfrak a\), then \(V/\mathfrak a\ne0\)
has a nonzero common zero vector. For the adjoint action such
a vector lifts to a normalizing vector. Thus self-normalization
forces \(V=\mathfrak a\).

For \(x\in\mathfrak a\), both \(\operatorname{ad}x_s\) and
\(\operatorname{ad}x_n\), being polynomials in
\(\operatorname{ad}x\), preserve \(\mathfrak a\).
Hence \(x_s,x_n\in\mathfrak a\).
On \(\mathfrak a\) the action of \(x\) is nilpotent,
so the semisimple restriction of \(\operatorname{ad}x_s\)
is zero. The semisimple elements of \(\mathfrak a\) therefore
form a central torus \(\mathfrak t\).
On \(V=C_{\mathfrak g}(\mathfrak t)\) every \(x\in\mathfrak a\)
acts as the nilpotent operator \(\operatorname{ad}x_n\).
The preceding consequence of Engel gives \(V=\mathfrak a\).
In particular \(B|_{\mathfrak a}\) is nondegenerate.
Lie's theorem triangularizes \(\mathfrak a\) on \(\mathfrak g\);
the nilpotent components \(x_n\) have zero diagonal, hence
\(B(x_n,\mathfrak a)=0\). Thus all \(x_n=0\),
and \(\mathfrak a=\mathfrak t\) is maximal toral.

For conjugacy, every regular semisimple \(x\) has the unique
Cartan \(C(x)\). Declare two such elements equivalent if their
Cartans are inner-conjugate. For any fixed Cartan
\(\mathfrak h\), the class in question is precisely its inner
saturation of \(\mathfrak h_{\rm reg}\).
The differential calculation in Theorem~\ref{res:D06}
makes each such class open. The regular semisimple set is
connected, so its partition into open classes has a single class.
Every Cartan contains a regular element; therefore all Cartans
are conjugate.
\end{proof}

\section{Matrix models}

\begin{example}[The classical root systems]\label{res:D08}
In the following table the \(\varepsilon_i\) are coordinate
functionals on diagonal Cartans:
\[
\begin{array}{c|c|c}
\mathfrak g&\mathfrak h&R\\ \hline
\mathfrak{sl}_{n+1}&
 \operatorname{diag}(t_1,\ldots,t_{n+1}),\ \sum t_i=0&
 \varepsilon_i-\varepsilon_j\ (i\ne j)\\
\mathfrak{so}_{2n+1}&
 \operatorname{diag}(t_1,\ldots,t_n,0,-t_1,\ldots,-t_n)&
 \pm\varepsilon_i,\ \pm\varepsilon_i\pm\varepsilon_j\\
\mathfrak{sp}_{2n}&
 \operatorname{diag}(t_1,\ldots,t_n,-t_1,\ldots,-t_n)&
 \pm2\varepsilon_i,\ \pm\varepsilon_i\pm\varepsilon_j\\
\mathfrak{so}_{2n}&
 \operatorname{diag}(t_1,\ldots,t_n,-t_1,\ldots,-t_n)&
 \pm\varepsilon_i\pm\varepsilon_j .
\end{array}
\]
The orthogonal forms are chosen in hyperbolic coordinates.
Coroots are \(2\alpha/(\alpha,\alpha)\) under the Euclidean
identification. The low rank coincidences are
\(B_1=C_1=A_1\), \(B_2=C_2\), \(D_2=A_1\sqcup A_1\),
and \(D_3=A_3\); \(D_1\) is not a semisimple case.
\end{example}
\begin{proof}
For a diagonal matrix \(t\),
\([t,E_{ij}]=(t_i-t_j)E_{ij}\).
In the symplectic case the equation \(X^TJ+JX=0\) gives
\[
 X=\begin{pmatrix}A&B\\C&-A^T\end{pmatrix},
 \qquad B=B^T,\quad C=C^T.
\]
The off-diagonal entries of \(A\) have weights
\(\varepsilon_i-\varepsilon_j\); the off-diagonal symmetric
entries of \(B,C\) have weights
\(\pm(\varepsilon_i+\varepsilon_j)\), and their diagonal
entries have weights \(\pm2\varepsilon_i\).
For the even orthogonal form
\(\left(\begin{smallmatrix}0&I\\I&0\end{smallmatrix}\right)\),
the same block equation makes \(B,C\) skew symmetric and
therefore removes the diagonal weights.
In odd dimension the extra row and column give
\(\pm\varepsilon_i\). These calculations account for every
matrix entry allowed by the defining equations, so the zero
spaces are exactly the displayed Cartans.
Pair opposite matrix entries to obtain the corresponding
root triples, which gives the stated coroot normalization.
Finally the listed small root sets are carried into each
other by permuting or rescaling the Euclidean coordinates;
their Lie algebra isomorphisms also follow from the
reconstruction theorem proved in Chapter~\ref{ch:classification}.
\end{proof}

The algebraic root calculations follow
\cite[Sections 8 and 15--16]{humphreys} and
\cite[Sections 19--20]{etingof}. The use of the characteristic
polynomial and its connected nonvanishing locus in the conjugacy
argument explains why the assertion about a dense regular set
concerns regular \emph{semisimple} elements, not every element
whose ordinary centralizer has minimal dimension.

\chapter{Root Systems and Weyl Groups}
\label{ch:root-systems}

\section{Positive roots and chambers}

\begin{definition}[Finite root systems]\label{res:E01}
A reduced crystallographic root system in a Euclidean space \(E\)
is a finite spanning subset \(R\subset E\setminus\{0\}\) such that
\[
 s_\alpha(R)=R,\qquad
 \langle\beta,\alpha^\vee\rangle
 =\frac{2(\beta,\alpha)}{(\alpha,\alpha)}\in\mathbb Z,\qquad
 R\cap\mathbb R\alpha=\{\alpha,-\alpha\}.
\]
The coroots \(2\alpha/(\alpha,\alpha)\) form another root system,
the dual one. Orthogonal irreducible components are unique.
The rank two possibilities are \(A_1\sqcup A_1,A_2,B_2,G_2\),
with respectively \(4,6,8,12\) roots.
\end{definition}
\begin{proof}
The coroot reflection is the same orthogonal reflection, and
\[
 \frac{2(\alpha^\vee,\beta^\vee)}
      {(\alpha^\vee,\alpha^\vee)}
 =\frac{2(\alpha,\beta)}{(\beta,\beta)}.
\]
This proves the dual assertions and double duality.
For two independent roots making angle \(\theta\),
the product of the two integral Cartan numbers is
\(4\cos^2\theta\), a nonnegative integer less than \(4\).
The nonorthogonal cases give the squared length ratios
\(1,2,3\), and the acute angles \(60,45,30\) degrees.
To see that the rank two list is exhaustive, take two adjacent
root rays making the smallest positive angle. Reflection in
their perpendicular lines repeats this sector around the
circle. The angle restrictions give \(2,3,4,6\) pairs of
opposite rays. No extra root ray can occur inside a sector,
by its minimality; a second length on a ray is excluded by
reducedness. Integrality fixes equal lengths in \(A_2\)
and the alternating ratios \(2,3\) in \(B_2,G_2\).
The coordinate models below verify their existence.
Uniqueness of the orthogonal decomposition will also follow
from the simple-root basis proved next: it is the decomposition
of its graph into connected components.
\end{proof}

\begin{theorem}[Simple roots]\label{res:E02}
Choose \(v\in E\) off the root hyperplanes and put
\(R^+=\{\alpha:(v,\alpha)>0\}\). The positive roots that are
not sums of two positive roots form a basis
\(\Pi=\{\alpha_1,\ldots,\alpha_r\}\) of \(E\).
Every root is an integral combination of \(\Pi\), with all
coefficients nonnegative or all nonpositive.
The simple reflections \(s_i=s_{\alpha_i}\) generate the Weyl
group \(W=\langle s_\alpha:\alpha\in R\rangle\), and
\(R=W\Pi\).
\end{theorem}
\begin{proof}
If a positive root is decomposable, split it into two positive
roots and repeat. The values of \((v,-)\) on \(R^+\) form a
finite set of positive numbers, so the process terminates.
Thus every positive root is a sum of simple ones.
Distinct simple roots have nonpositive inner product: if
\((\alpha,\beta)>0\), inspection of the four rank two systems
gives \(\alpha-\beta\) or \(\beta-\alpha\) as a root, making
one of the two simple roots decomposable.
For linear independence, split a nonzero relation into
\(u=\sum_{i\in I}a_i\alpha_i=\sum_{j\in J}b_j\alpha_j\),
where \(I,J\) are disjoint and all coefficients positive.
Then \(u\ne0\), as pairing with \(v\) shows, whereas
\((u,u)=\sum a_ib_j(\alpha_i,\alpha_j)\le0\), a contradiction.
This proves the basis and the sign assertion.

If \(\beta=\sum n_i\alpha_i\) is positive and not simple,
some \((\beta,\alpha_i)>0\), since
\((\beta,\beta)=\sum n_i(\beta,\alpha_i)>0\).
The root \(s_i\beta=\beta-
\langle\beta,\alpha_i^\vee\rangle\alpha_i\)
has smaller height. It cannot be negative unless \(\beta\)
is a multiple of \(\alpha_i\), because all other coefficients
are unchanged. In the latter case reducedness makes it simple.
Thus repeated simple reflections take any root to a simple
root, proving \(R=W_0\Pi\) for \(W_0=\langle s_i\rangle\).
For \(\beta=w\alpha_i\), \(s_\beta=ws_iw^{-1}\in W_0\),
so \(W_0=W\).
One also obtains the useful lowering assertion:
if \(\beta>0\) is not simple, then \(\beta-\alpha_i\)
is a positive root for some \(i\). Indeed the root string
through \(\beta\) in the direction just chosen contains that
root. These string assertions in an abstract root system
can equivalently be checked in its rank two subsystem.
\end{proof}

\begin{theorem}[Chambers, length, and a fundamental domain]\label{res:E03}
The Weyl group acts simply transitively on the chambers in
\(E\setminus\bigcup_{\alpha\in R}\alpha^\perp\).
The chamber for \(\Pi\) is
\(C=\{x:(x,\alpha_i)>0\ \forall i\}\).
For \(w\in W\) its shortest simple-reflection word has length
\[
 \ell(w)=|\{\alpha>0:w\alpha<0\}|.
\]
Every \(W\)-orbit has exactly one point in \(\overline C\).
\end{theorem}
\begin{proof}
The chamber corresponding to a sign choice is an intersection
of open half spaces, hence convex. The positive-root
expansions show that the walls of \(C\) are precisely the
simple-root hyperplanes. A generic segment between two
chambers crosses one wall at a time. Reflection in that wall
takes one adjacent chamber to the other, so \(W\) is transitive.

We give the word argument establishing freeness, rather than
assuming it in the gallery argument. A simple reflection
permutes \(R^+\setminus\{\alpha_i\}\): coefficients outside
\(\alpha_i\) remain nonnegative, and only \(\alpha_i\)
can change sign. Suppose \(u=s_{i_1}\cdots s_{i_m}\)
and \(u\alpha<0\) for a positive root \(\alpha\).
Follow \(\alpha\) through this word from right to left.
At the first sign change the positive root must equal the
simple root being reflected. If the suffix before that step
is \(v\), then \(v\alpha=\alpha_{i_k}\) and
\(s_\alpha=v^{-1}s_{i_k}v\). Thus \(us_\alpha\)
has an expression obtained by deleting the \(k\)-th letter.
In a shortest expression \(w=us_i\), this implies
\(u\alpha_i>0\), for otherwise two letters could be deleted.
Counting the positive roots sent negative gives
\(N(us_i)=N(u)+1\) in that case. Induction shows that every
shortest word has length \(N(w)\).
If \(wC=C\), it sends no positive root negative, hence its
shortest word has length zero and \(w=1\).

Existence of a representative in \(\overline C\) follows
by continuity from transitivity on chambers, or by reflecting
successively across a wall on whose negative side the point lies.
For uniqueness we record a consequence of a reduced word
\(w=s_{i_1}\cdots s_{i_m}\). Each
\(\gamma_k=s_{i_1}\cdots s_{i_{k-1}}\alpha_{i_k}\)
is positive, by the preceding deletion argument. For dominant
\(\lambda\),
\[
 \lambda-w\lambda
 =\sum_{k=1}^m
   \langle\lambda,\alpha_{i_k}^\vee\rangle\gamma_k
 \quad\text{belongs to }\sum_i\mathbb R_{\ge0}\alpha_i.
\]
If \(\mu=w\lambda\) is also dominant, applying the same
assertion to \(w^{-1}\) and \(\mu\) puts \(\mu-\lambda\)
in that same pointed cone. Their sum is zero, so
\(\lambda=\mu\).
\end{proof}

\section{Dynkin classification}

\begin{proposition}[Cartan matrices and reconstruction]\label{res:E04}
Let \(a_{ij}=\langle\alpha_j,\alpha_i^\vee\rangle\).
Then \(a_{ii}=2\), \(a_{ij}\in\mathbb Z_{\le0}\) for
\(i\ne j\), and \(a_{ij}a_{ji}\in\{0,1,2,3\}\).
The matrix \(\operatorname{diag}((\alpha_i,\alpha_i))A\)
is symmetric positive definite. The Dynkin graph joins \(i,j\)
by \(a_{ij}a_{ji}\) edges, with its arrow pointing toward
the shorter root. It determines the root system up to
isometry and one positive scale on each irreducible component.
\end{proposition}
\begin{proof}
The angle and simple-root arguments already give the stated
entries and products. Symmetrization is twice the Gram matrix.
On a connected component the ratio of adjacent squared lengths
is \(a_{ji}/a_{ij}\), so the graph determines every length
relative to one chosen length. It therefore determines the
Gram matrix of a basis. The induced isometry identifies
the simple reflections, and Theorem~\ref{res:E02} then
identifies all roots as their orbits.
Orthogonal graph components give invariant orthogonal
subspaces and \(W\) is the product of their Weyl groups.
This also proves the uniqueness of the irreducible decomposition.
\end{proof}

\begin{theorem}[Exhaustion of finite Dynkin diagrams]\label{res:E05}
The connected finite reduced crystallographic diagrams are
\[
 A_n\ (n\ge1),\quad B_n,C_n\ (n\ge2),\quad D_n\ (n\ge4),
 \quad E_6,E_7,E_8,F_4,G_2,
\]
with \(B_2=C_2\) up to exchanging the two root lengths.
\end{theorem}
\begin{proof}
Normalize the simple roots to squared length \(2\).
Their symmetric Gram matrix \(M\) has diagonal entries \(2\)
and off-diagonal entries \(-\sqrt m\), where
\(m=a_{ij}a_{ji}\). Every principal submatrix is positive
definite. A cycle is impossible: on its vertices the vector
with all entries \(1\) has quadratic value at most zero.
A vertex with four neighbors is impossible: put \(2\) at
the center and \(1\) at its four neighbors; again the
quadratic value is at most zero.
Thus the graph is a tree of valence at most three.

There cannot be two trivalent vertices. Take the path joining
two consecutive ones and retain two neighboring leaves at
each end. Give all vertices on the path value \(2\) and
the four leaves value \(1\). For single edges its quadratic
value is zero; heavier edges can only decrease it.
A triple edge cannot have any attached vertex: the
three-vertex path with edge products \(3,1\) has determinant
\(8-2(3+1)=0\). Hence a triple edge gives just \(G_2\).

A trivalent vertex and a double edge cannot coexist.
Retain the path from that vertex to the nearest double edge,
two additional leaves at the trivalent end, and the vertex
just beyond the double edge. Eliminate the two leaves by
Schur complement: each subtracts \(1/2\) from the central
diagonal, leaving \(1\). Eliminating the last vertex across
the double edge subtracts \(2/2=1\) from the other end.
The remaining chain has endpoint diagonals \(1\), all interior
diagonals \(2\), and off-diagonals \(-1\), so its all-ones
vector has quadratic value zero. If the double edge is
incident to the trivalent vertex, these two endpoint
subtractions occur on the same diagonal, leaving the
one-by-one zero matrix instead. Two double edges are
excluded by the same endpoint elimination along
the path between them; if they share a vertex,
both subtractions again leave one zero diagonal.

It remains to classify simply laced trees with at most
one branch and chains with one double edge.
A simple chain is \(A_n\).
For a branch with arms of lengths \(a\le b\le c\),
the inverse endpoint entry of the length \(a\) chain
matrix is \(a/(a+1)\); this follows by the determinant
recurrence \(d_m=2d_{m-1}-d_{m-2}\), \(d_m=m+1\).
Its Schur complement at the branch is positive precisely when
\[
 \frac1{a+1}+\frac1{b+1}+\frac1{c+1}>1.
\]
Thus \(a=1\); if \(b=1\) one gets \(D_{c+3}\), and if
\(b=2\), one gets \(c=2,3,4\), namely \(E_6,E_7,E_8\).
No \(b\ge3\) is possible.
For a double edge separating chains of \(p,q\ge1\) vertices,
the determinant is
\[
 (p+1)(q+1)-2pq=2-(p-1)(q-1).
\]
Positivity gives \(p=1\) or \(q=1\), the diagrams \(B_n,C_n\),
or \(p=q=2\), the diagram \(F_4\).
The two orientations of \(F_4\) are interchanged by reversing
the chain. These cases exhaust all graphs and orientations.
Their actual root systems are constructed next, so the list
has neither missing cases nor merely formal candidates.
\end{proof}

\begin{theorem}[Coordinate constructions]\label{res:E06}
Every diagram in Theorem~\ref{res:E05} is realized by a root
system. The classical systems are
\[
\begin{array}{c|l|r}
\text{type}&\text{roots}&|R|\\\hline
A_n&\varepsilon_i-\varepsilon_j\text{ in }\sum x_i=0
 &(n+1)n\\
B_n&\pm\varepsilon_i,\ \pm\varepsilon_i\pm\varepsilon_j&2n^2\\
C_n&\pm2\varepsilon_i,\ \pm\varepsilon_i\pm\varepsilon_j&2n^2\\
D_n&\pm\varepsilon_i\pm\varepsilon_j&2n(n-1).
\end{array}
\]
The exceptional cardinalities are
\[
 |G_2|=12,\quad |F_4|=48,\quad
 |E_6|=72,\quad |E_7|=126,\quad |E_8|=240.
\]
\end{theorem}
\begin{proof}
Signed coordinate permutations verify the reflection axioms
for the classical lists. Their simple bases are
\(\varepsilon_i-\varepsilon_{i+1}\), with the final root
\(\varepsilon_n\) in \(B_n\), \(2\varepsilon_n\) in \(C_n\),
and with the two final roots
\(\varepsilon_{n-1}-\varepsilon_n,\varepsilon_{n-1}+\varepsilon_n\)
in \(D_n\). For \(A_n\) use \(n+1\) coordinates.
The dot products give the indicated graphs.

In \(\sum_{i=1}^3x_i=0\), take the \(G_2\) roots
\[
 \{\varepsilon_i-\varepsilon_j:i\ne j\}
 \ \cup\
 \{\pm(2\varepsilon_i-\varepsilon_j-\varepsilon_k):
                       \{i,j,k\}=\{1,2,3\}\}.
\]
A short reflection interchanges two coordinates.
The reflection in \(2\varepsilon_1-\varepsilon_2-\varepsilon_3\)
sends each short root to a short root and permutes the six
displayed long roots; inserting
\(x\mapsto x-(x,\alpha)\alpha/3\) verifies this directly.
The simple roots can be
\(\varepsilon_1-\varepsilon_2\) and
\(-2\varepsilon_1+\varepsilon_2+\varepsilon_3\),
with Cartan products \(3\).

In \(\mathbb R^4\), take
\[
 F_4=\{\pm\varepsilon_i\}\cup
 \{\pm\varepsilon_i\pm\varepsilon_j:i<j\}\cup
 \{\tfrac12(\pm\varepsilon_1\pm\varepsilon_2
                    \pm\varepsilon_3\pm\varepsilon_4)\}.
\]
The three parts have \(8,24,16\) members.
Signed permutations preserve the list. Reflection in a
half root preserves the long roots; on the short roots it
either negates that root, fixes it, or interchanges a
coordinate root with a half root. This follows by the
possible inner products \(0,\pm\tfrac12,\pm1\);
for long roots the nonzero inner product is \(\pm1\).
These computations prove both reflection invariance and
integrality. One simple basis is
\[
 \varepsilon_2-\varepsilon_3,\quad
 \varepsilon_3-\varepsilon_4,\quad
 \varepsilon_4,\quad
 \tfrac12(\varepsilon_1-\varepsilon_2-\varepsilon_3-\varepsilon_4).
\]
Its expansion coefficients for \(x\) are
\(x_1+x_2,\ 2x_1+x_2+x_3,\
3x_1+x_2+x_3+x_4,\ 2x_1\).
For every displayed root they are integers of one sign;
this also verifies directly that this is a simple basis.

For \(E_8\), let
\[
 L=\{z\in\mathbb Z^8:\textstyle\sum z_i\text{ even}\}
 \cup
 \{z+(\tfrac12,\ldots,\tfrac12):z\in\mathbb Z^8,\
                                      \textstyle\sum z_i\text{ even}\}.
\]
This is an integral even lattice. Its norm \(2\) vectors are
the \(112\) vectors \(\pm\varepsilon_i\pm\varepsilon_j\)
and the \(128\) half-coordinate vectors with an even
number of minus signs. For a norm \(2\) vector \(\alpha\),
the map \(x\mapsto x-(x,\alpha)\alpha\) preserves \(L\)
and norm. It therefore preserves precisely this root list.
Take the basis
\[
\begin{split}
 \alpha_1&=\tfrac12(\varepsilon_1+\varepsilon_8
                         -\varepsilon_2-\cdots-\varepsilon_7),\\
 \alpha_2&=\varepsilon_1+\varepsilon_2,\qquad
 \alpha_3=\varepsilon_2-\varepsilon_1,\\
 \alpha_4&=\varepsilon_3-\varepsilon_2,\quad
 \alpha_5=\varepsilon_4-\varepsilon_3,\quad
 \alpha_6=\varepsilon_5-\varepsilon_4,\quad
 \alpha_7=\varepsilon_6-\varepsilon_5,\quad
 \alpha_8=\varepsilon_7-\varepsilon_6 .
\end{split}
\]
Its graph has arms \(1,2,4\).
Here is an explicit way to check simplicity, without relying
only on a drawing. For \(x=\sum c_i\alpha_i\), recover
\[
\begin{aligned}
c_1&=2x_8,&c_8&=x_7+x_8,&c_7&=x_6+x_7+2x_8,\\
c_6&=x_5+x_6+x_7+3x_8,&
c_5&=x_4+x_5+x_6+x_7+4x_8,\\
c_4&=x_3+x_4+x_5+x_6+x_7+5x_8,&
c_2&=(x_1+x_2+c_4)/2,&c_3&=c_2+x_8-x_1 .
\end{aligned}
\]
For a half root the parity condition makes these integers,
all nonnegative when \(x_8=1/2\), and all nonpositive
when \(x_8=-1/2\). For an integral root, the same formulas
with its two nonzero coordinates give one sign throughout.
This proves simplicity and the \(E_8\) graph.

Intersect with \(x_7+x_8=0\) to obtain \(E_7\),
with basis \(\alpha_1,\ldots,\alpha_7\).
There are \(60+2\) integral roots and \(64\) half roots.
Intersect further with \(x_6=x_7=-x_8\) to obtain \(E_6\),
with basis \(\alpha_1,\ldots,\alpha_6\).
There are \(40\) integral roots and \(32\) half roots.
The intersections are root systems because the reflection
in a vector of the subspace preserves that subspace.
Their graphs have arms \(1,2,3\) and \(1,2,2\), respectively.
\end{proof}

\section{Lattices and useful dominant elements}

\begin{proposition}[Root and weight lattices]\label{res:E07}
Set
\[
 Q=\sum_i\mathbb Z\alpha_i,\quad
 Q^\vee=\sum_i\mathbb Z\alpha_i^\vee,\quad
 P=(Q^\vee)^*,\quad P^\vee=Q^*.
\]
Then \(Q\subset P\), \(Q^\vee\subset P^\vee\), and
the fundamental weights \(\omega_i\), characterized by
\(\langle\omega_i,\alpha_j^\vee\rangle=\delta_{ij}\),
form a basis of \(P\).
Dominant integral weights are precisely
\(\sum_i n_i\omega_i\), \(n_i\in\mathbb Z_{\ge0}\).
Each real orbit has a unique dominant representative.
Moreover \([P:Q]=\det A\).
\end{proposition}
\begin{proof}
The simple coroots are a basis of the coroot lattice:
the proof of Theorem~\ref{res:E02} applies to the dual
root system. Taking dual bases proves the assertion about
fundamental weights. Integral Cartan pairings put every root
in \(P\), and the dual argument gives the other inclusion.
The inequalities for the closed positive chamber are
exactly the nonnegativity of the fundamental-weight
coefficients. Orbit uniqueness is Theorem~\ref{res:E03}.
Finally the simple roots have coordinate columns \(A\)
in the fundamental-weight basis, so the index is the
absolute determinant, which is positive by symmetrization.
For example these determinants are \(n+1,2,2,4,3,2,1,1,1\)
for \(A_n,B_n,C_n,D_n,E_6,E_7,E_8,F_4,G_2\).
They can be computed by the chain recurrence and branch
Schur complements in Theorem~\ref{res:E05}.
\end{proof}

\begin{proposition}[The longest element and finite weight sets]\label{res:E08}
There is a unique \(w_0\in W\) carrying \(R^+\) to \(-R^+\).
It satisfies \(w_0^2=1\) and \(\ell(w_0)=|R^+|\).
The element
\[
 \rho=\tfrac12\sum_{\alpha>0}\alpha
\]
satisfies \(\langle\rho,\alpha_i^\vee\rangle=1\).
Similarly \(\rho^\vee=\frac12\sum_{\alpha>0}\alpha^\vee\)
satisfies \(\langle\alpha_i,\rho^\vee\rangle=1\).
For a finite-dimensional \(\mathfrak g\)-module its
\(\mathfrak h\)-weights are integral, their multiplicities
are \(W\)-invariant, and their set is saturated in the sense
that, if \(\mu\) occurs and \(n=\langle\mu,\alpha^\vee\rangle
\ge0\), then \(\mu-j\alpha\) occurs for \(0\le j\le n\).
\end{proposition}
\begin{proof}
The negative chamber is a chamber; simple transitivity gives
\(w_0\), and its square stabilizes the positive chamber.
The length formula proves maximality and the value stated.
A simple reflection permutes the positive roots other
than \(\alpha_i\), so \(s_i\rho=\rho-\alpha_i\).
Comparing with the reflection formula gives its pairing \(1\).
Apply the same argument to the dual system.

The action of \(\mathfrak h\) is diagonalizable on a finite
module by Theorem~\ref{res:C07}. Restriction to every root
triple and the \(\mathfrak{sl}_2\) classification give
integer coroot eigenvalues. The finite direct sum of
weight spaces in a coset \(\mu+\mathbb Z\alpha\)
decomposes into \(\mathfrak{sl}_2\) strings.
Each string is symmetric about zero and contains every
intermediate weight of the same parity. Thus the endpoints
are exchanged by \(s_\alpha\), with their multiplicities,
and a string containing the eigenvalue \(n\ge0\)
contains \(n,n-2,\ldots,-n\).
This proves both multiplicity invariance and saturation.
\end{proof}

\begin{theorem}[Highest roots and the principal triple]\label{res:E09}
For a simple \(\mathfrak g\), there is a unique highest root
\(\theta\), and every root is at most \(\theta\) in the
simple-root order. Choose simple triples \(e_i,f_i,h_i\).
If \(2\rho^\vee=\sum_i c_i h_i\), then
\[
 e=\sum_i e_i,\qquad h=2\rho^\vee,\qquad f=\sum_i c_i f_i
\]
form an \(\mathfrak{sl}_2\)-triple.
For \(m>0\), let \(r_m\) be the number of positive roots of
height \(m\), and let \(r_0=r\). Define the exponents by
giving \(m\) multiplicity \(r_m-r_{m+1}\). Then
\[
 \mathfrak g|_{\langle e,h,f\rangle}
       \simeq\bigoplus_{i=1}^r L_{2m_i},\qquad
 \sum_i m_i=|R^+|.
\]
Here \(L_a\) has highest weight \(a\) and dimension \(a+1\).
The Coxeter number used in this volume is
\(h_R=1+\operatorname{ht}\theta=1+\max m_i\).
\end{theorem}
\begin{proof}
Choose a maximal positive root \(\theta\).
Its root vector is killed by all positive root operators.
By PBW, the submodule of the adjoint module generated by
that vector is spanned by negative root monomials applied
to it: move positive factors right and Cartan factors
to the middle. All its weights are therefore at most
\(\theta\). Since the adjoint module of a simple algebra
is simple, this submodule is all of \(\mathfrak g\).
Thus every root is at most \(\theta\); a second maximal root
must equal it. This argument uses PBW, not the highest-weight
classification of Chapter~7.

Distinct simple positive and negative root spaces bracket
to zero, whereas \([e_i,f_i]=h_i\). Therefore
\([e,f]=h\). The identities
\(\alpha_i(2\rho^\vee)=2\) give the other two triple
relations. In the adjoint module the \(h\)-weights are
twice the root heights, with zero space \(\mathfrak h\).
All simple \(\mathfrak{sl}_2\)-summands consequently have
even highest weight. Each contributes one dimension to
zero weight and one to each positive even weight up to
its highest weight. Hence there are \(r\) summands and
the number of summands with highest weight exactly \(2m\)
is \(r_m-r_{m+1}\). Summing their contributions to positive
weights gives \(\sum m_i=|R^+|\).
There are no trivial summands since \(r_0=r_1=r\).
The top nonzero weight is \(2\operatorname{ht}\theta\),
proving the assertion about the largest exponent.
\end{proof}

For reference the complete height calculation is summarized here:
\[
\begin{array}{c|l|c}
 R&\{m_i\}&h_R\\ \hline
 A_n&1,2,\ldots,n&n+1\\
 B_n,C_n&1,3,\ldots,2n-1&2n\\
 D_n&1,3,\ldots,2n-3,\ n-1&2n-2\\
 G_2&1,5&6\\
 F_4&1,5,7,11&12\\
 E_6&1,4,5,7,8,11&12\\
 E_7&1,5,7,9,11,13,17&18\\
 E_8&1,7,11,13,17,19,23,29&30
\end{array}
\]
These entries can be checked without a representation formula:
use the coordinate lists of Theorem~\ref{res:E06}, expand each
positive root in the indicated basis, and form
\(\sum_{\alpha>0}t^{\operatorname{ht}\alpha}\).
For each row it equals
\(\sum_{m\text{ in the middle column}}(t+\cdots+t^m)\);
the explicit coefficient-recovery formulas there also cover
the exceptional systems. In type \(D_n\), the exponent \(n-1\)
has multiplicity two when \(n\) is even.
The definition of \(h_R\) here does not silently identify it
with the order of a product of simple reflections.

The diagram and lattice arguments follow
\cite[Sections 9--13]{humphreys} and
\cite[Sections 21--23 and 32]{etingof}.

\chapter{Reconstruction and Classification}
\label{ch:classification}

\section{Generators, relations, and the nonzero construction}

Throughout this chapter \(A=(a_{ij})\) is the Cartan matrix of
a finite reduced root system, with
\(a_{ij}=\langle\alpha_j,\alpha_i^\vee\rangle\).
We establish the existence theorem before using its presentation
to prove uniqueness. Thus the reconstruction result labelled F01
below is a consequence of the full construction, rather than an
assumption used to prove it.

\begin{proposition}[The Serre relations and the triangular construction]
\label{res:F02}
In a semisimple algebra with simple-root triples the elements
\(e_i,f_i,h_i\) generate the algebra and satisfy
\[
\begin{gathered}
[h_i,h_j]=0,\quad [h_i,e_j]=a_{ij}e_j,\quad
[h_i,f_j]=-a_{ij}f_j,\quad [e_i,f_j]=\delta_{ij}h_i,\\
(\operatorname{ad}e_i)^{1-a_{ij}}e_j=0,\qquad
(\operatorname{ad}f_i)^{1-a_{ij}}f_j=0\quad(i\ne j).
\end{gathered}\tag{6.1}\label{eq:serre-relations}
\]
Conversely form the free Lie algebra on these symbols and impose
these relations. The resulting algebra \(\mathfrak g(A)\)
has a vector space decomposition
\[
 \mathfrak g(A)=\mathfrak n^-\oplus\mathfrak h\oplus\mathfrak n^+,
\]
where \(\mathfrak h\) has basis \(h_i\), and
\(\mathfrak n^\pm\) have precisely their respective positive
or negative Serre relations. In particular no \(e_i,f_i,h_i\)
has been killed by the presentation.
\end{proposition}
\begin{proof}
In an existing algebra, a positive root \(\beta\) of height
greater than one has \(\beta-\alpha_i>0\) for some \(i\).
The nonzero root-string bracket
\([e_i,\mathfrak g_{\beta-\alpha_i}]=\mathfrak g_\beta\)
gives generation by induction. Negative roots work the same way,
and the \(h_i\) form a basis of \(\mathfrak h\).
For \(i\ne j\), \(f_j\) is a highest vector of weight
\(-a_{ij}\) for the \(i\)-th triple. The
\(\mathfrak{sl}_2\) calculation gives the second Serre relation;
the first follows by reversing signs. The other relations
are the root-space bracket calculations.

We now justify carefully that the abstract presentation is
nonzero. Initially omit the Serre relations and call the algebra
\(\widetilde{\mathfrak g}\). Repeated use of
\([e_i,f_j]=\delta_{ij}h_i\), followed by Jacobi, moves any
mixed bracket into a sum of brackets using only \(e\)'s,
only \(f\)'s, or one \(h\). For example apply
\([[u,v],w]=[u,[v,w]]-[v,[u,w]]\) to reduce the number
of mixed letters adjacent to an innermost bracket.
The root-lattice grading separates positive degree, zero,
and negative degree. Hence
\(\widetilde{\mathfrak g}=
\widetilde{\mathfrak n}^-\oplus\widetilde{\mathfrak h}
\oplus\widetilde{\mathfrak n}^+\) once freeness and
independence in these pieces are verified.

For that verification let \(F\) be the free Lie algebra on
symbols \(f'_i\), and act on
\[
 V=T(f'_1,\ldots,f'_r)\otimes\mathbb C[H_1,\ldots,H_r].
\]
If a word \(w\) has weight \(-\gamma\), put
\[
 F_i(wP)=f'_iwP,\qquad
 H_i^{\rm op}(wP)=w(H_i-\gamma(h_i))P.
\]
Define \(E_i(1P)=0\) and recursively
\[
 E_i(f'_j wP)=f'_jE_i(wP)+\delta_{ij}H_i^{\rm op}(wP).
 \tag{6.2}
\]
The recursion gives \([E_i,F_j]=\delta_{ij}H_i^{\rm op}\).
The degree of \(F_j\) changes the weight by \(-\alpha_j\),
so \([H_i^{\rm op},F_j]=-a_{ij}F_j\).
Every summand of \(E_jwP\) has weight
\(-\gamma+\alpha_j\); comparing the polynomial factors
therefore gives \([H_i^{\rm op},E_j]=a_{ij}E_j\).
The \(H_i^{\rm op}\) commute. These are all the imposed
relations of \(\widetilde{\mathfrak g}\), so this is an action.
Applying a bracket of the \(F_i\)'s to \(1\) gives the
identical free-Lie bracket inside \(T(f')=U(F)\);
PBW embeds \(F\) there. Thus the negative piece is free.
Also \(H_i^{\rm op}(1)=H_i\) proves independence of the \(h_i\).
Interchanging \(e,f\) and replacing \(h\) by \(-h\)
proves the same freeness for the positive piece.

Put \(S^+_{ij}=(\operatorname{ad}e_i)^{1-a_{ij}}e_j\).
For \(k\ne i,j\), \([f_k,S^+_{ij}]=0\) follows at once.
For \(k=i\), the usual iterated \(\mathfrak{sl}_2\)
commutator has scalar
\((1-a_{ij})(a_{ij}+1-a_{ij}-1)=0\).
For \(k=j\), it is a multiple of
\((\operatorname{ad}e_i)^{1-a_{ij}}h_j\).
If \(a_{ij}=0\), this vanishes since \(a_{ji}=0\);
otherwise the exponent is at least two and the second
bracket of \(e_i\) with \(h_j\) is already zero.
Thus all \(f_k\) annihilate \(S^+_{ij}\).
The ideal \(I^+\) generated by the \(S^+_{ij}\) within the
free positive piece is stable under all \(f_k\): apply Jacobi
inductively to a bracket of positive generators with an
\(S^+_{ij}\). The terms involving a bracket with \(f_k\)
have fewer positive generators; Cartan terms only multiply
homogeneous vectors by scalars. It is also stable under
\(\mathfrak h\). The corresponding statement holds for
the negative ideal \(I^-\).
Therefore the full ideal of Serre relations is exactly
\(I^-\oplus I^+\). It has no Cartan component and no
component of degree \(\pm\alpha_i\).
Quotienting gives precisely the decomposition asserted.
\end{proof}

\begin{theorem}[Finite type and simplicity of the Serre algebra]
\label{res:F03}
The algebra \(\mathfrak g(A)\) is finite dimensional. Its nonzero
weights are exactly the prescribed roots, each with multiplicity
one, and its zero space is the \(r\)-dimensional \(\mathfrak h\).
For connected \(A\) it is simple. For arbitrary \(A\) it is
the direct sum of the simple algebras of the connected components.
\end{theorem}
\begin{proof}
Fix \(i\). The \(i\)-th triple acts locally finitely on every
generator: on itself it is the adjoint three-dimensional module;
on \(e_j\) or \(f_j\), \(j\ne i\), the Serre relation bounds
the raising or lowering string; a Cartan generator belongs
to the sum of a trivial module and that adjoint module.
If \(x,y\) lie in finite-dimensional invariant spaces, their
brackets lie in the image of the tensor product of those
spaces. Induction on bracket length shows that every element
is in a finite-dimensional module for the \(i\)-th triple.
Consequently the finite exponentials
\[
 n_i=\exp(\operatorname{ad}e_i)
       \exp(-\operatorname{ad}f_i)
       \exp(\operatorname{ad}e_i)
\]
are well-defined automorphisms. Their effect on \(\mathfrak h\)
is \(h\mapsto h-\alpha_i(h)h_i\), so they take the weight
space of \(\gamma\) to that of \(s_i\gamma\).

The triangular decomposition says that a nonzero weight
belongs to \(Q^+\) or \(-Q^+\), and that the zero space is
exactly \(\mathfrak h\). Each individual weight space is
finite dimensional, since only finitely many words have
a specified nonnegative multidegree. Suppose a positive
weight \(\gamma=\sum k_j\alpha_j\) occurs.
If its support has one element, it is a simple root:
a free Lie algebra on one generator has no degree greater
than one. Otherwise positivity of the Euclidean form gives
an \(i\) with \(\langle\gamma,\alpha_i^\vee\rangle>0\).
The weight \(s_i\gamma\) occurs and has lower height.
At least one coefficient other than the \(i\)-th remains
positive; hence it cannot be in \(-Q^+\).
It must be in \(Q^+\), and induction on height takes it
to a simple root. Thus \(\gamma\) is a root of the
given root system. The negative case is identical.
Conversely the \(n_i\)'s and the nonzero simple generators
produce all \(W\Pi=R\).
The weight spaces of simple roots are one dimensional;
the reflection isomorphisms make all root spaces
one dimensional. It follows that
\(\dim\mathfrak g(A)=r+|R|\).

Suppose first that the diagram is connected and
\(0\ne I\triangleleft\mathfrak g(A)\).
The commuting diagonal actions of the \(h_i\) decompose \(I\)
into weight spaces. If \(I\) has only a nonzero Cartan
element \(h\), some \(\alpha_i(h)\ne0\), since \(A\)
is invertible; hence \([h,e_i]\in I\) supplies a nonzero
root vector. A root vector in \(I\) can be reflected by
the \(n_i\)'s, which preserve ideals, to some \(e_j\).
Then \(h_j=[e_j,f_j]\in I\), and
\([h_j,f_j]=-2f_j\in I\). Whenever \(a_{jk}\ne0\),
\([h_j,e_k]=a_{jk}e_k\) puts \(e_k\) in \(I\).
Connectedness propagates this to every generator.
Thus \(I=\mathfrak g(A)\), proving simplicity.
For disconnected diagrams, generators in different components
commute by the relations, and the independent Cartan and
root decomposition makes their sum direct.
\end{proof}

\begin{theorem}[Reconstruction and the isomorphism theorem]\label{res:F01}
An isomorphism of the root systems of two complex semisimple
Lie algebras extends, after choosing compatible simple-root
triples, to a Lie algebra isomorphism. In particular the root
system determines the algebra.
\end{theorem}
\begin{proof}
For an algebra with root system \(R\), the generation and
relations in Proposition~\ref{res:F02} give a surjective map
\(\mathfrak g(A)\to\mathfrak g\) sending the indicated
generators to the selected triples. Both sides have dimension
\(r+|R|\), by Theorem~\ref{res:F03} and the root decomposition.
Thus this map is an isomorphism.
Apply the same construction to the other algebra, using
the root-system isomorphism to identify the ordered simple
roots and hence their Cartan matrix. Composing the two
isomorphisms gives the required extension.
The choices of simple root vectors are compatible because
each negative vector can be normalized to have bracket
equal to its coroot.
\end{proof}

\begin{corollary}[Complex semisimple classification]\label{res:F04}
The complex simple Lie algebras are the algebras of types
\[
 A_n,\ B_n,\ C_n,\ D_n,\ E_6,E_7,E_8,F_4,G_2
\]
with the rank restrictions and repetitions specified in
Theorem~\ref{res:E05}. Every complex semisimple Lie algebra
is a direct sum of these, uniquely up to permutation.
\end{corollary}
\begin{proof}
Every semisimple algebra gives a root system by
Proposition~\ref{res:D05}; its irreducible components
are the finite list in Theorem~\ref{res:E05}.
The coordinate models and Serre construction prove that
each component actually occurs, and reconstruction proves
uniqueness. The simple-ideal decomposition
of Theorem~\ref{res:B06} agrees with the root-system
decomposition: different ideals have orthogonal Killing
forms and disjoint root systems, and an irreducible
component yields a simple Serre algebra.
For \(A,B,C,D\), Example~\ref{res:D08} identifies these
algebras with the stated classical matrix algebras.
\end{proof}

\section{Integral structure constants}

\begin{theorem}[Chevalley bases]\label{res:F05}
There are nonzero root vectors \(x_\alpha\) such that
\[
 [x_\alpha,x_{-\alpha}]=h_\alpha,\qquad
 [x_\alpha,x_\beta]=N_{\alpha\beta}x_{\alpha+\beta},
\]
and
\[
 N_{-\alpha,-\beta}=-N_{\alpha\beta},\qquad
 N_{\alpha\beta}=\pm(p+1)
\]
when \(\alpha+\beta\) is a root and the \(\alpha\)-string
through \(\beta\) begins at \(\beta-p\alpha\).
Together with the simple coroots these vectors are a basis
whose structure constants are integers.
\end{theorem}
\begin{proof}
The Serre relations admit the involution
\[
 \sigma(e_i)=-f_i,\quad \sigma(f_i)=-e_i,\quad
 \sigma(h_i)=-h_i.
\]
Substitution in (6.1) verifies every relation, and its square
is the identity. Hence it carries
\(\mathfrak g_\alpha\) to \(\mathfrak g_{-\alpha}\).
For each pair \(\{\alpha,-\alpha\}\), choose \(x_\alpha\)
and set \(x_{-\alpha}=-\sigma x_\alpha\).
Opposite root spaces pair perfectly, so their bracket is
a nonzero multiple of \(h_\alpha\). Multiplying both vectors
by a suitable complex scalar changes that multiple by its
square; choose it to obtain bracket \(h_\alpha\).
Applying \(\sigma\) to a bracket immediately gives
\(N_{-\alpha,-\beta}=-N_{\alpha\beta}\).

We verify the magnitude, since a sign choice alone would
not prove integrality. In a root string
\(\beta-p\alpha,\ldots,\beta+q\alpha\),
the \(\mathfrak{sl}_2\) calculation gives
\[
 [x_{-\alpha},[x_\alpha,x_\beta]]
       =q(p+1)x_\beta .
 \tag{6.3}
\]
Apply Jacobi twice to
\([[x_\alpha,x_\beta],[x_{-\alpha},x_{-\beta}]]\).
Using (6.3) for both strings, and the negative-constant
identity, gives
\[
 N_{\alpha\beta}^2h_{\alpha+\beta}
       =q'(p'+1)h_\alpha+q(p+1)h_\beta,
 \tag{6.4}
\]
where \(p',q'\) describe the \(\beta\)-string through
\(\alpha\). Compare the coefficient of \(t_\beta\),
using \(h_\gamma=2t_\gamma/(\gamma,\gamma)\):
\[
 N_{\alpha\beta}^2
   =q(p+1)\frac{(\alpha+\beta,\alpha+\beta)}{(\beta,\beta)}.
 \tag{6.5}
\]
The rank two lists give the following complete possibilities
when \(q>0\):
\[
\begin{array}{c|rrrrrr}
(p,q)&(0,1)&(0,2)&(0,3)&(1,1)&(1,2)&(2,1)\\
(\alpha+\beta)^2/\beta^2&1&1/2&1/3&2&1&3 .
\end{array}
\]
In each case the right side of (6.5) is \((p+1)^2\).
Thus \(N_{\alpha\beta}=\pm(p+1)\).
If \(\alpha+\beta\) is neither zero nor a root the bracket
is zero; opposite roots have the normalized coroot bracket.
Finally every coroot is an integral combination of the simple
coroots, by applying the simple-root theorem to the dual root
system. The remaining Cartan brackets have their integral
Cartan entries. All structure constants are therefore integers.

The construction chooses all signs simultaneously through
one involution; it does not attempt to assign pairwise signs
independently. Under \(x_\alpha\mapsto c_\alpha x_\alpha\)
the opposite normalization requires \(c_\alpha c_{-\alpha}=1\),
and a bracket constant changes by
\(c_\alpha c_\beta/c_{\alpha+\beta}\).
These equations specify precisely the comparison of normalizations.
\end{proof}

\section{Borel subalgebras without an omitted fixed-point argument}

We include the two analytic ingredients used below, so that
the conjugacy theorem does not depend on a later classification
of compact groups.

\begin{lemma}[A matrix-group fact]\label{supp:closed-matrix-groups}
A closed subgroup \(H\subset\operatorname{GL}(V_{\mathbb R})\)
is a Lie group with Lie algebra
\(\mathfrak h=\{X:\exp(tX)\in H\text{ for all real }t\}\).
Its identity component is generated by exponentials from
\(\mathfrak h\).
\end{lemma}
\begin{proof}
The product limits
\[
 (e^{tX/n}e^{tY/n})^n\longrightarrow e^{t(X+Y)},\qquad
 (e^{sX}e^{sY}e^{-sX}e^{-sY})^{n}
       \longrightarrow e^{t[X,Y]},\quad s=\sqrt{t/n},
\]
first for \(t\ge0\) and then using inverses, follow by
expanding the matrix products through orders \(1/n\).
They prove that \(\mathfrak h\) is a vector space and
closed under brackets; closure of \(H\) is used in passing
to the limits. It is also closed as a set.
Choose a vector-space complement \(M\) to \(\mathfrak h\).
The map \((X,Y)\mapsto e^Xe^Y\) from
\(\mathfrak h\oplus M\) has invertible differential at zero,
and hence is a local coordinate system near the identity.
If there were arbitrarily small nonzero \(Y_j\in M\)
with \(e^{Y_j}\in H\), a subsequence of
\(Y_j/\|Y_j\|\) would approach a unit \(Y\in M\).
Powers with exponent \(\lfloor t/\|Y_j\|\rfloor\)
would give \(e^{tY}\in H\) for every \(t\), a contradiction.
Thus in a sufficiently small coordinate neighborhood an
element \(e^Xe^Y\in H\) has \(Y=0\).
This gives the submanifold chart \(\exp(\mathfrak h)\)
and the asserted tangent space. Translates give the global
Lie structure. The subgroup generated by these local
exponentials is open in \(H\); its cosets are also open,
so it is exactly the identity component.
\end{proof}

\begin{lemma}[The compact form used for the flag orbit]
\label{supp:compact-form-flag}
There is a conjugate-linear involution \(\tau\) of
\(\mathfrak g\) whose fixed algebra \(\mathfrak u\)
has negative definite Killing form. For each positive system
and \(\mathfrak b=\mathfrak h\oplus\mathfrak n^+\),
\[
 \mathfrak g_{\mathbb R}=\mathfrak u+\mathfrak b_{\mathbb R}.
\]
The inner orbit of \(\mathfrak b\) in the Grassmannian of
\(\dim\mathfrak b\) dimensional subspaces is compact.
\end{lemma}
\begin{proof}
Use the real span of the Chevalley basis and define
\(\tau(h_i)=-h_i\), \(\tau(x_\alpha)=-x_{-\alpha}\),
conjugating scalars. The real structure constants and their
negative-constant identity prove that \(\tau\) preserves
brackets and has square one.
Its fixed real basis consists of
\[
 ih_i,\qquad x_\alpha-x_{-\alpha},\qquad
 i(x_\alpha+x_{-\alpha})\quad(\alpha>0).
 \tag{6.6}
\]
On the real coroot space \(B\) is positive definite,
and \(B(x_\alpha,x_{-\alpha})=2/(\alpha,\alpha)>0\).
Orthogonality of root spaces makes (6.6) a negative
definite decomposition. In particular
\((v,w)_{\rm H}=-B(v,\tau w)\) is a positive definite
Hermitian form on \(\mathfrak g\).
The displayed real sum follows by writing a negative root
vector as the sum of an element from (6.6) and a positive
root vector, and by splitting Cartan coefficients into
real and imaginary parts.

Let \(G=\operatorname{Aut}(\mathfrak g)^0\).
It is a closed matrix group: preserving each bracket is
a finite set of polynomial equations in the matrix entries.
The matrix-group lemma applies, and its Lie algebra is
\(\operatorname{Der}(\mathfrak g)=\operatorname{ad}\mathfrak g\),
by Theorem~\ref{res:B07}. Thus it is the inner group used
in Chapter~4.
Let \(K\) be the identity component of the automorphisms
preserving the positive Hermitian form. It is compact.
Its Lie algebra is \(\operatorname{ad}\mathfrak u\):
the adjoint of \(\operatorname{ad}z\) is
\(-\operatorname{ad}(\tau z)\), so skew adjointness is
equivalent, by the zero center, to \(\tau z=z\).
The tangent to the \(G\)-orbit at \(\mathfrak b\) is the
image of \(\mathfrak g_{\mathbb R}\) in the tangent
space of the Grassmannian, with kernel its real normalizer.
The sum \(\mathfrak g_{\mathbb R}=\mathfrak u+
\mathfrak b_{\mathbb R}\) shows that \(K\)'s orbit has
the same tangent space. The inverse function theorem
therefore makes this compact orbit open in the \(G\)-orbit.
It is closed there by compactness. The \(G\)-orbit is
connected, so both orbits are equal. This proves compactness.
\end{proof}

\begin{lemma}[Solvable fixed points on a compact projective set]
\label{supp:solvable-projective-fixed}
Let a complex solvable Lie algebra \(\mathfrak a\) act on a
finite-dimensional space \(V\). Every nonempty compact subset
\(X\subset\mathbb P(V)\), invariant under
\(\exp(zx)\) for every \(x\in\mathfrak a,z\in\mathbb C\),
contains a point fixed by \(\mathfrak a\).
\end{lemma}
\begin{proof}
Induct on \(\dim\mathfrak a\). Choose a codimension one ideal
\(\mathfrak a'\), which exists because a nonzero solvable
algebra has a proper derived algebra.
By induction its fixed-point subset \(Y\subset X\)
is nonempty. It is closed, hence compact, and is invariant
under \(\mathfrak a\): conjugation by \(e^{zx}\) preserves
\(\mathfrak a'\). Choose \(a\) projecting to a generator
of \(\mathfrak a/\mathfrak a'\) and a point \([v]\in Y\).
Multiply \(a\) by a nonzero complex scalar so that its
distinct eigenvalues have distinct real parts.
In the Jordan decomposition of \(v\), take the eigenvalue
of largest real part present, and, within it, the highest
nonzero nilpotent power. As \(t\to+\infty\), the line
\([e^{ta}v]\) tends to the line of that highest-power
vector. This is an eigenline: lower real parts disappear
exponentially and lower powers disappear after division
by the leading power of \(t\). The limit belongs to \(Y\)
and is fixed by \(a\). It was already fixed by
\(\mathfrak a'\), so is fixed by \(\mathfrak a\).
\end{proof}

\begin{theorem}[Borel subalgebras]\label{res:F06}
The algebra \(\mathfrak b=\mathfrak h\oplus
\bigoplus_{\alpha>0}\mathfrak g_\alpha\) is maximal solvable.
Every solvable subalgebra lies in an inner conjugate of
\(\mathfrak b\); consequently all maximal solvable
subalgebras, called Borel subalgebras, are conjugate.
Moreover \(N_{\mathfrak g}(\mathfrak b)=\mathfrak b\).
\end{theorem}
\begin{proof}
Brackets increase positive height, so \(\mathfrak n^+\)
is nilpotent. The quotient \(\mathfrak b/\mathfrak n^+\)
is abelian; thus \(\mathfrak b\) is solvable.
A subalgebra strictly containing it is \(\mathfrak h\)-stable,
and hence contains some negative root space.
It then contains that root's full \(\mathfrak{sl}_2\),
so cannot be solvable. The same weight argument shows that
an element normalizing \(\mathfrak b\) has no negative
root component, proving its stated normalizer.

Embed the Grassmannian into
\(\mathbb P(\bigwedge^d\mathfrak g)\), \(d=\dim\mathfrak b\),
by \(L\mapsto[\bigwedge^d L]\).
The inner orbit \(X\) of \(\mathfrak b\) is nonempty and
compact by Lemma~\ref{supp:compact-form-flag}.
A solvable subalgebra \(\mathfrak a\) acts on this
exterior power, preserves \(X\), and the preceding fixed-point
lemma supplies an \(\mathfrak a\)-fixed line belonging to
some conjugate \(\mathfrak b'\).
A line \(\bigwedge^d\mathfrak b'\) is fixed precisely when
\([\mathfrak a,\mathfrak b']\subset\mathfrak b'\):
the infinitesimal action modulo \(\mathfrak b'\) is
the corresponding tangent map of the Grassmannian.
Thus \(\mathfrak a\subset N_{\mathfrak g}(\mathfrak b')
=\mathfrak b'\). If \(\mathfrak a\) was maximal solvable,
this inclusion is equality.
\end{proof}

\begin{theorem}[Parabolic subalgebras]\label{res:F07}
Subalgebras containing \(\mathfrak b\) are in bijection with
subsets \(S\subset\Pi\):
\[
 \mathfrak p_S=
 \mathfrak b\oplus\bigoplus_{\alpha\in R_S^+}
                       \mathfrak g_{-\alpha},
 \qquad R_S=R\cap\operatorname{span}_{\mathbb R}S .
\]
They are the standard parabolic subalgebras.
Every subalgebra containing a Borel is inner conjugate
to one of these. There is a decomposition
\[
 \mathfrak p_S=\mathfrak l_S\ltimes\mathfrak u_S,\quad
 \mathfrak l_S=\mathfrak h\oplus
                 \bigoplus_{\alpha\in R_S}\mathfrak g_\alpha,
 \quad
 \mathfrak u_S=\bigoplus_{\alpha\in R^+\setminus R_S}
                                       \mathfrak g_\alpha.
\]
Here \(\mathfrak l_S\) is reductive,
\(\mathfrak u_S\) is the nilpotent radical,
and replacing positive by negative gives the opposite
parabolic with the same Levi subalgebra.
\end{theorem}
\begin{proof}
Root addition verifies that \(\mathfrak p_S\) is a subalgebra
and \(\mathfrak u_S\) an ideal; height makes the latter
nilpotent. The generators from \(S\), their negatives,
and their coroots give the semisimple Serre algebra
for that root subsystem. Its complementary Cartan space
\(\{h:\alpha(h)=0\ \forall\alpha\in S\}\)
is central in \(\mathfrak l_S\). Thus \(\mathfrak l_S\)
is their direct sum and is reductive.

Conversely let \(\mathfrak p\supset\mathfrak b\).
It is a sum of \(\mathfrak h\) and root spaces, and define
\(S=\{\alpha_i:\mathfrak g_{-\alpha_i}\subset\mathfrak p\}\).
If \(\mathfrak g_{-\beta}\subset\mathfrak p\) for a nonsimple
positive root \(\beta\), choose \(\alpha_i\) with
\(\gamma=\beta-\alpha_i>0\) a root.
The two nonzero brackets
\[
 [\mathfrak g_{\alpha_i},\mathfrak g_{-\beta}]
       =\mathfrak g_{-\gamma},\qquad
 [\mathfrak g_\gamma,\mathfrak g_{-\beta}]
       =\mathfrak g_{-\alpha_i}
\]
belong to \(\mathfrak p\).
Induction on height puts every simple root in the support
of \(\gamma\), and also \(\alpha_i\), in \(S\).
Thus \(\beta\in R_S\).
The converse inclusion follows by generation of the
negative subsystem from its simple roots.
This proves \(\mathfrak p=\mathfrak p_S\) and uniqueness of \(S\).
Borel conjugacy proves the assertion about arbitrary parabolics.

To verify that the displayed ideal is the entire nilpotent
radical, let \(I\) be a nilpotent ideal of \(\mathfrak p_S\).
Its image in the reductive quotient lies in
\(\mathfrak z(\mathfrak l_S)\); any ideal with a nonzero
simple-factor image is not nilpotent.
As \(I\) is \(\mathfrak h\)-stable, a nonzero zero-weight
part of that image is represented by
\(0\ne z\in I\cap\mathfrak z(\mathfrak l_S)\).
Some root outside \(R_S\) has \(\alpha(z)\ne0\),
because all roots span and the roots in \(S\) kill \(z\).
Replacing it by its positive sign, its root vector
\([z,x_\alpha]/\alpha(z)\) belongs to \(I\).
But then \(\operatorname{ad}_I z\) has a nonzero eigenvalue,
contradicting nilpotence of \(I\).
Thus its quotient image is zero and \(I\subset\mathfrak u_S\).
The root formulas with negative roots give the opposite
parabolic and the common Levi factor.
\end{proof}

The Serre proof is based on
\cite[Sections 18--19]{humphreys} and
\cite[Section 24]{etingof}, including the free negative-word
module that prevents a degenerate presentation.
The integral normalization is
\cite[Section 25]{humphreys}. The last section supplies the
fixed-point and compact-orbit arguments behind the Borel
and parabolic assertions in \cite[Sections 51.3--51.5]{etingof}.

\chapter[Highest Weights and Finite-Dimensional Modules]{Highest Weights\\and Finite-Dimensional Modules}
\label{ch:highest-weights}

In this chapter \(\mathfrak g\) is a complex semisimple Lie algebra.
Fix a Cartan subalgebra \(\mathfrak h\), positive roots \(R^+\),
simple roots \(\alpha_1,\ldots,\alpha_r\), and normalized triples
\((e_i,f_i,h_i)\). Put
\[
 Q_+=\sum_i\mathbb Z_{\geq0}\alpha_i,\qquad
 P=\{\lambda\in\mathfrak h^*:\lambda(h_i)\in\mathbb Z\},\qquad
 P_+=\{\lambda\in P:\lambda(h_i)\geq0\}.
\]
The order \(\mu\leq\lambda\) means \(\lambda-\mu\in Q_+\).
The weight space in a module \(V\) is
\(V_\mu=\{v:hv=\mu(h)v\text{ for all }h\in\mathfrak h\}\).
The root-space results give
\(\mathfrak g_\alpha V_\mu\subseteq V_{\mu+\alpha}\).
Our arguments do not use the integration of a representation to a
Lie group. The sources are \cite[Sections 20--21]{humphreys} and
\cite[Section 25]{etingof}.

\section{Induction and the universal highest-weight module}

\begin{proposition}[Triangular multiplication and Verma modules]
\label{res:G01}
Multiplication is a vector-space isomorphism
\[
 U(\mathfrak n^-)\otimes U(\mathfrak h)\otimes U(\mathfrak n^+)
 \longrightarrow U(\mathfrak g).
\]
For \(\lambda\in\mathfrak h^*\), let \(\mathbb C_\lambda\) be the
one-dimensional left module over
\(\mathfrak b=\mathfrak h\oplus\mathfrak n^+\), with \(\mathfrak n^+\)
acting by zero and \(h\) acting by \(\lambda(h)\). Then
\[
 M(\lambda)=U(\mathfrak g)\otimes_{U(\mathfrak b)}\mathbb C_\lambda
 \simeq U(\mathfrak g)/
 \left(\sum_{h\in\mathfrak h}U(\mathfrak g)(h-\lambda(h))
       +\sum_iU(\mathfrak g)e_i\right).
\]
The map \(u\mapsto uv_\lambda\) identifies \(U(\mathfrak n^-)\)
linearly with \(M(\lambda)\), where \(v_\lambda=1\otimes1\).
\end{proposition}
\begin{proof}
Choose ordered bases of the negative root spaces, \(\mathfrak h\),
and the positive root spaces, in that order. PBW
(Theorem~\ref{res:A05}) says exactly that multiplication sends the
tensor product of their ordered monomial bases bijectively to the
ordered monomial basis of \(U(\mathfrak g)\).
In particular \(U(\mathfrak g)\) is free as a right
\(U(\mathfrak b)\)-module on the negative-root monomials.

The prescribed action on \(\mathbb C_\lambda\) is a Lie action:
\([\mathfrak b,\mathfrak b]\subseteq\mathfrak n^+\), so its image
vanishes. The induced associative algebra character has kernel the
left ideal of \(U(\mathfrak b)\) generated by
\(h-\lambda(h)\) and \(e_i\). To check this description, commute all
positive generators to the right in PBW order. Every positive-degree
monomial in the positive generators lies in their generated left ideal;
the remaining Cartan polynomial reduces to its value at \(\lambda\).
The simple \(e_i\) generate \(\mathfrak n^+\), so they suffice.
Tensoring the right-free decomposition with that one-dimensional
quotient proves both identifications and the final assertion.
\end{proof}

\begin{proposition}[Weights and the universal property]
\label{res:G02}
The module \(M(\lambda)\) is a direct sum of finite-dimensional weight
spaces, has weights precisely \(\lambda-Q_+\), and has
\(M(\lambda)_\lambda=\mathbb C v_\lambda\). A vector \(v\) in any
\(\mathfrak g\)-module satisfying
\[
 hv=\lambda(h)v,\qquad e_iv=0\quad(1\leq i\leq r)
\]
determines a unique homomorphism \(M(\lambda)\to V\) taking
\(v_\lambda\) to \(v\). Every module generated by such a nonzero
vector is a quotient of \(M(\lambda)\).
\end{proposition}
\begin{proof}
A negative-root PBW monomial
\(\prod_{\alpha>0}f_\alpha^{a_\alpha}v_\lambda\) has weight
\(\lambda-\sum_{\alpha>0}a_\alpha\alpha\).
If the latter sum has height \(n\), each \(a_\alpha\leq n\);
there are only finitely many choices. This proves finite weight
spaces. Every element of \(Q_+\) can be obtained using the simple
negative root vectors, and their PBW monomials are nonzero, which
proves the assertion about the weight set. The empty monomial is
the sole one of weight \(\lambda\).
The map \(u\mapsto uv\) kills the displayed ideal in
Proposition~\ref{res:G01}, hence factors through \(M(\lambda)\).
Its value on the generator forces uniqueness and its image is
the submodule generated by \(v\).
\end{proof}

\begin{lemma}[Submodules of a weight module]
\label{supp:weight-submodules}
If \(V=\bigoplus_\mu V_\mu\) and \(W\subseteq V\) is
\(\mathfrak h\)-stable, then
\(W=\bigoplus_\mu(W\cap V_\mu)\), and \(V/W\) has the corresponding
quotient weight spaces.
\end{lemma}
\begin{proof}
A vector has finitely many weight components. Choose \(h\in\mathfrak h\)
which distinguishes those finitely many weights. Lagrange interpolation
in the operator \(h\) extracts each component of the vector. If the
vector lies in \(W\), so does each extracted component. This proves
the direct sum assertion. Applying it to representatives in \(V/W\)
proves the quotient assertion.
\end{proof}

\begin{theorem}[The unique simple highest-weight quotient]
\label{res:G03}
For every \(\lambda\in\mathfrak h^*\), \(M(\lambda)\) has a unique
maximal proper submodule and hence a unique simple quotient
\(L(\lambda)\). Every nonzero highest-weight module of weight
\(\lambda\) surjects onto \(L(\lambda)\).
Every simple finite-dimensional \(\mathfrak g\)-module is one of
these quotients; its highest weight determines its isomorphism class.
\end{theorem}
\begin{proof}
A proper submodule of \(M(\lambda)\) has zero intersection with
its top weight line: containing \(v_\lambda\) would make it the
whole module. Lemma~\ref{supp:weight-submodules} therefore implies
that the sum \(J(\lambda)\) of all proper submodules still has zero
top weight space. It is proper, contains every proper submodule,
and is the unique maximal one. If a nonzero highest-weight quotient
has kernel \(K\), then \(K\subseteq J(\lambda)\), giving the asserted
surjection. For a simple quotient this surjection is an isomorphism.

For a finite-dimensional module, Theorem~\ref{res:C07} supplies
the Cartan weight decomposition. Choose a weight maximal for the
order defined by \(R^+\); it exists because there are finitely many
weights. Every positive-root operator kills that weight space.
A nonzero vector in it generates the module when the module is
simple. The universal property applies.
Finally if two highest-weight generators have weights \(\lambda,\mu\),
all weights of the generated module lie both in \(\lambda-Q_+\)
and in \(\mu-Q_+\). Each of \(\lambda,\mu\) is a weight, so
\(\lambda\leq\mu\leq\lambda\), giving \(\lambda=\mu\).
\end{proof}

\section{The finite-dimensionality criterion}

\begin{proposition}[Necessity of dominance and integrality]
\label{res:G04}
If \(L(\lambda)\) is finite dimensional, then \(\lambda\in P_+\).
More generally, every finite-dimensional highest-weight module of
weight \(\lambda\) satisfies this conclusion and
\(f_i^{\lambda(h_i)+1}v_\lambda=0\).
\end{proposition}
\begin{proof}
Restrict to the \(i\)-th \(\mathfrak{sl}_2\).
Its highest vector \(v_\lambda\) has \(h_i\)-weight \(\lambda(h_i)\).
The rank-one classification, Theorem~\ref{res:C01}, forces this
weight to be a nonnegative integer and gives the asserted last
nonzero step of the \(f_i\)-string. This holds for every \(i\).
\end{proof}

\begin{lemma}[Local finiteness from the generator]
\label{supp:local-integrability}
For \(\lambda\in P_+\), the quotient
\[
 Q(\lambda)=M(\lambda)/
     \sum_i U(\mathfrak g)f_i^{\lambda(h_i)+1}v_\lambda
\]
is nonzero. Both \(e_i\) and \(f_i\) act locally nilpotently on it.
Every vector lies in a finite-dimensional module for the
\(i\)-th \(\mathfrak{sl}_2\).
\end{lemma}
\begin{proof}
Write \(m_i=\lambda(h_i)\).
The rank-one commutator gives
\(e_i f_i^{m_i+1}v_\lambda=0\).
For \(j\ne i\), \([e_j,f_i]=0\), so
\(e_j f_i^{m_i+1}v_\lambda=0\) as well.
These vectors are thus highest vectors of weights strictly below
\(\lambda\). Each of their generated submodules has only weights
below its own highest weight. Their sum has no \(\lambda\) weight,
so \(v_\lambda\) survives in \(Q(\lambda)\).

The adjoint operators of \(e_i,f_i\) are nilpotent on
\(\mathfrak g\), by the root strings. Their derivation extensions
are locally nilpotent on \(U(\mathfrak g)\): on a product of finitely
many generators the repeated derivative eventually differentiates
one factor more often than its nilpotence bound. The identity
\[
 x^n(uv)=\sum_{j=0}^n\binom nj(\operatorname{ad}x)^j(u)x^{n-j}v
\]
therefore proves local nilpotence on \(Q(\lambda)\), starting from
\(e_iv_\lambda=0\) and \(f_i^{m_i+1}v_\lambda=0\).

For a weight vector \(v\), PBW for \(\mathfrak{sl}_2\) spans its
generated module by \(f_i^a h_i^b e_i^c v\).
Only finitely many \(c\) occur; each \(e_i^cv\) is a weight vector,
so \(h_i^b\) contributes only scalar factors; and for each of these
finitely many vectors only finitely many \(a\) occur. Thus the span
is finite. A general vector is a finite sum of weight vectors,
which proves the last assertion.
\end{proof}

\begin{theorem}[Existence of finite-dimensional highest-weight modules]
\label{res:G05}
For every \(\lambda\in P_+\), \(Q(\lambda)\) is finite dimensional.
Its weight multiplicities are invariant under the Weyl group.
Consequently \(L(\lambda)\) is finite dimensional.
\end{theorem}
\begin{proof}
Consider a weight module locally finite for the \(i\)-th
\(\mathfrak{sl}_2\). If \(\mu(h_i)=m\geq0\), then
\(f_i^m:V_\mu\to V_{s_i\mu}\) is injective: any nonzero vector
generates a finite-dimensional rank-one module, and the explicit
rank-one decomposition shows that this operator is injective on
its \(m\) weight space. If \(m\leq0\), use \(e_i^{-m}\) instead.
Applying the same argument to \(s_i\mu\) gives equality of
dimensions. Weights occurring here are integral because restriction
to every rank-one algebra is a sum of its finite-dimensional modules.
Since the \(s_i\) generate \(W\), all multiplicities are
\(W\)-invariant.

We prove finiteness explicitly. If a weight \(\mu\) is dominant,
write \(\mu=\sum_i a_i\omega_i\) with \(a_i\in\mathbb Z_{\geq0}\).
Let \(\rho^\vee\) be half the sum of the positive coroots. Each
\(\omega_i(\rho^\vee)>0\): the coefficient of \(\alpha_i^\vee\)
in this positive sum is positive. Since
\(\alpha_i(\rho^\vee)=1\) and \(\lambda-\mu\in Q_+\), we have
\[
 0\leq\sum_i a_i\omega_i(\rho^\vee)
      =\mu(\rho^\vee)\leq\lambda(\rho^\vee).
\]
This bounds every nonnegative integer \(a_i\); thus there are only
finitely many dominant weights below \(\lambda\).
Every integral weight has a Weyl conjugate in the dominant chamber,
and \(W\) is finite. The weight set of \(Q(\lambda)\) is therefore
finite. Its weight spaces are finite as quotients of those of
\(M(\lambda)\). This proves finite dimensionality, without assuming
it in the local-finiteness argument.
The map \(M(\lambda)\to L(\lambda)\) kills the indicated singular
vectors (or use Proposition~\ref{res:G03}); hence \(L(\lambda)\)
is a quotient of \(Q(\lambda)\) and is finite dimensional.
\end{proof}

\begin{theorem}[Classification and defining relations]
\label{res:G06}
The map \(\lambda\mapsto L(\lambda)\) is a bijection from \(P_+\)
to the simple finite-dimensional \(\mathfrak g\)-modules up to
isomorphism. In addition, \(Q(\lambda)=L(\lambda)\). Thus
\(L(\lambda)\) is generated by one vector \(v\) with precisely the
defining relations
\[
 e_iv=0,\quad hv=\lambda(h)v,\quad
 f_i^{\lambda(h_i)+1}v=0 .
\]
Every finite-dimensional module is a direct sum of these modules.
\end{theorem}
\begin{proof}
Necessity, existence and uniqueness follow from
Propositions~\ref{res:G03}--\ref{res:G04} and
Theorem~\ref{res:G05}. To verify the stronger presentation assertion,
apply Weyl complete reducibility (Theorem~\ref{res:C05}) to
\(Q(\lambda)\). Project its cyclic highest vector onto each simple
summand. A nonzero projection is again a highest vector of weight
\(\lambda\), so that summand is \(L(\lambda)\); a zero projection
would contradict cyclic generation of the whole summand.
The \(\lambda\) weight space has dimension one. Exactly one summand
can therefore occur. The presentation follows.
The last assertion is Weyl complete reducibility with the just-proved
classification of its simple summands.
\end{proof}

\section{Strings, duals and examples}

\begin{proposition}[Weight strings and the convex bound]
\label{res:G07}
For a finite-dimensional module, multiplicities are \(W\)-invariant.
For every root \(\alpha\), the subspace
\(\bigoplus_{j\in\mathbb Z}V_{\mu+j\alpha}\) decomposes into rank-one
strings, each symmetric under \(s_\alpha\).
In particular, if \(m=\langle\mu,\alpha^\vee\rangle\geq0\)
and \(V_\mu\ne0\), all \(\mu-j\alpha\), \(0\leq j\leq m\),
are weights. Every weight of \(L(\lambda)\) belongs to
\[
 (\lambda+Q)\cap\operatorname{conv}(W\lambda).
\]
\end{proposition}
\begin{proof}
The direct sum over the \(\alpha\)-coset is stable under the
root triple. Decompose it into simple rank-one modules using
Theorem~\ref{res:C06}; the Cartan directions killed by \(\alpha\)
act by fixed scalars on this sum and do not affect that
decomposition. In each rank-one module the weights are an unbroken
string, with the explicit raising and lowering coefficients of
Theorem~\ref{res:C01}. Reflection exchanges its two ends, with
equal multiplicities. The claims about strings and multiplicities
follow by summing these decompositions.

All weights are in \(\lambda-Q_+\), so are in \(\lambda+Q\).
For the convex bound let \(H\) lie in the real Cartan space.
Move \(H\) by \(W\) to the closed dominant chamber. For any weight
\(\nu\), \((\lambda-\nu)(H)\geq0\) in that chamber.
The \(W\)-invariance already proved therefore implies, for arbitrary
\(H\),
\(\nu(H)\leq\max_{w\in W}(w\lambda)(H)\).
The separating-hyperplane characterization of the convex hull of
a finite set gives the required inclusion. Equivalently that
characterization follows by minimizing distance to the compact
convex polytope: a point outside it is strictly separated by the
vector from its nearest point.
\end{proof}

\begin{proposition}[Dual highest weights and Cartan products]
\label{res:G08}
The dual of \(L(\lambda)\) is \(L(-w_0\lambda)\).
The tensor product \(L(\lambda)\otimes L(\mu)\) contains one copy
of \(L(\lambda+\mu)\); its highest line is generated by
\(v_\lambda\otimes v_\mu\).
More generally the tensor product of the fundamental modules,
with \(\omega_i\) repeated \(\lambda(h_i)\) times, contains
\(L(\lambda)\) as the submodule generated by their top vectors.
\end{proposition}
\begin{proof}
The dual action is \((x\varphi)(v)=-\varphi(xv)\), so the
dual weights are negatives of the original weights with the same
multiplicities. The lowest weight of \(L(\lambda)\) is \(w_0\lambda\):
apply \(w_0\) to the inequalities \(\nu\leq\lambda\), using that
it changes positive roots to negative roots.
That weight has multiplicity one by Weyl invariance.
Consequently the dual has top weight \(-w_0\lambda\);
it is simple since a proper submodule would annihilate a nonzero
proper submodule quotient of \(L(\lambda)\).

In the tensor product all weights are below \(\lambda+\mu\),
and its top space is exactly the tensor product of the two top
lines. That vector is killed by every \(e_i\).
The generated submodule is a finite-dimensional cyclic highest-weight
module, so Theorem~\ref{res:G06} and its complete-reducibility
argument show that it is \(L(\lambda+\mu)\).
No second copy can occur because the top line is one dimensional.
The assertion about fundamental modules follows by iteration;
for \(\lambda=0\) use the empty tensor product \(\mathbb C\).
\end{proof}

\begin{example}[Type \(A\) and the polynomial comparison]
\label{res:G09}
For \(\mathfrak{sl}_2\), \(L(n\omega_1)=\operatorname{Sym}^n\mathbb C^2\).
For \(\mathfrak{sl}_d\), the defining module has highest weight
\(\omega_1\), \(\bigwedge^k\mathbb C^d\) has highest weight
\(\omega_k\) for \(1\leq k<d\), and
\(\operatorname{Sym}^a\mathbb C^d=L(a\omega_1)\).
For \(\mathfrak{sl}_3\),
\[
 \mathbb C^3\otimes(\mathbb C^3)^*
   =\mathbb C\oplus L(\omega_1+\omega_2),
 \qquad \dim L(\omega_1+\omega_2)=8 .
\]
A polynomial \(\operatorname{GL}_d\)-module indexed by a partition
\(\tau=(\tau_1\geq\cdots\geq\tau_d\geq0)\) restricts to the
\(\mathfrak{sl}_d\) highest weight
\(\sum_{i<d}(\tau_i-\tau_{i+1})\omega_i\).
Two such parameters have the same restricted highest weight exactly
when they differ by an integral power of the determinant.
\end{example}
\begin{proof}
In the defining matrix action \(e_i=E_{i,i+1}\) and
\(h_i=E_{ii}-E_{i+1,i+1}\). On polynomials,
\(e_i=x_i\partial/\partial x_{i+1}\).
A homogeneous polynomial killed by all \(e_i\) is a scalar
multiple of \(x_1^a\): each derivative with respect to
\(x_2,\ldots,x_d\) must vanish in characteristic zero.
Thus a completely reducible symmetric power has a unique highest
line and is simple with the asserted weight.

The weights in the exterior power are the distinct sums of \(k\)
coordinate weights. A basis wedge is killed by all \(e_i\) exactly
when its indices are \(1,\ldots,k\). There cannot be a cancellation
of their images between distinct weights. Therefore it also has
a unique highest line and is simple; evaluating \(h_i\) gives
\(\omega_k\).
The tensor product with the dual is the conjugation representation
on \(\operatorname{End}(\mathbb C^3)\).
Trace splits it into scalar and trace-zero matrices.
The latter is the adjoint module of the simple algebra
\(\mathfrak{sl}_3\); its highest root is
\(\alpha_1+\alpha_2=\omega_1+\omega_2\), and its dimension is eight.

For the final comparison, the polynomial module construction and
its simple highest line are those of
\cite[Chapter 4]{finite-representations}. Differentiation evaluates
the diagonal highest character \(t_1^{\tau_1}\cdots t_d^{\tau_d}\)
on \(h_i\) as \(\tau_i-\tau_{i+1}\).
A Lie-stable subspace is group-stable here: the elementary
unipotent matrices act by exponentials of \(E_{ij}\), and diagonal
matrices act on the same integral weight spaces.
These matrices generate \(\operatorname{GL}_d(\mathbb C)\).
Thus a simple polynomial group module remains simple on
restriction to the Lie algebra after removing its scalar central
action. Finally equality of all successive differences is
equivalent to \(\tau_i-\sigma_i\) being a constant integer for
all \(i\). Multiplication by the corresponding determinant
character has exactly this effect.
\end{proof}

\chapter{Characters, Multiplicities, and the Center}
\label{ch:characters}

Keep the complex semisimple algebra and positive system of
Chapter~\ref{ch:highest-weights}. We use a nondegenerate invariant
form \(B\), positive on the real root space after identifying
\(\mathfrak h\) with its dual. In particular every squared norm
in a dominance argument is taken on that real space.
Let
\[
 \rho=\tfrac12\sum_{\alpha>0}\alpha,\quad
 \varepsilon(w)=\det(w),\quad
 A_\nu=\sum_{w\in W}\varepsilon(w)e^{w\nu},\quad
 D=e^\rho\prod_{\alpha>0}(1-e^{-\alpha}).
\]
The symbols \(e^\mu\) are formal basis elements, with
\(e^\mu e^\nu=e^{\mu+\nu}\), rather than a choice of a Lie-group
exponential.

\section{Formal characters and a finite truncation argument}

\begin{proposition}[Characters and their completion]
\label{res:H01}
For a finite-dimensional module set
\(\operatorname{ch}V=\sum_\mu(\dim V_\mu)e^\mu\in\mathbb Z[P]\).
Characters add in exact sequences, multiply in tensor products,
and are \(W\)-invariant. The characters of the \(L(\lambda)\),
\(\lambda\in P_+\), are linearly independent.
There is also a well-defined ring of series over the additive group
\(\mathfrak h^*\) whose support lies in a finite union of sets
\(\lambda-Q_+\). Integral-weight series form a subring; in the
larger ring
\[
 \operatorname{ch}M(\lambda)
   =e^\lambda\prod_{\alpha>0}(1-e^{-\alpha})^{-1}.
\]
\end{proposition}
\begin{proof}
Exactness holds weight by weight by
Lemma~\ref{supp:weight-submodules}, and
\((V\otimes U)_\eta=\bigoplus_{\mu+\nu=\eta}V_\mu\otimes U_\nu\).
These facts prove additivity and multiplication.
Weyl invariance is Proposition~\ref{res:G07}.
In a finite linear relation between simple characters, choose a
maximal highest weight among the terms. Its coefficient receives
no contribution from a different term, since all other weights
lie below that term's own highest weight. Its multiplicity in its
own character is one, so its coefficient is zero. Repeating proves
independence.

For the series ring, a coefficient of a product can receive only
finitely many contributions. Indeed, for fixed upper bounds
\(\lambda,\mu\), the decompositions
\(\beta+\gamma=\lambda+\mu-\eta\), \(\beta,\gamma\in Q_+\),
are finite: the simple-root coefficients are nonnegative integers
with bounded sum. A finite union of upper bounds makes no
difference to finiteness. The same observation proves associativity
by rearranging finite coefficient sums.
Finally PBW in \(U(\mathfrak n^-)\) makes a negative-root exponent
independent for every positive root. Summing each geometric series
gives the displayed formula. No Weyl action on this entire
completion is assumed; we use that action on finite expressions
when taking alternating sums.
\end{proof}

\begin{proposition}[The denominator, the Casimir, and truncated expansion]
\label{res:H02}
The denominator identity is \(D=A_\rho\).
The Casimir for \(B\) acts on any highest-weight module of weight
\(\lambda\) by
\[
 c_\lambda=(\lambda,\lambda+2\rho)
           =\|\lambda+\rho\|^2-\|\rho\|^2 .
\]
If a weight module \(V\), supported in \(\lambda-Q_+\), has finite
weight spaces and Casimir scalar \(c\), then, to any fixed height
bound, its character is an integral linear combination of
Verma characters \(M(\mu)\) with \(c_\mu=c\).
\end{proposition}
\begin{proof}
First \(s_i\) permutes the positive roots other than \(\alpha_i\),
and \(s_i\rho=\rho-\alpha_i\). Substitution in the product for \(D\)
gives \(s_iD=-D\).
In a finite alternating expression a coefficient on a reflection
hyperplane is zero. Every remaining orbit has a unique strictly
dominant representative \(\nu\), and its contribution is a
multiple of \(A_\nu\).
Every exponent of \(D\) is in \(\rho-Q_+\).
If such a \(\nu\) is strictly dominant and integral, then
\(\nu-\rho\) is dominant. If it were nonzero, its pairing with
\(\rho\) would be positive: \(\rho\) is a positive sum of simple
roots containing every simple root, and a nonzero dominant vector
pairs nonnegatively with all of them and positively with at least
one. On the other hand \(\nu-\rho\in-Q_+\), which pairs
nonpositively with \(\rho\), since
\((\rho,\alpha_i)>0\). This is impossible.
Thus only \(A_\rho\) occurs; its top coefficient is one.
This proves the denominator identity independently of the character
formula.

Choose dual bases \(h_a,h^a\) of \(\mathfrak h\), and root vectors
\(x_\alpha,y_\alpha\) with \(B(x_\alpha,y_\alpha)=1\).
Invariance gives
\([x_\alpha,y_\alpha]=t_\alpha\), where
\(B(t_\alpha,h)=\alpha(h)\).
The Casimir is
\[
 \Omega=\sum_a h_ah^a+
       \sum_{\alpha>0}(x_\alpha y_\alpha+y_\alpha x_\alpha)
 =\sum_a h_ah^a+2t_\rho+
       2\sum_{\alpha>0}y_\alpha x_\alpha.
\]
The dual tensor is invariant, so its multiplication commutes with
every \(x\in\mathfrak g\); this is also the calculation in
Theorem~\ref{res:C02}. Acting on a highest vector gives
\(\|\lambda\|^2+2(\rho,\lambda)\).
Since the module is generated by that vector and \(\Omega\) is
central, this is its action on the whole module.

For the expansion, the height of a weight \(\lambda-\beta\)
is \(\operatorname{ht}\beta\). If \(V\ne0\), take all weight
spaces of minimal height, and choose a basis of each.
There are finitely many of these spaces, each finite dimensional.
Every selected vector is highest, because raising a weight
would give a smaller height. The universal property gives a map
\[
 T=\bigoplus_j M(\mu_j)\longrightarrow V
\]
which is an isomorphism on the entire minimal-height layer.
The Casimir scalar on each summand is \(c_{\mu_j}=c\), as may
be checked on its nonzero image of the highest vector.
Both its kernel \(K\) and cokernel \(C\) therefore have Casimir
scalar \(c\) and have zero spaces on that layer. Exactness gives
\[
 \operatorname{ch}V=\operatorname{ch}T-\operatorname{ch}K+
                       \operatorname{ch}C.
\]
Repeat separately for \(K\) and \(C\). Their minimal heights have
strictly increased. To compute all coefficients through height
\(N\), stop after at most \(N+1\) such layers on each branch.
There are finitely many branches, and finitely many Verma modules
at each node. This proves the truncated assertion.
Taking the whole minimal layer is essential: removing only one
highest vector could leave the same highest weight in the cokernel.
The proof does not claim that an arbitrary module has a finite
untruncated Verma resolution.
\end{proof}

\begin{theorem}[Weyl character formula]
\label{res:H03}
For \(\lambda\in P_+\),
\[
 \operatorname{ch}L(\lambda)=
 \frac{A_{\lambda+\rho}}{A_\rho}
 =\frac{\sum_w\varepsilon(w)e^{w(\lambda+\rho)}}
        {e^\rho\prod_{\alpha>0}(1-e^{-\alpha})}.
\]
The fraction denotes the equality after multiplication by the
denominator, or its expansion in the completion of
Proposition~\ref{res:H01}.
\end{theorem}
\begin{proof}
The finite expression \(D\operatorname{ch}L(\lambda)\) is
alternating. Its exponents lie in \(\lambda+\rho-Q_+\), and the
coefficient of \(e^{\lambda+\rho}\) is one.
Apply the truncated expansion in Proposition~\ref{res:H02} at
arbitrarily large height bounds. Multiplication by
\(\prod_{\alpha>0}(1-e^{-\alpha})\) only lowers weights,
so any fixed coefficient is computed by a finite truncation.
Whenever \(e^{\mu+\rho}\) has a nonzero coefficient, some
Verma term contributing to it has \(c_\mu=c_\lambda\). Thus
\[
 \|\mu+\rho\|^2=\|\lambda+\rho\|^2.
\]
Choose the strictly dominant representative \(\nu\) of any orbit
in this alternating expression. It too is an exponent in the
support, so \(\delta=\lambda+\rho-\nu\in Q_+\).
If \(\delta\ne0\), strict dominance of both \(\lambda+\rho\)
and \(\nu\) gives
\[
 \|\lambda+\rho\|^2-\|\nu\|^2
     =(\delta,\lambda+\rho+\nu)>0,
\]
contrary to the Casimir equality. Therefore the sole possible
orbit is that of \(\lambda+\rho\).
Its top coefficient determines it to be \(A_{\lambda+\rho}\).
Singular orbits have zero coefficients by alternation.
This proves the formula without using Harish--Chandra's theorem.
\end{proof}

\begin{corollary}[Weyl dimension formula]
\label{res:H04}
For every dominant integral weight,
\[
 \dim L(\lambda)
  =\prod_{\alpha>0}\frac{(\lambda+\rho,\alpha)}{(\rho,\alpha)}
  =\prod_{\alpha>0}
       \frac{\langle\lambda+\rho,\alpha^\vee\rangle}
            {\langle\rho,\alpha^\vee\rangle}.
\]
The formula includes \(\lambda=0\), products of simple factors,
and the empty root system.
\end{corollary}
\begin{proof}
Evaluate the character at \(e^\mu=\exp(2t(\mu,\rho))\).
Invariance of the inner product and substitution \(w\mapsto w^{-1}\)
give
\[
 A_{\lambda+\rho}(2t\rho)
  =\sum_w\varepsilon(w)\exp(2t(w\rho,\lambda+\rho))
  =A_\rho(2t(\lambda+\rho)).
\]
By the denominator identity this is
\(\prod_{\alpha>0}
 (\exp(t(\alpha,\lambda+\rho))-\exp(-t(\alpha,\lambda+\rho)))\).
The denominator has the same expression with \(\lambda=0\).
Divide, and let \(t\to0\), using
\((e^{ta}-e^{-ta})/(e^{tb}-e^{-tb})\to a/b\) for \(b\ne0\).
Every denominator is positive by the choice of positive system.
The left side tends to the sum of the weight multiplicities.
Replacing roots by coroots cancels the same nonzero factor in
each quotient. All factors are independent on different simple
components, so the asserted conventions follow as well.
\end{proof}

\section{Weight multiplicities and tensor coefficients}

\begin{theorem}[Freudenthal's recursion]
\label{res:H05}
Let \(m_\lambda(\mu)=\dim L(\lambda)_\mu\), extended by zero
away from the weight set. Then
\[
 \bigl(\|\lambda+\rho\|^2-\|\mu+\rho\|^2\bigr)m_\lambda(\mu)
 =2\sum_{\alpha>0}\sum_{j\geq1}
        (\mu+j\alpha,\alpha)m_\lambda(\mu+j\alpha).
\]
Together with \(m_\lambda(\lambda)=1\), this is a terminating
recursion on the finite set
\((\lambda+Q)\cap\operatorname{conv}(W\lambda)\).
\end{theorem}
\begin{proof}
Set \(T_\mu^\alpha=\operatorname{tr}
(y_\alpha x_\alpha|_{L(\lambda)_\mu})\).
For maps between two different finite-dimensional spaces,
\(\operatorname{tr}(AB)=\operatorname{tr}(BA)\);
this follows by writing their two rectangular matrices.
Hence, using \([x_\alpha,y_\alpha]=t_\alpha\),
\[
 T_\mu^\alpha
   =\operatorname{tr}(x_\alpha y_\alpha|_{L(\lambda)_{\mu+\alpha}})
   =T_{\mu+\alpha}^\alpha+
                 (\mu+\alpha,\alpha)m_\lambda(\mu+\alpha).
\]
Iteration terminates because there are finitely many weights,
and gives
\(T_\mu^\alpha=\sum_{j\geq1}
(\mu+j\alpha,\alpha)m_\lambda(\mu+j\alpha)\).
Now take traces of the normal-ordered Casimir in
Proposition~\ref{res:H02} on the \(\mu\) weight space.
Its Cartan and linear terms have trace
\((\|\mu\|^2+2(\rho,\mu))m_\lambda(\mu)\).
Subtracting them from \(c_\lambda m_\lambda(\mu)\)
proves the identity.

By Proposition~\ref{res:G07} the weight set lies in the indicated
bounded polytope. For \(\mu\) in that polytope,
\(\|\mu\|^2\leq\|\lambda\|^2\), since a convex combination of
vectors of norm \(\|\lambda\|\) has no larger norm.
If also \(\mu\ne\lambda\) and \(\lambda-\mu\in Q_+\), then
\((\rho,\lambda-\mu)>0\). Consequently
\[
 \|\lambda+\rho\|^2-\|\mu+\rho\|^2
 =\|\lambda\|^2-\|\mu\|^2+2(\rho,\lambda-\mu)>0.
\]
All terms on the right have smaller height below \(\lambda\).
The lattice intersection with the polytope is finite. These
observations prove both that division is legitimate away from
the top weight and that the recursion terminates. Candidates
outside the cone or polytope have multiplicity zero.
\end{proof}

\begin{theorem}[Kostant's multiplicity formula]
\label{res:H06}
Let \(p(\beta)\) count the unordered decompositions of \(\beta\)
as a sum of positive roots, with \(p(0)=1\), and set \(p(\beta)=0\)
when \(\beta\notin Q_+\). Then
\[
 m_\lambda(\mu)=
   \sum_{w\in W}\varepsilon(w)
      p\bigl(w(\lambda+\rho)-(\mu+\rho)\bigr).
\]
\end{theorem}
\begin{proof}
Each root may occur with any nonnegative multiplicity, so in the
completed ring
\[
 \prod_{\alpha>0}(1-e^{-\alpha})^{-1}
     =\sum_{\beta\in Q_+}p(\beta)e^{-\beta}.
\]
Every coefficient is a finite sum by the height bound used in
Proposition~\ref{res:G02}.
Insert this expression into Theorem~\ref{res:H03}, cancel
\(e^\rho\), and compare the coefficient of \(e^\mu\).
The equality of exponents requires precisely
\(\beta=w(\lambda+\rho)-(\mu+\rho)\).
The zero convention covers weights outside the root-lattice
coset and outside the positive cone.
\end{proof}

\section{Polynomial invariants and central characters}

Polynomial functions on a vector space \(E\) form \(S(E^*)\).
The form \(B\) identifies \(S(\mathfrak g)\) with polynomial
functions on \(\mathfrak g\), compatibly with the adjoint action;
similarly on \(\mathfrak h\). When speaking of restriction we
always use these identifications. The group of inner automorphisms
is generated by \(\exp(t\operatorname{ad}x)\), \(x\in\mathfrak g\).
For root vectors these exponentials are finite on \(\mathfrak g\);
their rank-one products represent the Weyl reflections.

\begin{theorem}[Chevalley restriction]
\label{res:H07}
Restriction is an isomorphism
\[
 S(\mathfrak g^*)^{\mathfrak g}\longrightarrow
 S(\mathfrak h^*)^W .
\]
The source is generated as an algebra by trace polynomials
\(x\mapsto\operatorname{tr}_{L(\lambda)}(x^m)\),
\(\lambda\in P_+\), \(m\geq1\), and constants.
\end{theorem}
\begin{proof}
Infinitesimal invariance is equivalent to invariance under all
inner exponentials. If every infinitesimal action kills a
polynomial, its value along \(\exp(t\operatorname{ad}x)y\) has
zero derivative for every \(t\), and hence is constant.
Conversely differentiate such invariance at \(t=0\).
The root directions suffice to test infinitesimal invariance,
because their operators generate the Lie algebra action.
The representatives of the simple Weyl reflections are products
of these exponentials, so restriction lands in the stated
invariants. If restriction is zero, invariance makes the polynomial
vanish on all conjugates of \(\mathfrak h\).
They include the dense regular semisimple set from
Proposition~\ref{res:D06}; a polynomial vanishing on a dense set
is zero. This proves injectivity.

Here is a direct surjectivity proof, including the invariant
polynomials which must be constructed. The lattice \(P\) spans
\(\mathfrak h^*\), and pure powers of lattice elements span its
symmetric powers. For example polarization expresses a product
\(\eta_1\cdots\eta_m\) as
\[
 \frac1{m!}\sum_{J\subseteq\{1,\ldots,m\}}
       (-1)^{m-|J|}\left(\sum_{j\in J}\eta_j\right)^m.
\]
Averaging over \(W\), and then moving a lattice element to the
dominant chamber, shows that the invariant polynomials are spanned
by orbit sums
\(q_{\lambda,m}=\sum_{\nu\in W\lambda}\nu^m\),
\(\lambda\in P_+\). Constants can be treated separately.

For such a \(\lambda\), take \(L(\lambda)\), whose existence has
already been proved. The trace polynomial
\(T_{\lambda,m}(x)=\operatorname{tr}(\pi_\lambda(x)^m)\)
is infinitesimally invariant: its derivative in the direction
\([y,x]\) is
\(m\operatorname{tr}(\pi_\lambda(x)^{m-1}
[\pi_\lambda(y),\pi_\lambda(x)])=0\) by trace cyclicity.
The preceding paragraph gives its inner invariance as well.
Its restriction is
\[
 T_{\lambda,m}|_{\mathfrak h}
   =q_{\lambda,m}
     +\sum_{\substack{\nu\in P_+\\\nu<\lambda}}
                  m_\lambda(\nu)q_{\nu,m}.
\]
Only weights of \(L(\lambda)\) occur. There are finitely many
dominant weights below any fixed \(\lambda\), by the
\(\rho^\vee\) bound in Theorem~\ref{res:G05}.
Induction on that finite lower set therefore lifts every
\(q_{\lambda,m}\) as a linear combination of trace polynomials.
This proves surjectivity and the generation assertion.
No conclusion about surjectivity has been inferred just from
root-direction differentiation.
\end{proof}

\begin{theorem}[Harish--Chandra isomorphism and central characters]
\label{res:H08}
Let \(Z=Z(U(\mathfrak g))\). In the triangular PBW ordering,
project onto the monomials entirely in \(\mathfrak h\):
\(p:U(\mathfrak g)\to U(\mathfrak h)=S(\mathfrak h)\).
For \(z\in Z\), define
\[
 \gamma(z)(\eta)=p(z)(\eta-\rho),\qquad \eta\in\mathfrak h^*.
\]
Then \(\gamma:Z\longrightarrow S(\mathfrak h)^W\) is an
algebra isomorphism. The central character on \(M(\lambda)\)
and its nonzero highest-weight quotients is
\(\chi_\lambda(z)=\gamma(z)(\lambda+\rho)\), and
\[
 \chi_\lambda=\chi_\mu
    \quad\Longleftrightarrow\quad
 \mu+\rho\in W(\lambda+\rho).
\]
Every complex algebra character of \(Z\) is obtained in this way.
\end{theorem}
\begin{proof}
A central \(z\) has Cartan weight zero. In its PBW expansion,
a term with a positive-root factor on the right kills a highest
vector. A term with negative-root factors and none positive
cannot have weight zero unless those negative factors are absent.
Thus \(zv_\lambda=p(z)(\lambda)v_\lambda\).
Centrality extends this scalar action to the generated module.
It follows for all \(\lambda\) that
\(p(zz')(\lambda)=p(z)(\lambda)p(z')(\lambda)\).
Two polynomials equal at all \(\lambda\) are equal, proving
multiplicativity on \(Z\).

We check the shifted invariance without presuming the theorem.
If \(\lambda(h_i)=m\in\mathbb Z_{\geq0}\), the vector
\(f_i^{m+1}v_\lambda\) in \(M(\lambda)\) is nonzero by PBW.
The rank-one relation makes it killed by \(e_i\), and it is
killed by \(e_j\), \(j\ne i\), because \([e_j,f_i]=0\).
It has weight \(s_i(\lambda+\rho)-\rho\).
There is therefore a nonzero highest-weight map from the Verma
module of that weight into \(M(\lambda)\).
The center has the same scalar on the two modules, so
\[
 p(z)(\lambda)
    =p(z)\bigl(s_i(\lambda+\rho)-\rho\bigr).
\]
For each nonnegative integer \(m\) this is a polynomial identity
on the entire affine hyperplane \(\lambda(h_i)=m\);
the other coordinates are arbitrary. A polynomial in
\(\lambda(h_i)\), with coefficients in the remaining coordinates,
cannot have infinitely many distinct roots unless it is zero.
The identity thus holds everywhere. After translation by
\(\rho\) it says that \(\gamma(z)\) is \(s_i\)-invariant,
and hence \(W\)-invariant.

Equip \(U(\mathfrak g)\) with PBW degree.
The leading symbol of a central element is in \(S(\mathfrak g)^{\mathfrak g}\).
The leading symbol of \(p(z)\), if not zero, is its restriction
to \(\mathfrak h\). By Chevalley restriction that restriction
cannot be zero for a nonzero invariant symbol. Translation
by \(-\rho\) does not change the leading symbol. It follows that
\(\gamma(z)\ne0\) for \(z\ne0\), proving injectivity.

For surjectivity, take a homogeneous invariant polynomial
\(q\in S(\mathfrak h)^W\) of degree \(d\).
By Theorem~\ref{res:H07} and the form \(B\), lift it to a
homogeneous \(F\in S(\mathfrak g)^{\mathfrak g}\).
The symmetrization map
\[
 x_1\cdots x_d\longmapsto
       \frac1{d!}\sum_{\sigma\in S_d}x_{\sigma(1)}\cdots x_{\sigma(d)}
\]
is a filtered \(\mathfrak g\)-module isomorphism to
\(U(\mathfrak g)\) by PBW.
It carries \(F\) to a central element \(z\) whose leading
symbol restricts to \(q\). Thus \(\gamma(z)-q\) is an invariant
of smaller degree. Induction on degree corrects this difference;
degree zero is immediate. This proves the isomorphism.

Finally finite-group invariants separate distinct orbits.
On two disjoint finite orbits choose, by polynomial interpolation,
a polynomial equal to one on the first and zero on the second.
One construction makes a delta polynomial at each point by
multiplying linear functions separating it from all other points,
then sums those delta polynomials. Averaging over \(W\) keeps
the two prescribed constant values and makes it invariant.
This proves the equivalence of central characters.
To see that no characters are missing, the polynomial ring
\(A=S(\mathfrak h)\) is integral and finite over \(A^W\):
each of its finitely many linear generators satisfies the monic
polynomial \(\prod_{w\in W}(t-w\ell)\), and adjoining finitely
many integral generators is finite.
Lying-over lifts a maximal ideal of \(A^W\) to one of \(A\).
The weak Nullstellensatz identifies that ideal of \(A\) with
evaluation at some \(\eta\in\mathfrak h^*\).
These elementary integral-extension facts are among the
prerequisites in \cite[Chapter 4]{p103}.
Take \(\lambda=\eta-\rho\).
\end{proof}

\begin{theorem}[Steinberg's tensor multiplicities and small-rank checks]
\label{res:H09}
If \(\lambda,\mu,\nu\in P_+\), the multiplicity of \(L(\nu)\)
in \(L(\lambda)\otimes L(\mu)\) is
\[
 c_{\lambda\mu}^{\nu}
 =\sum_{u,w\in W}\varepsilon(u)\varepsilon(w)
   p\bigl(u(\lambda+\rho)+w(\mu+\rho)-(\nu+2\rho)\bigr).
\]
For type \(A_1\) this gives
\(L(a)\otimes L(b)=\bigoplus_{j=0}^{\min(a,b)}L(a+b-2j)\).
For \(A_2\),
\[
 \dim L(a\omega_1+b\omega_2)
    =\frac{(a+1)(b+1)(a+b+2)}2 ,
 \qquad p(s\alpha_1+t\alpha_2)=
      \min(s,t)+1\quad(s,t\geq0).
\]
\end{theorem}
\begin{proof}
By complete reducibility and the character formula,
\[
 \operatorname{ch}L(\lambda) A_{\mu+\rho}
       =\sum_{\xi\in P_+}c_{\lambda\mu}^{\xi}A_{\xi+\rho}.
\]
For dominant \(\nu\), the coefficient of \(e^{\nu+\rho}\)
on the right is just \(c_{\lambda\mu}^\nu\):
two strictly dominant weights in the same Weyl orbit agree.
The coefficient on the left is
\(\sum_w\varepsilon(w)m_\lambda(\nu+\rho-w(\mu+\rho))\).
Insert Kostant's formula and simplify its argument to obtain
the stated double sum.

For \(A_1\), the rank-one character is
\(e^a+e^{a-2}+\cdots+e^{-a}\).
The coefficient counts in the product, followed by successive
subtraction of its top rank-one character, give the indicated
list, each once. This also checks the double sum since
\(p(j\alpha)=1\) for \(j\geq0\) and zero otherwise.
For \(A_2\) the positive roots are
\(\alpha_1,\alpha_2,\alpha_1+\alpha_2\).
The dimension product has factors
\(a+1,b+1,(a+b+2)/2\).
A decomposition of \(s\alpha_1+t\alpha_2\) uses
\(j\) copies of \(\alpha_1+\alpha_2\), with
\(0\leq j\leq\min(s,t)\), and then forced numbers of the
two simple roots. This proves the partition formula.
As an internal multiplicity check, for the adjoint highest
weight \(\lambda=\alpha_1+\alpha_2\), Kostant's formula at
\(\mu=0\) has the identity term \(p(\alpha_1+\alpha_2)=2\);
every other Weyl term lies outside \(Q_+\).
It gives the two-dimensional zero weight space \(\mathfrak h\),
consistent with the eight-dimensional decomposition in
Example~\ref{res:G09}.
\end{proof}

\chapter{Cohomology, Levi Decomposition, and Ado}
\label{ch:cohomology-levi-ado}

An extension remembers an action as well as a kernel. Throughout the first
section, \(\mathfrak g\) is a Lie algebra over a field \(k\), and \(M\) is
a specified left \(\mathfrak g\)-module. The Chevalley--Eilenberg convention is
\[
 (db)(x,y)=x b(y)-y b(x)-b([x,y]).
\]
We use the exterior-power complex, its coefficient long exact sequence,
and its comparison with \(\operatorname{Ext}_{U(\mathfrak g)}\) proved in
\cite[Chapter 6, A06--A08]{homological}. In particular, all coefficient
modules in an application of Whitehead's lemmas below are finite dimensional.

\section{Extensions and deformations}

\begin{theorem}[Abelian extensions with a specified action]\label{res:I01}
Equivalence classes of extensions
\[
 0\longrightarrow M\longrightarrow E\longrightarrow\mathfrak g
 \longrightarrow0
\]
whose kernel is abelian and whose induced action on it is the given one
are parametrized by \(H^2(\mathfrak g,M)\). Equivalence induces the identity
on both \(M\) and \(\mathfrak g\). Central extensions are the special case
of the trivial action. An extension has a Lie algebra section if and only
if its class is zero.
\end{theorem}
\begin{proof}
Choose a linear section \(s\). The error
\(\omega(x,y)=[s(x),s(y)]-s([x,y])\) lies in \(M\).
Under the resulting identification \(E=\mathfrak g\oplus M\), the bracket is
\[
 [(x,m),(y,n)]=([x,y],\,xn-ym+\omega(x,y)).
 \tag{9.1}\label{eq:abelian-extension}
\]
Jacobi on three elements of \(\mathfrak g\) has \(M\)-component
\[
 x\omega(y,z)-y\omega(x,z)+z\omega(x,y)
 -\omega([x,y],z)+\omega([x,z],y)-\omega([y,z],x).
\]
Thus it is precisely \(d\omega=0\). Jacobi with one entry in \(M\)
is the module identity, and with two entries in \(M\) it is zero.
Consequently every \(2\)-cocycle produces an extension by
\eqref{eq:abelian-extension}.

Changing the section to \(s+b\) changes \(\omega\) to \(\omega+db\).
Conversely any linear equivalence of extensions has the form
\((x,m)\mapsto(x,m+b(x))\). Substitution in \eqref{eq:abelian-extension}
shows that it preserves brackets exactly when the two cocycles differ
by the corresponding coboundary. This proves both directions of the
classification. The section \(x\mapsto(x,0)\) is a Lie section exactly
when \(\omega=0\); changing sections proves the last assertion.
See \cite[Section 7.6]{weibel} and \cite[Proposition 48.3]{etingof}.
\end{proof}

\begin{proposition}[First-order bracket deformations]\label{res:I02}
Fix a bracket \(\mu(x,y)=[x,y]\). Alternating brackets
\(\mu+t c\) on \(\mathfrak g\otimes_k k[t]/(t^2)\), up to changes of
coordinates \(1+t a\), are classified by
\(H^2(\mathfrak g,\mathfrak g)\), for the adjoint action.
\end{proposition}
\begin{proof}
The coefficient of \(t\) in Jacobi is
\[
 \sum_{\mathrm{cyc}}\bigl([c(x,y),z]+c([x,y],z)\bigr)=-dc(x,y,z).
\]
It vanishes exactly when \(c\) is a cocycle. Pulling back the bracket
by \(T=1+t a\), with inverse \(1-t a\), replaces \(c\) by
\[
 c(x,y)+[a(x),y]+[x,a(y)]-a([x,y])=c(x,y)+da(x,y).
\]
Every allowed change of coordinates has this form. Hence its orbits
are exactly cocycles modulo coboundaries. The action of the kernel
\(t\mathfrak g\) is adjoint; it is not the trivial action of a central
extension. This also recovers the assertion from
Theorem~\ref{res:I01}. Compare \cite[Proposition 48.5]{etingof}.
\end{proof}

\begin{proposition}[Obstructions and formal rigidity]\label{res:I03}
Suppose \(k\) has characteristic zero. A bracket
\(\mu_t=\sum_{i=0}^{N-1}t^i c_i\), where \(c_0=\mu\), that satisfies
Jacobi modulo \(t^N\) extends modulo \(t^{N+1}\) exactly when the class
of
\[
 O_N(x,y,z)=
 \sum_{\substack{i+j=N\\i,j>0}}
 \sum_{\mathrm{cyc}}c_i(c_j(x,y),z)
 \tag{9.2}\label{eq:lie-obstruction}
\]
vanishes in \(H^3(\mathfrak g,\mathfrak g)\). In particular
\(H^2(\mathfrak g,\mathfrak g)=0\) implies that every one-parameter
formal deformation is formally equivalent to the constant bracket.
\end{proposition}
\begin{proof}
The terms involving the unknown \(c_N\) are \(-dc_N\), by the
calculation in Proposition~\ref{res:I02}. The equation to solve is
therefore \(dc_N=O_N\).
Here is the closure check on the obstruction. On alternating maps use
the insertion operation
\[
 (f\circ h)(x_1,\ldots,x_{p+q+1})
 =\sum_{\sigma\in\operatorname{Sh}(q+1,p)}
 \operatorname{sgn}(\sigma)
 f\bigl(h(x_{\sigma(1)},\ldots,x_{\sigma(q+1)}),
          x_{\sigma(q+2)},\ldots,x_{\sigma(p+q+1)}\bigr),
\]
where \(f\) has arity \(p+1\) and \(h\) has arity \(q+1\).
Its commutator
\([f,h]_{\mathrm{NR}}=f\circ h-(-1)^{pq}h\circ f\)
satisfies graded Jacobi. In fact, for \(\ell\) of arity
\(r+1\), the shuffle expansion gives the associator identity
\[
 (f\circ h)\circ \ell-f\circ(h\circ \ell)
   =(-1)^{qr}\bigl((f\circ \ell)\circ h-f\circ(\ell\circ h)\bigr).
\]
On each side the nested-block terms cancel. The remaining terms
correspond to disjoint input sets for the two insertions.
Reordering those blocks swaps \(q+1\) and \(r+1\) inputs,
and their outputs occupy two slots of the alternating map \(f\).
Moving each output past the other block before the final slot
exchange gives the additional signs
\((-1)^{q+1}\), \((-1)^{r+1}\) and \(-1\); their product
with \((-1)^{(q+1)(r+1)}\) is \((-1)^{qr}\).
Expanding the three nested commutators in graded Jacobi pairs
these associators with precisely this sign, so they cancel.
For a binary alternating \(\nu\), its Jacobiator is
\(\tfrac12[\nu,\nu]_{\mathrm{NR}}=\nu\circ\nu\).
The same associator identity with \(f=h=\ell=\nu\) gives
\([\nu,[\nu,\nu]_{\mathrm{NR}}]_{\mathrm{NR}}=0\).
Apply this to the truncated polynomial \(\mu_t\). Its Jacobiator has
no coefficients below \(N\); the coefficient of \(t^N\) in this
identity is \([\mu,O_N]_{\mathrm{NR}}=0\), which is \(dO_N=0\),
up to the fixed overall sign in the exterior-complex convention.
Thus the displayed equation is solvable exactly when its cohomology
class vanishes. Altering a proposed \(c_N\) by a coboundary changes
coordinates and not the obstruction class.

For the rigidity assertion, suppose the first nonzero coefficient is
\(t^m c_m\). It is a \(2\)-cocycle. Write \(c_m=da_m\), and pull back
by \(1-t^m a_m\) to kill it without changing lower coefficients.
Repeat at degrees \(m+1,m+2,\ldots\). The product of these changes
converges in the \(t\)-adic topology: each coefficient is altered
only finitely often, and the constant term remains the identity,
so its inverse exists by formal recursion. The limiting change
removes every positive coefficient.
This proves formal, not analytic-convergence, rigidity.
See \cite[Section 48.1]{etingof} and \cite[Sections 1 and 3]{fialowski}.
\end{proof}

\section{Whitehead's second lemma and Levi factors}

\begin{theorem}[Second Whitehead lemma for trivial coefficients]\label{res:I04}
If \(\mathfrak g\) is finite dimensional and semisimple over a
characteristic-zero field, then \(H^2(\mathfrak g,k)=0\).
\end{theorem}
\begin{proof}
The Killing form \(B\) is nondegenerate by Theorem~\ref{res:B05}.
Given an alternating \(2\)-cocycle \(\omega\), define \(D\) uniquely by
\(\omega(x,y)=B(Dx,y)\). Alternation says that \(D\) is skew-adjoint
for \(B\). Pairing
\(D[x,y]-[Dx,y]-[x,Dy]\) with \(z\), using invariance of \(B\),
and using skew-adjointness, gives
\[
 \omega([x,y],z)+\omega([y,z],x)+\omega([z,x],y)=0.
\]
Thus \(D\) is a derivation. By Theorem~\ref{res:B07}, \(D=\operatorname{ad}a\).
It follows that
\(\omega(x,y)=B(a,[x,y])=d\ell(x,y)\) for
\(\ell(x)=-B(a,x)\). Every cocycle is a coboundary.
This is the trivial-coefficient case of
\cite[Corollary 7.8.12]{weibel}.
\end{proof}

\begin{theorem}[Second Whitehead lemma with coefficients]\label{res:I05}
For a finite-dimensional semisimple characteristic-zero Lie algebra
\(\mathfrak g\) and a finite-dimensional module \(M\),
\[
 H^2(\mathfrak g,M)=0.
\]
Every abelian extension of \(\mathfrak g\) by such a specified module
has a Lie algebra section.
\end{theorem}
\begin{proof}
By Weyl's theorem, Theorem~\ref{res:C05}, \(M\) is a finite direct
sum of simple modules. The exterior-complex differential respects
this finite direct sum. A nontrivial simple summand has zero
cohomology in every degree by the Casimir homotopy in
Theorem~\ref{res:C03}; a trivial simple summand is covered by
Theorem~\ref{res:I04}. Thus degree two vanishes for their direct sum.
The extension assertion follows from Theorem~\ref{res:I01}.
This argument retains finite-dimensionality and does not claim
that all higher cohomology with trivial coefficients vanishes.
See \cite[Corollary 7.8.12]{weibel}.
\end{proof}

\begin{theorem}[Levi decomposition]\label{res:I06}
Let \(\mathfrak g\) be a finite-dimensional Lie algebra over a
characteristic-zero field, with solvable radical \(\mathfrak r\).
There is a semisimple subalgebra \(\mathfrak s\) such that
\[
 \mathfrak g=\mathfrak s\ltimes\mathfrak r,\qquad
 \mathfrak s\xrightarrow{\;\sim\;}\mathfrak g/\mathfrak r.
\]
\end{theorem}
\begin{proof}
Induct on \(\dim\mathfrak r\). If it is zero, take \(\mathfrak s=\mathfrak g\).
Otherwise take the last nonzero term \(\mathfrak a\) of the derived
series of \(\mathfrak r\); if \(\mathfrak r\) is abelian, take
\(\mathfrak a=\mathfrak r\). This is a nonzero abelian characteristic
ideal of \(\mathfrak r\), hence an ideal of \(\mathfrak g\).
The radical of \(\mathfrak g/\mathfrak a\) is
\(\mathfrak r/\mathfrak a\): the latter is solvable, and the quotient
by it is semisimple. Induction gives a Levi subalgebra
\(\overline{\mathfrak s}\) in this quotient.

Let \(E\) be its inverse image in \(\mathfrak g\).
Then \(0\to\mathfrak a\to E\to\overline{\mathfrak s}\to0\)
is an abelian extension, with its actual adjoint action on
\(\mathfrak a\). By Theorem~\ref{res:I05} its class is zero.
A Lie section has image \(\mathfrak s\).
Its intersection with \(\mathfrak r\) is zero because
\(\overline{\mathfrak s}\cap(\mathfrak r/\mathfrak a)=0\) and the
section has no intersection with \(\mathfrak a\).
Moreover its map to \(\mathfrak g/\mathfrak r\) is surjective,
so dimensions give the required vector-space direct sum.
The radical is an ideal, so this is the stated semidirect product.
See \cite[Theorem 7.8.13]{weibel} and
\cite[Theorem 48.9]{etingof}.
\end{proof}

\begin{theorem}[Conjugacy of Levi factors; reductive algebras]\label{res:I07}
Any two Levi factors are conjugate by
\(\exp(\operatorname{ad}z)\) for an element \(z\) of the nilradical
\(\mathfrak n\). The exponential is a finite sum.
A finite-dimensional characteristic-zero Lie algebra whose radical
is its center is the direct sum of a semisimple algebra and its center.
\end{theorem}
\begin{proof}
Fix one Levi factor \(\mathfrak s\) and write the second as
\(\{x+f(x):x\in\mathfrak s\}\), where \(f:\mathfrak s\to\mathfrak r\)
is linear. By Theorem~\ref{res:B08},
\([\mathfrak g,\mathfrak r]\subseteq\mathfrak n\).
Modulo \(\mathfrak n\), therefore,
\(\mathfrak g/\mathfrak n=\mathfrak s\oplus(\mathfrak r/\mathfrak n)\)
has central second summand. Since \(\mathfrak s=[\mathfrak s,\mathfrak s]\),
the homomorphism into that central summand is zero. Thus \(f\)
already takes values in \(\mathfrak n\).

Write \(\gamma_1\mathfrak n=\mathfrak n\) and
\(\gamma_{j+1}\mathfrak n=[\mathfrak n,\gamma_j\mathfrak n]\).
Suppose a conjugate second factor has discrepancy in
\(\gamma_j\mathfrak n\). In
\(\gamma_j\mathfrak n/\gamma_{j+1}\mathfrak n\) the equation for its
bracket is
\[
 f([x,y])=x f(y)-y f(x);
\]
the additional term \([f(x),f(y)]\) belongs to
\(\gamma_{2j}\mathfrak n\subseteq\gamma_{j+1}\mathfrak n\).
By Theorem~\ref{res:C04} there is \(v_j\in\gamma_j\mathfrak n\)
such that \(f(x)=[x,v_j]\) modulo \(\gamma_{j+1}\mathfrak n\).
Conjugating the second factor by \(\exp(\operatorname{ad}v_j)\)
changes the discrepancy by \([v_j,x]\), and kills it modulo
\(\gamma_{j+1}\mathfrak n\). Higher exponential terms have values
in that next term as well.

The lower central series terminates. Also
\(\operatorname{ad}v_j\) is nilpotent on all of \(\mathfrak g\):
its first application lands in the nilpotent ideal \(\mathfrak n\),
and its subsequent applications follow that ideal's lower central
series. Thus all the exponentials used are finite. Their product
is \(\exp(\operatorname{ad}z)\) for \(z\in\mathfrak n\), by the
finite BCH identity of Theorem~\ref{res:A09}. It sends the
second factor to the first.

For the last assertion, Levi decomposition writes the algebra as
\(\mathfrak s\ltimes Z(\mathfrak g)\). The action on the center is
zero, hence the sum is direct. Conversely such a direct sum has
exactly its abelian central summand as radical. This is the
meaning of \emph{reductive} used here; it is not an assertion that
every representation of a reductive algebra is completely reducible.
The conjugacy statement expands the Malcev observation following
\cite[Theorem 7.8.13]{weibel}.
\end{proof}

\section{Finite faithful representations}

\begin{proposition}[A derivation-stable nilpotent truncation]\label{res:I08}
Every finite-dimensional nilpotent Lie algebra \(\mathfrak n\)
has a finite-dimensional faithful representation by nilpotent
operators. The construction can be made compatible with every
derivation of \(\mathfrak n\).
\end{proposition}
\begin{proof}
Choose a basis adapted to the lower central series. Give a basis
element weight \(j\) if it belongs to
\(\gamma_j\mathfrak n\) but not to \(\gamma_{j+1}\mathfrak n\).
For an ordered PBW monomial give its weight as the sum of the
weights of its factors. Let \(F_{\geq N}U(\mathfrak n)\) be the span
of monomials of weight at least \(N\).

This is a two-sided ideal. In fact, inserting a new factor and
reordering by PBW either preserves the sum of weights or replaces
two factors by a bracket. The latter has weight at least their
sum because
\([\gamma_i\mathfrak n,\gamma_j\mathfrak n]\subseteq
\gamma_{i+j}\mathfrak n\). No resulting term has smaller weight.
Every derivation \(\delta\) preserves each \(\gamma_j\mathfrak n\),
as follows by induction from the derivation identity. Its extension
to \(U(\mathfrak n)\), obtained by applying it in every factor, also
never decreases weight. Thus it preserves this ideal.

The quotient \(V_N=U(\mathfrak n)/F_{\geq N}\) is finite dimensional:
the positive integer weights bound each exponent in a monomial of
weight less than \(N\). If \(\gamma_{c+1}\mathfrak n=0\), take
\(N>c\). The distinct PBW vectors \(1\) and all the individual
basis vectors survive and remain independent. Hence left
multiplication \(L_x\) is faithful on \(\mathfrak n\), since
\(L_x(1)=x\). It raises weight by at least one, so \(L_x^N=0\).
Finally every \(\delta\) acts on \(V_N\), kills \(1\), and satisfies
\([D_\delta,L_x]=L_{\delta x}\).
This is a PBW proof of the faithful nilpotent-module conclusion
of \cite[Proposition 50.5 and Corollary 50.7]{etingof}.
\end{proof}

\begin{proposition}[A reductive--nilpotent enlargement]\label{res:I09}
Every finite-dimensional complex Lie algebra embeds into an algebra
\[
 \widehat{\mathfrak g}=\mathfrak k\ltimes\mathfrak n,\qquad
 \mathfrak k=\mathfrak s\oplus\mathfrak t,
\]
where \(\mathfrak s\) is semisimple, \(\mathfrak t\) is abelian
and acts semisimply, and \(\mathfrak n\) is nilpotent.
The central toral actions introduced in the construction have
integer weights in a chosen basis of \(\mathfrak t\).
\end{proposition}
\begin{proof}
We give an entirely Lie-algebraic version of the reduction in
\cite[Proposition 50.3]{etingof}. Inductively maintain a decomposition
\(\mathfrak k\ltimes\mathfrak a\), with \(\mathfrak a\) solvable,
\(\mathfrak k=\mathfrak s\oplus\mathfrak t\) reductive, and
\(\mathfrak t\) acting diagonally on \(\mathfrak a\).
Such an algebra starts the induction: take a Levi factor for
\(\mathfrak k=\mathfrak s\), and \(\mathfrak a=\mathfrak r\).
Weyl's theorem makes the initial \(\mathfrak k\)-module
\(\mathfrak a\) completely reducible. More generally, a module
on which the center acts diagonally is a direct sum of its central
weight spaces; on each, Weyl's theorem for \(\mathfrak s\)
gives complete reducibility.

Let \(\mathfrak n_0\) be the nilradical of \(\mathfrak a\).
Use the nonnegative integer \(q=\dim(\mathfrak a/\mathfrak n_0)\)
as decreasing measure. If \(q=0\), we are done.
Every derivation of \(\mathfrak a\) acts trivially on
\(\mathfrak a/\mathfrak n_0\), by Theorem~\ref{res:B08}.
Complete reducibility therefore supplies a
\(\mathfrak k\)-invariant element \(d\in\mathfrak a\)
whose image in that quotient is nonzero. Choose a linear
functional \(\chi\) on \(\mathfrak a/\mathfrak n_0\) with
\(\chi(d)=1\), and regard it as a functional on \(\mathfrak a\).
It is a character because
\([\mathfrak a,\mathfrak a]\subseteq\mathfrak n_0\).

Decompose \(\mathfrak a=\bigoplus_\beta\mathfrak a_{[\beta]}\)
into generalized eigenspaces of \(D=\operatorname{ad}d\).
The derivation identity gives
\[
 [\mathfrak a_{[\beta]},\mathfrak a_{[\gamma]}]
 \subseteq\mathfrak a_{[\beta+\gamma]}.
 \tag{9.3}\label{eq:additive-generalized-weights}
\]
Indeed \((D-(\beta+\gamma))^m[x,y]\) is the binomial sum of brackets
of \((D-\beta)^i x\) and \((D-\gamma)^{m-i}y\), and is zero for
large \(m\). The finitely many eigenvalues generate a free
abelian subgroup \(\Gamma\subset\mathbb C\); choose a
\(\mathbb Z\)-basis \(\alpha_1,\ldots,\alpha_m\).
Write \(\beta=\sum_j b_{\beta j}\alpha_j\), with integers
\(b_{\beta j}\).
An abelian algebra \(\mathfrak t_1=\mathbb C^m\) acts on
\(\mathfrak a_{[\beta]}\) by
\((z_j)\mapsto\sum_j b_{\beta j}z_j\).
Equation~\eqref{eq:additive-generalized-weights} proves that these
operators are derivations. They commute with \(\mathfrak k\),
because \(d\) is \(\mathfrak k\)-invariant.

For \(\alpha=(\alpha_1,\ldots,\alpha_m)\in\mathfrak t_1\),
the action of \(\alpha\) is the semisimple part \(S\) of \(D\).
Both \(\mathfrak n_0\) and \(\ker\chi\) are stable under
\(\mathfrak t_1\). For the first, \(D\)-stability gives a
decomposition into generalized eigenspaces; for the second,
\(D\) acts trivially on \(\mathfrak a/\mathfrak n_0\), so that
only weight zero occurs in this quotient.
Also \(\mathfrak t_1\) kills \(d\), since \(D d=0\).

In \((\mathfrak k\oplus\mathfrak t_1)\ltimes\mathfrak a\), put
\[
 y=d-\alpha,\qquad
 \mathfrak a'=\ker\chi\oplus\mathbb C y.
\]
It is an ideal complementary to
\(\mathfrak k'=\mathfrak k\oplus\mathfrak t_1\).
For example \([y,x]=(D-S)x\) for \(x\in\ker\chi\), and this
lies in \(\mathfrak n_0\), since both \(D\) and \(S\) act
trivially on \(\mathfrak a/\mathfrak n_0\).
All other required brackets follow from the stability just proved.
The new algebra contains the old \(\mathfrak k\ltimes\mathfrak a\);
we have only adjoined \(\mathfrak t_1\).

The ideal \(\mathfrak m=\mathfrak n_0+\mathbb C y\) of
\(\mathfrak a'\) is nilpotent. To check this carefully, every
\(x\in\mathfrak n_0\) acts nilpotently on \(\mathfrak a'\):
its first bracket lands in \(\mathfrak n_0\), where repeated
bracketing terminates. The element \(y\) acts on \(\ker\chi\)
by the nilpotent operator \(D-S\) and kills itself.
Lie's theorem simultaneously upper triangularizes the adjoint
representation of the solvable algebra \(\mathfrak m\) on
\(\mathfrak a'\). Its diagonal characters vanish on
\(\mathfrak n_0\) and on \(y\), by the two nilpotency observations,
so every element has zero diagonal. Engel's theorem gives the
claim. Hence the new nilradical contains \(\mathfrak m\), and
\[
 \dim(\mathfrak a'/\operatorname{nilrad}\mathfrak a')
 \leq \dim(\mathfrak a/\mathfrak n_0)-1.
\]

Finally the central torus of \(\mathfrak k'\) acts diagonally
on \(\mathfrak a'\): it does so on the old generalized
eigenspace blocks, and it kills the extra line \(\mathbb C y\).
Its new weights are the integers \(b_{\beta j}\).
Thus complete reducibility needed in the next step is still
available. The integer \(q\) decreases until it is zero.
This proves the enlargement without assuming that arbitrary
modules of a reductive algebra are semisimple, or assuming an
integration theorem for an algebraic group.
\end{proof}

\begin{proposition}[The semidirect-product representation]\label{res:I10}
If \(\mathfrak k\) is a finite-dimensional complex reductive algebra
and it acts by derivations on a finite-dimensional nilpotent
algebra \(\mathfrak n\), then \(\mathfrak k\ltimes\mathfrak n\)
has a finite-dimensional faithful module.
\end{proposition}
\begin{proof}
Take \(V_N\) from Proposition~\ref{res:I08}. The formula
\[
 (a,x)\longmapsto D_a+L_x
\]
is a representation: \([D_a,D_b]=D_{[a,b]}\),
\([D_a,L_y]=L_{a y}\), and \([L_x,L_y]=L_{[x,y]}\)
give exactly the semidirect-product bracket.
If \(D_a+L_x\) is zero, applying it to \(1\) gives \(x=0\).
Its restriction to the surviving copy of \(\mathfrak n\)
then says that \(a\) acts as the zero derivation on \(\mathfrak n\).

To detect this possible remaining kernel, take a faithful module
\(W\) of \(\mathfrak k=\mathfrak s\oplus\mathfrak t\).
The adjoint module of \(\mathfrak s\) is faithful because its center
is zero. If \(t_1,\ldots,t_m\) is a basis of \(\mathfrak t\),
let \(t_j\) act on \(\mathbb C^m\) by the diagonal matrix with
its only nonzero entry \(1\) in position \(j\); let \(\mathfrak s\)
act trivially on that summand and \(\mathfrak t\) trivially on
the adjoint summand. Their direct sum is the required \(W\).
Inflate it along the projection
\(\mathfrak k\ltimes\mathfrak n\to\mathfrak k\).
On \(V_N\oplus W\), the kernel condition first gives \(x=0\)
and then \(a=0\). This also covers \(\mathfrak n=0\) and pure
reductive direct factors; no faithfulness of the action of
\(\mathfrak k\) on \(\mathfrak n\) has been presumed.
Compare the final construction in \cite[Theorem 50.8]{etingof}.
\end{proof}

\begin{theorem}[Ado over the complex and real fields]\label{res:I11}
Every finite-dimensional complex or real Lie algebra admits an
injective homomorphism into \(\mathfrak{gl}(V)\) for a
finite-dimensional vector space over the same field.
\end{theorem}
\begin{proof}
Over \(\mathbb C\), embed into the enlargement of
Proposition~\ref{res:I09}, apply Proposition~\ref{res:I10},
and restrict the faithful representation to the original subalgebra.
For a real algebra, apply the complex assertion to
\(\mathfrak g\otimes_{\mathbb R}\mathbb C\) and regard its
representation as a real vector space. The original map
\(\mathfrak g\to\mathfrak g\otimes_{\mathbb R}\mathbb C\)
is injective, so its action remains faithful. The zero algebra
may act on the zero-dimensional space, or trivially on a line.
The arbitrary characteristic-zero version is proved in
Proposition~\ref{res:J02}, with its descent step included there.
\end{proof}

\chapter{Forms and Characteristic Boundaries}
\label{ch:real-forms}

The complex classification does not identify a real algebra merely from
its Dynkin diagram. A real form also specifies a conjugation. We first
translate conjugations into involutions of a compact form, then solve
the resulting finite classification problem. The final section gives
counterexamples to some hypotheses used earlier.

\section{Automorphisms and descent}

\begin{theorem}[Automorphisms of a complex semisimple algebra]\label{res:J01}
For a complex semisimple algebra with Dynkin diagram \(D\), including
its connected components,
\[
 \operatorname{Aut}(\mathfrak g)
 =\operatorname{Int}(\mathfrak g)\rtimes\operatorname{Aut}(D).
\]
Here \(\operatorname{Int}\) is generated by exponentials of inner
derivations; automorphisms of \(D\) may permute isomorphic components.
\end{theorem}
\begin{proof}
A permutation of the simple generators preserving the Cartan matrix
preserves every Serre relation, so gives a Lie automorphism by
Theorem~\ref{res:F03}. It gives a section of the asserted quotient.
For an arbitrary automorphism, Cartan conjugacy first lets us compose
with an inner automorphism to preserve a fixed Cartan. A Weyl-group
representative then makes it preserve the positive roots. Its action
on simple roots is an automorphism of \(D\).
Remove that diagram automorphism.

The remaining map fixes the Cartan pointwise, sends each simple
root vector \(e_i\) to \(c_i e_i\), and sends \(f_i\) to
\(c_i^{-1}f_i\), because \([e_i,f_i]=h_i\).
Choose \(h\) in the Cartan with \(\exp(\alpha_i(h))=c_i\):
choose complex logarithms of the \(c_i\), and use that the
simple roots are a basis. Then \(\exp(\operatorname{ad}h)\)
has precisely these effects. The simple generators generate the
algebra, so the automorphism is this inner one.

Finally a nontrivial diagram automorphism permutes at least two
fundamental weights (also when it permutes simple factors).
The corresponding finite-dimensional simple modules are not
isomorphic by Theorem~\ref{res:G06}. Twisting by an inner
automorphism, however, preserves every module's isomorphism class:
for \(\exp(\operatorname{ad}x)\), the intertwiner is
\(\exp(\rho(x))\), as the exponential power series gives
\(\rho(\exp(\operatorname{ad}x)y)
=\exp(\rho(x))\rho(y)\exp(-\rho(x))\);
compose these intertwiners for a product of inner exponentials.
Thus no nontrivial diagram automorphism is inner. The section
has trivial intersection with the inner subgroup, proving the
semidirect product.
See \cite[Section 14.3]{humphreys} and
\cite[Proposition 40.3]{etingof}.
\end{proof}

\begin{proposition}[Finite Galois descent and characteristic-zero Ado]\label{res:J02}
Every finite reduced root system has a split semisimple model over
every characteristic-zero field, obtained from its Chevalley
integer lattice.
Let \(L/K\) be a finite Galois extension and let \(\mathfrak a\) be
a Lie algebra over \(K\). The \(K\)-forms of
\(\mathfrak a_L=L\otimes_K\mathfrak a\) are classified by
semilinear Galois descent data, or equivalently by
\[
 H^1(\operatorname{Gal}(L/K),\operatorname{Aut}_L(\mathfrak a_L)).
\]
The equivalence is by coboundary, not by equality of cocycles.
In addition, every finite-dimensional Lie algebra over an arbitrary
characteristic-zero field has a faithful finite-dimensional representation.
\end{proposition}
\begin{proof}
Start with the Chevalley basis of Theorem~\ref{res:F05}. Its
integer span \(\mathfrak g_{\mathbb Z}\) is closed under brackets.
All coefficients of its Jacobi and alternation identities are
integers and are zero after embedding in \(\mathbb C\), hence are
zero already over \(\mathbb Z\). Thus
\(\mathfrak g_K=K\otimes_{\mathbb Z}\mathfrak g_{\mathbb Z}\)
is a Lie algebra with that same basis over any field \(K\).
Its Killing matrix has integer entries, because each adjoint
matrix does. Its determinant is a nonzero integer: over
\(\mathbb C\) it is the nondegenerate Killing form of the
constructed semisimple algebra. In characteristic zero this
determinant remains nonzero, so \(\mathfrak g_K\) is semisimple
by Theorem~\ref{res:B05}. Its Cartan basis \(h_i\) acts diagonally
with the integer eigenvalues \(\alpha(h_i)\); the root spaces
are already the individual \(Kx_\alpha\), and the zero space
is \(Kh_1+\cdots+Kh_r\). This is the asserted split model.
The descent below applies in particular to its forms trivialized
over the stated extension \(L\); it does not identify all
nonsplit forms with the split model itself.

Write \(\Gamma=\operatorname{Gal}(L/K)\), and let \(\gamma_0\)
act on scalar coefficients. A semilinear bracket-preserving action
has the form \(\widetilde\gamma=c_\gamma\gamma_0\).
The action law is exactly
\(c_{\gamma\delta}=c_\gamma\gamma_0(c_\delta)\).
Changing an identification with \(\mathfrak a_L\) by \(b\)
replaces it by
\(c'_\gamma=b c_\gamma\gamma_0(b)^{-1}\).

We recall why the fixed subspace actually descends the vector space.
Distinct automorphisms of a field are linearly independent as functions:
in a minimal relation, compare the relation at \(au\) with that at
\(u\) multiplied by one automorphism's value at \(a\); choosing \(a\)
on which two automorphisms differ gives a shorter relation.
For a nonzero semilinear \(\Gamma\)-space \(V\), the expression
\(\sum_{\gamma\in\Gamma}\widetilde\gamma(a v)\) is invariant.
For some \(a,v\) it is nonzero, by that independence, applied
after an \(L\)-linear coordinate functional.
More generally, let \(W\) be the \(L\)-span of all invariant
vectors. If \(V/W\ne0\), the same expression computed in the
quotient produces a nonzero invariant sum that lifts to an
invariant vector of \(V\), a contradiction. Thus \(W=V\).
Choose an \(L\)-basis from the invariant vectors. An invariant
linear combination in that basis has all its coefficients in
\(K\). Hence
\(L\otimes_K V^\Gamma\to V\) is an isomorphism.
Taking fixed points of the bracket gives the asserted descended
Lie algebra, and recovers all forms. Any finite-dimensional form
that splits over an algebraic closure already splits over some
finite Galois extension: finitely many entries of an isomorphism
and its inverse generate a finite extension, which can be enlarged
to its Galois closure in characteristic zero.

Here is the additional descent needed for Ado; no embedding of the
whole original field into \(\mathbb C\) is required.
Let \(F\subseteq K\) be the finitely generated extension of
\(\mathbb Q\) generated by the finitely many structure constants
in a chosen basis of \(\mathfrak g\). The field \(F\) embeds into
\(\mathbb C\): send a transcendence basis to algebraically
independent complex numbers and extend through its finite
algebraic extension. By Theorem~\ref{res:I11}, the corresponding
complex algebra has a faithful representation of some dimension
\(n\). Unknown matrices \(X_i\in M_n\) for the basis elements
satisfy finitely many equations
\[
 [X_i,X_j]=\sum_\ell c_{ij}^{\ell}X_\ell.
\]
Faithfulness is the nonvanishing of one maximal minor \(\Delta\)
of the matrix whose columns are the entries of the \(X_i\).
Adjoin a variable \(u\) and the equation \(u\Delta-1=0\).
This polynomial system over \(F\) has a complex solution and
therefore its ideal is proper. A maximal ideal containing it
has residue field finite algebraic over \(F\), by Zariski's lemma.
For completeness, that lemma follows by choosing a transcendence
basis in a field that is finitely generated as an algebra:
clearing the denominators of the remaining algebraic generators
would make the field integral over
\(F[t_1,\ldots,t_r,1/f]\). An integral subring of a field over
which the field is integral must itself be a field, whereas for
\(r>0\) the displayed ring misses the inverse of some \(t_1-a\),
\(a\in F\). Thus \(r=0\), and the algebraic extension is finite.

Embed this residue field in an algebraic closure of \(K\).
Its compositum with \(K\) is a finite extension \(K'/K\), and
the resulting matrices are a faithful representation of
\(\mathfrak g_{K'}\). Restrict the scalars of the representation
space to \(K\). A nonzero original vector remains nonzero after
base extension, so the \(K\)-representation is faithful.
This argument explains why merely collecting finitely many
complex coefficients, without the polynomial-system step,
would not justify descent.
For forms compare \cite[Proposition 40.4]{etingof}.
\end{proof}

\section{The compact form and its involutions}

\begin{theorem}[Conjugations and the compact form]\label{res:J03}
A real form of a complex Lie algebra is the fixed algebra of an
antilinear bracket-preserving involution, and its complexification
is the original algebra. Every complex semisimple algebra has
a compact real form \(\mathfrak u\), on which the Killing form
is negative definite.
\end{theorem}
\begin{proof}
For an antilinear involution \(\sigma\), write
\[
 x=\frac{x+\sigma x}{2}
   +i\,\frac{x-\sigma x}{2i}.
\]
Both displayed summands before multiplication by \(i\) are fixed.
Their real direct sum is the whole complex vector space,
so complexification of the fixed algebra is an isomorphism.
Conversely the usual coefficient conjugation in a complexification
has exactly the original real algebra as fixed points.
Isomorphisms of real forms are precisely complex automorphisms
intertwining their conjugations.

Take the Chevalley normalization of Theorem~\ref{res:F05}
(equivalently its compact-normalization construction):
\([e_\alpha,e_{-\alpha}]=h_\alpha\),
\(B(e_\alpha,e_{-\alpha})>0\), and the conjugate-linear rule
\[
 \omega(h_i)=-h_i,\qquad
 \omega(e_i)=-f_i,\qquad \omega(f_i)=-e_i
\]
preserves the defining Serre relations. It is an involution.
Choose the root vectors compatibly so
\(\omega(e_\alpha)=-e_{-\alpha}\).
The fixed real space has basis
\[
 ih_i,\qquad
 e_\alpha-e_{-\alpha},\quad i(e_\alpha+e_{-\alpha})
 \quad(\alpha>0).
\]
It is a Lie algebra because it is a fixed algebra.
On the real span of the coroots the Killing form is positive
definite, by the positive-definite root form. Multiplication by
\(i\) makes it negative. Different pairs
\(\{\alpha,-\alpha\}\) are orthogonal, and on each of the two
displayed real root vectors its value is
\(-2B(e_\alpha,e_{-\alpha})<0\); their cross term is zero.
Thus \(B|_{\mathfrak u}\) is negative definite.
This construction uses the root normalization proved earlier,
not the existence of a compact group to prove that normalization.
Compare \cite[Proposition 41.1]{etingof} and
\cite[Section VI.1]{knapp}.
\end{proof}

\begin{theorem}[Polar reduction of real structures]\label{res:J04}
Fix the compact conjugation \(\omega\). Real forms are in bijection
with conjugacy classes of involutions
\(\theta\in\operatorname{Aut}_{\mathbb R}(\mathfrak u)\), by
\(\theta\mapsto\sigma=\theta\omega\).
If \(\mathfrak u=\mathfrak u^+\oplus\mathfrak u^-\) are its
eigenspaces, the corresponding real form is
\[
 \mathfrak g_\theta=\mathfrak u^+\oplus i\mathfrak u^-.
\]
The fixed complex algebra \(\mathfrak k=\mathfrak g^\theta\)
is reductive. The Killing signature of the real form, recorded
as positive dimension minus negative dimension, is
\(\dim_{\mathbb C}\mathfrak g-2\dim_{\mathbb C}\mathfrak k\).
\end{theorem}
\begin{proof}
Use the positive Hermitian form
\(\langle x,y\rangle=-B(x,\omega y)\).
Invariance of \(B\) gives, for every complex automorphism \(a\),
\[
 a^\dagger=\omega a^{-1}\omega,\qquad
 \omega(a)=(a^\dagger)^{-1}.
\]
Write any real structure as \(\sigma=a\omega\).
The condition \(\sigma^2=1\) says \(a\omega(a)=1\), hence
\(a=a^\dagger\). Its nonzero eigenvalues are real, and
\([V_\beta,V_\gamma]\subseteq V_{\beta\gamma}\).
Therefore \(|a|^t\), which multiplies \(V_\beta\) by
\(|\beta|^t\), is a Lie automorphism for every real \(t\):
the scalar on a bracket is the product of its scalars on the
two entries. Put \(\theta=a|a|^{-1}\).
It has eigenvalues \(\pm1\), is unitary, and commutes with
\(\omega\). Conjugating \(a\omega\) by \(|a|^{-1/2}\) gives
\(\theta\omega\). This proves existence of the reduction
without assuming that an arbitrary matrix power preserves brackets.

Suppose \(\theta\omega\) and \(\theta'\omega\) are intertwined
by \(x\). Thus \(\theta'=x\theta\omega(x)^{-1}=x\theta x^\dagger\).
Writing \(P=x^\dagger x\), the equation \((\theta')^2=1\)
gives \(\theta P\theta=P^{-1}\).
The operator \(v=xP^{-1/2}\) is a unitary Lie automorphism,
so it preserves \(\mathfrak u\), and
\[
 v\theta v^{-1}=xP^{-1/2}\theta P^{1/2}x^{-1}
              =x\theta P x^{-1}=\theta'.
\]
The converse is immediate. This proves the asserted bijection.

The real fixed algebra \(\mathfrak u^+\) has the invariant positive
form \(-B\). In any compact algebra with such a form, the
orthogonal complement of an ideal is an ideal. Moreover an
abelian ideal is central: for \(a,b\) in it and arbitrary \(x\),
\(\langle[x,a],b\rangle=\langle x,[a,b]\rangle=0\), while
\([x,a]\) still belongs to the ideal. Taking a last nonzero
derived ideal of the radical, splitting its orthogonal complement,
and inducting shows that the entire radical is central.
Thus the algebra, and its complexification \(\mathfrak k\),
are reductive by Theorem~\ref{res:I07}.
Finally the two eigenspaces are \(B\)-orthogonal;
\(B\) is negative on \(\mathfrak u^+\) and positive on
\(i\mathfrak u^-\). This gives the signature formula.
See \cite[Theorem 41.5]{etingof}.
\end{proof}

\begin{proposition}[A stable Cartan and Vogan data]\label{res:J05}
An involution of a compact form has a stable complex Cartan
\(\mathfrak h=\mathfrak h^+\oplus\mathfrak h^-\), with
\(\mathfrak h^+\) a Cartan in \(\mathfrak k=\mathfrak g^\theta\),
and no root vanishes identically on \(\mathfrak h^+\).
One can choose positive roots preserved by \(\theta\).
Thus \(\theta\) is encoded by an involution of the Dynkin
diagram, with a sign \(+1\) or \(-1\) on each fixed vertex.
Paired vertices need no additional sign.
\end{proposition}
\begin{proof}
Choose a maximal abelian subalgebra \(\mathfrak t^+\) of
\(\mathfrak u^+\). Choose a maximal subspace
\(\mathfrak t^-\subseteq\mathfrak u^-\) that commutes with
\(\mathfrak t^+\) and with itself. Then
\(\mathfrak t=\mathfrak t^+\oplus\mathfrak t^-\) is maximal
abelian in \(\mathfrak u\). Indeed an element centralizing it
has a positive and a negative part under \(\theta\); maximality
puts those parts in \(\mathfrak t^+\) and \(\mathfrak t^-\).
All its adjoint operators are commuting skew-adjoint operators,
so \(\mathfrak h=\mathfrak t_{\mathbb C}\) is toral and
self-centralizing, hence Cartan. The same argument in
\(\mathfrak u^+\) proves that
\(\mathfrak h^+=\mathfrak t^+_{\mathbb C}\) is Cartan in
\(\mathfrak k\).

If a root \(\alpha\) vanished on \(\mathfrak h^+\), its
coroot would lie in \(\mathfrak h^-\), using the
\(\theta\)-invariant root inner product. Hence
\(\theta\alpha=-\alpha\). For nonzero root vectors,
\(\theta(e_\alpha)=c e_{-\alpha}\) and
\(\theta(e_{-\alpha})=c^{-1}e_\alpha\).
The vector \(e_\alpha+c e_{-\alpha}\) would be fixed by
\(\theta\), would centralize \(\mathfrak h^+\), and would
not belong to it. This contradicts its being Cartan in
the reductive \(\mathfrak k\).

Choose a generic element in \(i\mathfrak t^+\), avoiding
the finitely many root hyperplanes. It orders the roots and
is fixed by \(\theta\), so the resulting positive system is
preserved. On simple generators
\(\theta h_i=h_{\sigma(i)}\).
If \(\sigma(i)=i\), involutivity forces
\(\theta e_i=\varepsilon_i e_i\),
\(\theta f_i=\varepsilon_i f_i\), \(\varepsilon_i=\pm1\).
If \(i,\sigma(i)\) are paired, choose \(e_{\sigma(i)}=\theta e_i\)
and the compatible negative root vector; the second application
of \(\theta\) is then the identity. Conversely every such
assignment respects the Serre relations, commutes with compact
conjugation, and squares to one.
These signs and pairs are the Vogan data.
See \cite[Proposition 41.7 and Lemma 41.8]{etingof}.
\end{proof}

\section{Equivalence and the classical list}

\begin{proposition}[Equivalence of the finite data]\label{res:J06}
Two normalized Vogan data define isomorphic real forms exactly when
their involutions are conjugate by a compact Cartan normalizer,
after possibly moving one stable Cartan by the connected fixed group.
Thus the test is finite: test root-system automorphisms and normalize
the scalar multiples on paired root vectors. In the inner case it
is the action of \(W\rtimes\operatorname{Aut}(D)\) on
\(\operatorname{Hom}(Q,\{\pm1\})\), where \(Q\) is the root lattice.
Writing a sign as \((-1)^{b_j}\), the simple reflection at \(i\) acts by
\[
 b_j\longmapsto b_j-a_{ij}b_i\pmod2,\qquad
 a_{ij}=\alpha_j(h_i).
 \tag{10.1}\label{eq:vogan-reflection}
\]
In particular a black fixed vertex changes the colors of its
fixed neighbors for which \(a_{ij}\) is odd.
\end{proposition}
\begin{proof}
We first justify the Cartan adjustment in the assertion.
Maximal abelian subalgebras of a compact reductive algebra
\(\mathfrak v\) are conjugate under its connected inner group.
Here is a useful finite-dimensional proof. Let \(H\) be the
identity component of the group of orthogonal automorphisms
of \(\mathfrak v\) that fix \(Z(\mathfrak v)\) pointwise.
It is a closed subgroup of an orthogonal group and hence compact.
Its Lie algebra consists of skew-adjoint derivations that vanish
on the center. Theorem~\ref{res:B07}, applied to the semisimple
summand of \(\mathfrak v_{\mathbb C}\), identifies these
derivations with \(\operatorname{ad}\mathfrak v\).
Choose regular elements \(X,Y\)
in two maximal abelian subalgebras; regular means that their
centralizers are those subalgebras, obtained by avoiding finitely
many weight hyperplanes. Maximize
\(\langle\operatorname{Ad}(g)X,Y\rangle\) in the compact group.
Differentiation in the direction \(Z\) gives
\(\langle Z,[\operatorname{Ad}(g)X,Y]\rangle=0\) for every \(Z\).
Thus \(\operatorname{Ad}(g)X\) commutes with \(Y\), and belongs
to the second maximal abelian subalgebra. Its centralizer is
the first subalgebra transported by \(g\), so the two maximal
abelian subalgebras coincide after transport.
Every element of \(H\) is a finite product of inner
exponentials. This follows by the inverse
function theorem applied to a product of exponentials of a
basis at the identity; their image contains a neighborhood,
and a subgroup containing a neighborhood is open and equals
the identity component.

Suppose \(c\theta c^{-1}=\theta'\) in
\(\operatorname{Aut}(\mathfrak u)\).
It carries \(\mathfrak t^+\) to a maximal abelian subalgebra
of \(\mathfrak u^{\theta'}\).
Apply the preceding result in this fixed compact algebra to
make \(c\mathfrak t^+=(\mathfrak t')^+\).
The required correcting automorphisms are exponentials of
elements of \(\mathfrak u^{\theta'}\), so commute with \(\theta'\).
By Proposition~\ref{res:J05}, no root vanishes on these
positive Cartan parts. Their centralizers in \(\mathfrak g\)
are consequently the full stable Cartans. Thus \(c\) can
be made to carry the full Cartan to the full Cartan.

The normalizer acts on roots through a finite root-system
automorphism. Its kernel scales each simple root vector
by a unit scalar and is the compact torus: solve the
simple-root logarithms as in Theorem~\ref{res:J01}.
There is therefore an explicit finite equivalence test.
For each root-system automorphism \(w\), use its lift formed
from root \(\mathfrak{sl}_2\) reflections and diagram permutations.
Retain it only if it conjugates the diagram action of the
first datum to that of the second. Compare the resulting
multipliers on the simple root vectors. Conjugation by the
kernel torus changes a multiplier by \(z_j/z_{\sigma(j)}\).
It cannot change a fixed-vertex sign, and it can normalize
both multipliers on any paired vertices, since they are
reciprocals. Thus comparison of the fixed signs is sufficient
after this finite set of tests.

For an inner datum the involution fixes the Cartan and is
the character \(b\) on \(Q\). A reflection changes its value
on \(\alpha_j\) to its value on
\(\alpha_j-a_{ij}\alpha_i\), proving
\eqref{eq:vogan-reflection}. It does not change \(b_i\), since
\(a_{ii}=2\). The same calculation at a fixed simple root in
an outer diagram gives the stated fixed-neighbor move,
after renormalizing paired generators.
This proves both the conjugacy criterion and the elementary
moves used below. Compare \cite[Lemma 42.1]{etingof}.
\end{proof}

\begin{theorem}[Classical simple real Lie algebras]\label{res:J07}
Up to the low-rank identifications specified below, the real forms
of the simple classical complex algebras are the following.
Parameters \(p,q\) satisfy \(p+q=n\), except in the orthogonal
row where \(p+q=N\); interchange of \(p,q\) gives the same form.
\[
\begin{array}{c|l|l}
 \mathfrak g&\text{real forms}&
       (\text{maximal compact algebra})_{\mathbb C}\\ \hline
 \mathfrak{sl}_n&
 \mathfrak{su}(p,q)&
 \mathfrak s(\mathfrak{gl}_p\oplus\mathfrak{gl}_q)\\
 &\mathfrak{sl}_n(\mathbb R)&\mathfrak{so}_n\\
 &\mathfrak{su}^*(2m),\ n=2m&
             \mathfrak{sp}_{2m}\\ \hline
 \mathfrak{so}_N&
 \mathfrak{so}(p,q)&\mathfrak{so}_p\oplus\mathfrak{so}_q\\
 &\mathfrak{so}^*(2n),\ N=2n&\mathfrak{gl}_n\\ \hline
 \mathfrak{sp}_{2n}&
 \mathfrak{sp}(p,q)&\mathfrak{sp}_{2p}\oplus\mathfrak{sp}_{2q}\\
 &\mathfrak{sp}_{2n}(\mathbb R)&\mathfrak{gl}_n
\end{array}
\]
The compact cases are included by \(q=0\).
The quaternionic row \(\mathfrak{su}^*(2m)\) is additional for
\(m\geq2\).
Every complex simple algebra, regarded as a real algebra, is
also real simple; its complexification has two isomorphic
simple components. These form the other branch of the
classification of real simple algebras.
\end{theorem}
\begin{proof}
The proof classifies involutions of the compact forms from
Theorem~\ref{res:J04}; all reductions used here are ordinary
spectral or bilinear-form calculations.

For type \(A_{n-1}\), an inner involution of
\(\mathfrak{su}(n)\) is \(\operatorname{Ad}A\).
Commutation with compact conjugation says \(A^\dagger A\)
is scalar, so scale \(A\) to be unitary.
Its square is scalar; multiplying by a unit scalar makes
\(A^2=1\). The Hermitian spectral theorem puts it in the
form \(I_p\oplus(-I_q)\).
The real structure on \(\mathfrak{sl}_n\) is
\(X\mapsto-A X^\dagger A\); its fixed algebra is exactly
\(\mathfrak{su}(p,q)\).
Its compact fixed algebra consists of the block-diagonal
matrices, giving the first row of the table.

For \(n>2\) the nontrivial diagram action is represented by
\(X\mapsto-X^T\). An outer involution therefore has the form
\(-A X^T A^{-1}\). Its square is inner by \(A A^{-T}\);
this is scalar exactly when \(A^T=\varepsilon A\), with
\(\varepsilon=\pm1\). In the compact case \(A\) can be
chosen unitary. A unitary symmetric matrix is congruent
to \(I\): its associated antiunitary map squares to \(1\),
whose fixed real orthonormal basis is obtained by writing
each vector as fixed real and imaginary parts and applying
real Gram--Schmidt. A unitary skew-symmetric matrix is
congruent to
\(J=\left(\begin{smallmatrix}0&I_m\\-I_m&0\end{smallmatrix}\right)\):
for its antiunitary map \(j\), choose a unit vector \(v\);
then \(v,jv\) are orthogonal because \(j^2=-1\), split off
this pair, and induct. Thus the skew case exists only
when \(n=2m\).
The resulting conjugations are \(X\mapsto\overline X\)
and \(X\mapsto J\overline XJ^{-1}\).
Their fixed algebras are
\(\mathfrak{sl}_n(\mathbb R)\) and
\(\mathfrak{su}^*(2m)=\mathfrak{sl}_m(\mathbb H)\), where
the latter notation means quaternionic matrices with real
part of the trace zero. The compact fixed subalgebras
of the involutions are \(\mathfrak{so}_n\) and
\(\mathfrak{sp}_{2m}\), by their defining matrix equations.

For orthogonal type, every compact involution is conjugate by a
compact automorphism to an involution represented by
\(A\in O(N,\mathbb R)\). By Theorem~\ref{res:J01}, the complex
inner subgroup is represented by \(SO(N,\mathbb C)\), and
in even dimension a real determinant \(-1\) representative
induces a diagram transposition. In \(D_4\) the diagram action
of an involution is either trivial or conjugate to this
transposition in the triality group \(S_3\). The permutations
of compact-normalized simple generators lift the required
diagram conjugacies, since they commute with compact conjugation
on those generators. Thus the compact involution, after this
conjugacy if needed, has a complex orthogonal representative \(A\).
As it preserves the compact real algebra, \(\overline A=cA\)
for a scalar \(c\), and orthogonality gives \(c=\pm1\).
The case \(c=-1\) would make \(A=iR\) with \(R\) real and
\(R^TR=-I\), which is impossible. Therefore \(A\) is real.
The condition \((\operatorname{Ad}A)^2=1\) is \(A^2=\pm I\).
If \(A^2=I\), the real orthogonal eigenspaces give
\(A=I_p\oplus(-I_q)\). The fixed real vector space of
the antiunitary \(A\) times coefficient conjugation is
\(\mathbb R^p\oplus i\mathbb R^q\); the original symmetric
form restricts to signature \((p,q)\). Hence the fixed
Lie algebra is \(\mathfrak{so}(p,q)\).
If \(A^2=-I\), \(N=2n\) and a real orthonormal basis
puts \(A\) into the standard complex-structure matrix
\(J\). This gives \(\mathfrak{so}^*(2n)\), the algebra
preserving both the complex symmetric form and the
quaternionic structure \(J\) times conjugation. Its
compact centralizer is \(\mathfrak u(n)\).
These are all orthogonal cases. A matrix with square
\(-I\) has determinant \(1\), so no missing
determinant-\(-1\) case occurs.

For symplectic type all diagram automorphisms are trivial.
An involution of the compact algebra is represented by
\(A\in Sp(n)\), modulo its central \(\{\pm I\}\);
thus \(A^2=\pm I\). For \(A^2=I\), the eigenspaces
are quaternionic orthogonal subspaces of dimensions
\(p,q\). They give
\(\mathfrak{sp}(p,q)\) and compact centralizer
\(\mathfrak{sp}(p)\oplus\mathfrak{sp}(q)\).
For \(A^2=-I\), a quaternionic orthonormal basis makes
\(A=iI\): choose a complex eigenvector of eigenvalue
\(i\), whose quaternionic mate has eigenvalue \(-i\),
and induct on its orthogonal complement.
The compact centralizer is \(U(n)\).
If \(j\) denotes the compact quaternionic structure on
\(\mathbb C^{2n}\), then the new antiunitary map \(Aj\)
squares to \(+I\). Its fixed real space has a real
nondegenerate alternating form, so a real symplectic
basis identifies the fixed algebra with
\(\mathfrak{sp}_{2n}(\mathbb R)\).
This also proves the last two entries of the table directly.

Within each family the matrix normal forms are unique up
to the stated exchange of eigenspaces. Distinct entries
are separated by the compact fixed algebra and its
center, except for the following actual root identifications.
One has \(B_1=C_1=A_1\), with compact
\(\mathfrak{so}(3)\cong\mathfrak{su}(2)\) and split
\(\mathfrak{so}(2,1)\cong\mathfrak{sl}_2(\mathbb R)\).
The isomorphism \(B_2=C_2\) identifies
\(\mathfrak{so}(5),\mathfrak{so}(4,1),\mathfrak{so}(3,2)\)
with \(\mathfrak{sp}(2),\mathfrak{sp}(1,1),
\mathfrak{sp}_4(\mathbb R)\).
For \(D_3=A_3\) the five forms correspond as
\[
\begin{array}{c|ccccc}
 &\mathfrak{so}(6)&\mathfrak{so}(5,1)&\mathfrak{so}(4,2)&
       \mathfrak{so}(3,3)&\mathfrak{so}^*(6)\\ \hline
 &\mathfrak{su}(4)&\mathfrak{su}^*(4)&\mathfrak{su}(2,2)&
       \mathfrak{sl}_4(\mathbb R)&\mathfrak{su}(3,1).
\end{array}
\]
These identifications follow by transporting their compact
involutions under the indicated root-system isomorphisms.
For \(D_4\) take the simple roots \(e_1-e_2,e_2-e_3,
e_3-e_4,e_3+e_4\) for rotations on four ordered planes.
The compact involution for \(\mathfrak{so}^*(8)\) has
angles \((\pi/2,\pi/2,\pi/2,\pi/2)\), so the eigenvalue
on the root line for the last simple root is \(-1\)
and that on the other three is \(1\). The involution for
\(\mathfrak{so}(6,2)\) has angles \((\pi,0,0,0)\), so
only the first simple root line has eigenvalue \(-1\).
Triality exchanges these two end roots. The eigenvalues
on the simple root lines determine those on all root
lines, and each automorphism acts trivially on the
common torus. This identifies the compact involutions
and therefore identifies their real forms.
The nonsimple case \(D_2=A_1+A_1\) has
\(\mathfrak{so}(4)=\mathfrak{su}(2)\oplus\mathfrak{su}(2)\),
\(\mathfrak{so}(2,2)=\mathfrak{sl}_2(\mathbb R)^{\oplus2}\),
and \(\mathfrak{so}(3,1)=\mathfrak{sl}_2(\mathbb C)\)
as a real algebra. The remaining mixed form is
\(\mathfrak{so}^*(4)\cong
\mathfrak{su}(2)\oplus\mathfrak{sl}_2(\mathbb R)\):
its compact fixed algebra is
\(\mathfrak{su}(2)\oplus\mathfrak u(1)\), so the two
\(A_1\) factors have respectively compact and split involutions.
The usual simple ranks are
\(A_{n-1}\) for \(n\geq2\), \(B_n,C_n\) for \(n\geq2\),
and \(D_n\) for \(n\geq4\), with these identifications understood.

Finally let \(\mathfrak l\) be real simple. Its complexification
is semisimple: a nonzero radical would be conjugation-stable
and would descend to a solvable ideal. Conjugation permutes
its complex simple ideals transitively; otherwise a union
of orbits descends to a proper real ideal. An involution
has orbits of size one or two, so either the complexification
is simple, or it is a pair \(\mathfrak a\oplus\mathfrak a\)
with factors exchanged. In the latter case the fixed algebra
is \(\mathfrak a\) with scalars restricted to \(\mathbb R\).
Conversely an ideal in such a real algebra complexifies
to a conjugation-stable ideal in the pair and hence is
zero or the whole algebra. This proves the second branch.
The matrix classification agrees with
\cite[Section 41.4]{etingof} and
\cite[Theorem 6.105]{knapp}.
\end{proof}

\section{Exceptional real forms by a finite root calculation}

The following conventions specify the finite calculation completely.
For \(E_r\), \(r=6,7,8\), label the horizontal chain
\(1-2-3-\cdots-(r-1)\), and attach the last vertex \(r\)
to vertex \(3\). Thus the three arms have lengths
\((2,r-4,1)\). For \(F_4\), vertices \(1,2\) are short,
vertices \(3,4\) long, and
\[
 A_{F_4}=\begin{pmatrix}
 2&-1&0&0\\-1&2&-2&0\\0&-1&2&-1\\0&0&-1&2
 \end{pmatrix},\qquad
 A_{G_2}=\begin{pmatrix}2&-3\\-1&2\end{pmatrix}.
\]
All \(E_r\) off-diagonal entries on an edge are \(-1\).
The convention remains \(a_{ij}=\alpha_j(h_i)\).

\begin{theorem}[The exceptional list]\label{res:J08}
The inner involutions, and the complexifications of their compact
fixed algebras, are exactly those in the following table.
\(S\) is the set of black vertices in the specified numbering,
and the orbit size refers to all \(2^r\) colorings.
The signature \(s\) is the difference of the two Killing indices.
\[
\begin{array}{c|c|r|l|r|r}
 D&S&\text{orbit}&\mathfrak k&\dim\mathfrak k&s\\ \hline
 G_2&\varnothing&1&G_2&14&-14\\
 G_2&\{1\}&3&A_1+A_1&6&2\\ \hline
 F_4&\varnothing&1&F_4&52&-52\\
 F_4&\{1\}&3&B_4&36&-20\\
 F_4&\{3\}&12&C_3+A_1&24&4\\ \hline
 E_6&\varnothing&1&E_6&78&-78\\
 E_6&\{1\}&27&D_5+\mathbb C&46&-14\\
 E_6&\{2\}&36&A_5+A_1&38&2\\ \hline
 E_7&\varnothing&1&E_7&133&-133\\
 E_7&\{1\}&63&D_6+A_1&69&-5\\
 E_7&\{4\}&36&A_7&63&7\\
 E_7&\{2,4\}&28&E_6+\mathbb C&79&-25\\ \hline
 E_8&\varnothing&1&E_8&248&-248\\
 E_8&\{1\}&135&D_8&120&8\\
 E_8&\{2\}&120&E_7+A_1&136&-24
\end{array}
\]
For \(E_6\) there are additionally two outer classes. The all-white
diagram involution has \(\mathfrak k=F_4\), dimension \(52\),
and signature \(-26\). Its nonzero fixed-vertex color class
has \(\mathfrak k=C_4\), dimension \(36\), and signature \(6\).
No other exceptional real forms occur.
The usual notation \(E_{r(s)}\), \(F_{4(s)}\), or \(G_{2(s)}\)
uses the signature column as its subscript.
\end{theorem}
\begin{proof}
For an inner datum \(b\in\mathbb F_2^r\), its fixed algebra is
\[
 \mathfrak h\oplus
 \bigoplus_{\substack{\alpha=\sum c_i\alpha_i\\
                     \sum b_i c_i\equiv0\ (2)}}\mathfrak g_\alpha.
 \tag{10.2}\label{eq:even-fixed-roots}
\]
Thus its dimension is \(r\) plus the number of even roots.
The even roots form a root system in their real span: they
are closed under reflections in their own roots, because
\(s_\alpha\beta=\beta-\beta(h_\alpha)\alpha\)
still has even parity. Its simple roots are exactly the
positive even roots that are not sums of two positive
even roots. The missing rank is the center dimension of
the fixed reductive algebra.

Here are explicit exhaustive procedures for both the orbits
and this root-system calculation. They use only integers.
Represent a root by its simple-root coefficient tuple \(c\).
Start with all unit tuples, close the set under
\[
 c_i\longmapsto c_i-\sum_j a_{ij}c_j
 \quad\text{with all other coordinates unchanged}.
 \tag{10.3}\label{eq:root-reflection-enumeration}
\]
This terminates at the known finite root system and obtains
every root: each root is in a Weyl orbit of a simple root,
as proved in Chapter~\ref{ch:root-systems}.
Independently, start with each as-yet-unseen binary tuple \(b\)
and close it under \eqref{eq:vogan-reflection}.
For each orbit retain its least integer representative
\(\sum_i2^{i-1}b_i\).
The resulting orbit sizes are respectively
\[
 (1,3),\quad(1,3,12),\quad(1,27,36),\quad
 (1,63,36,28),\quad(1,135,120).
\]
Their sums are \(2^2,2^4,2^6,2^7,2^8\).
For clarity, the core enumeration can be reproduced by
the following elementary pseudocode; \(\operatorname{closure}\)
adds images until no new tuple remains.
\[
\begin{aligned}
 \Phi&=\operatorname{closure}
       (\{e_1,\ldots,e_r\},\{s_1,\ldots,s_r\}),\\
 U&=\{0,1\}^r,\\
 \text{while }U\ne\varnothing:\quad
 b&=\operatorname{least}(U),\quad
 O=\operatorname{closure}(\{b\},\{s_1,\ldots,s_r\}),\\
 U&\leftarrow U\setminus O,\qquad
 \Phi_b=\{c\in\Phi:\sum b_i c_i\equiv0\pmod2\}.
\end{aligned}
\]
Here \(s_i\) in the first line is \eqref{eq:root-reflection-enumeration}
and in the orbit line is \eqref{eq:vogan-reflection}.
To check the fixed type without any Lie algebra identification
being assumed, take the positive elements of \(\Phi_b\),
delete every one that is a sum of two of them, and compute
their Cartan entries by
\[
 \frac{2c^T G d}{c^T G c},\qquad
 G=\operatorname{diag}(d_i)A,\quad
 (d_i)=
 \begin{cases}
 (1,3)&G_2,\\
 (1,1,2,2)&F_4,\\
 (1,\ldots,1)&E_r.
 \end{cases}
\]
This gives, in the table order, the root diagrams
\[
\begin{gathered}
 G_2,\ 2A_1;\qquad F_4,\ B_4,\ C_3+A_1;\\
 E_6,\ D_5,\ A_5+A_1;\qquad
 E_7,\ D_6+A_1,\ A_7,\ E_6;\qquad
 E_8,\ D_8,\ E_7+A_1 .
\end{gathered}
\]
The ranks in the two \(D_5\) and \(E_6\) rows are one less
than their ambient ranks, accounting for the central
\(\mathbb C\) in those rows. Counting the roots using the
same procedure gives ambient root counts
\(12,48,72,126,240\); the fixed dimensions then are exactly
those displayed. This procedure specifies every loop,
initial set, arithmetic operation, and stopping test in
the finite calculation, so the orbit and fixed-type
columns can be checked without using the real-form list
as input.
Each reflection is an actual conjugacy by
Proposition~\ref{res:J06}, so the calculation bounds the
number of inner forms from above. Different rows of one
type have different fixed dimensions, hence different
Killing signatures by Theorem~\ref{res:J04}; they cannot
be conjugate. Thus the bound is attained.

Only \(E_6\) has a nontrivial exceptional diagram involution.
It exchanges \(1\leftrightarrow5\), \(2\leftrightarrow4\)
and fixes \(3,6\). Reflection at a black fixed vertex
flips the other fixed vertex. Hence the four possible
colorings have the two classes \(00\) and
\(\{10,01,11\}\).
To identify the first, put
\[
 \beta_1=\tfrac12(\alpha_1+\alpha_5),\quad
 \beta_2=\tfrac12(\alpha_2+\alpha_4),\quad
 \beta_3=\alpha_3,\quad \beta_4=\alpha_6.
\]
For each paired orbit take the sums of its two \(e\)'s,
two \(f\)'s and two \(h\)'s; for each fixed vertex take
the individual generators. Each resulting
\((E_i,F_i,H_i)\) is an \(\mathfrak{sl}_2\) triple:
the roots in each paired orbit are nonadjacent, so all
cross brackets vanish. Direct evaluation on the simple
generators gives the Cartan matrix \(F_4\) in this order,
and \([F_i,E_j]=0\) for \(i\ne j\).
In the finite-dimensional adjoint module of the \(i\)-th
triple, \(E_j\) is therefore a lowest vector of weight
\(a_{ij}\). The rank-one classification gives
\((\operatorname{ad}E_i)^{1-a_{ij}}E_j=0\);
the negative relations follow in the same way.
All Serre relations have now been checked.
Theorem~\ref{res:F03} gives a nonzero homomorphism from
the simple \(F_4\) Serre algebra into the fixed algebra.
Its kernel is an ideal and cannot be the whole algebra,
so the homomorphism is injective and supplies a
\(52\)-dimensional fixed subalgebra. Direct application of
\eqref{eq:root-reflection-enumeration} shows that \(24\)
of the \(72\) roots of \(E_6\) are fixed by the diagram
involution. The other \(48\) roots form \(24\) pairs,
each contributing one fixed root vector. Thus the
fixed algebra has dimension at most \(4+24+24=52\).
Equality proves that it is precisely \(F_4\), and also
that the involution is \(+1\) on each fixed root line.

For the second class choose vertex \(3\) black. A fixed
root has an even coefficient at vertex \(3\) in exactly
eight of these \(24\) cases, by
\eqref{eq:root-reflection-enumeration}. The paired roots
still contribute \(24\) fixed vectors, so the fixed
algebra has dimension \(4+8+24=36\).
Its roots relative to the fixed Cartan are the projections
of the paired roots and these eight retained fixed roots.
The positive indecomposable roots in the \(\beta\) coordinates are
\[
 \beta_2,\quad\beta_1,\quad\beta_2+\beta_3,\quad\beta_4.
\]
Their squared lengths are \(1,1,1,2\), and their
successive inner products are
\(-\tfrac12,-\tfrac12,-1\), with all other inner
products zero. Their Cartan matrix is \(C_4\).
This verifies the fixed type as well as its dimension.
The two signatures are \(-26\) and \(6\), so these
two outer rows are distinct and are also distinct
from all inner rows.

Finally the split real algebra has signature equal to
the complex rank. Indeed its real Chevalley Cartan
has positive Killing form, and each real plane
\(\mathbb R e_\alpha+\mathbb R e_{-\alpha}\) contributes
one positive and one negative square. This identifies
the rows of signatures \(2,4,6,7,8\) as the split
forms. All remaining rows are constructed by their
explicit sign involutions. Combined with
Proposition~\ref{res:J05}, this proves existence,
exhaustiveness, and nonisomorphism of the exceptional
list. Its entries agree with
\cite[Section 42]{etingof} and
\cite[Theorem 6.105]{knapp}.
\end{proof}

\section{What the hypotheses prevent}

\begin{proposition}[Representative field and characteristic boundaries]
\label{res:J09}
The characteristic-zero, splitting, and finite-dimensionality
conditions in the preceding structure and representation
theorems cannot in general be removed.
\end{proposition}
\begin{proof}
\emph{A nonsplit real action.}
The one-dimensional abelian real algebra acting on
\(\mathbb R^2\) by
\(\left(\begin{smallmatrix}0&-1\\1&0\end{smallmatrix}\right)\)
has no invariant real line. After complexification it
has two eigenlines. Thus Lie's simultaneous-triangularization
conclusion needs a splitting hypothesis on the field.
Likewise a Cartan in \(\mathfrak{su}(2)\) has imaginary
nonzero root eigenvalues; its real algebra has the same
complex Dynkin type \(A_1\) as
\(\mathfrak{sl}_2(\mathbb R)\), but its Killing form is
negative definite rather than split.

\emph{An infinite-dimensional module for a semisimple algebra.}
In the Verma module \(M(0)\) of \(\mathfrak{sl}_2\), the
vectors \(f^jv\), \(j\geq0\), are independent by PBW.
The quotient sending \(v\) to \(1\) and all \(f^jv\),
\(j>0\), to zero is the trivial simple module.
It has no invariant section: a vector of weight zero
must be a multiple of \(v\), and \(fv\ne0\).
Hence this module is not completely reducible.
This does not contradict Weyl's theorem, which assumes
finite dimension.

\emph{Lie's theorem in positive characteristic.}
Let \(k\) have characteristic \(p>0\), and act on
\(V=k[t]/(t^p)\) by \(D=d/dt\), \(T=t\), and \(I\).
One has \([D,T]=I\), including on \(t^{p-1}\), and
\([D,I]=[T,I]=0\). Thus the three-dimensional
Heisenberg algebra is represented on \(V\).
If a nonzero vector were a common eigenvector for
\(D,T,I\), their commutator would kill it, whereas
\(I\) fixes it. There is no common eigenvector,
despite the acting algebra being nilpotent and
therefore solvable. Notice that \(I\) is not
nilpotent; Engel's theorem's hypothesis has not
been satisfied.

\emph{Cartan's Killing criterion in positive characteristic.}
For \(p>3\) let
\(W=\operatorname{Der}(k[t]/(t^p))\), with basis
\(e_i=t^{i+1}D\), \(-1\leq i\leq p-2\).
Its bracket is
\[
 [e_i,e_j]=(j-i)e_{i+j},
\]
with the term zero outside the indicated basis range.
This algebra is simple. Any nonzero ideal is stable
under \(\operatorname{ad}e_0\), whose eigenvalues
\(-1,0,\ldots,p-2\) are distinct in \(k\);
polynomial projections therefore put some \(e_i\)
in the ideal. Repeated bracketing with \(e_{-1}\)
puts \(e_{-1}\) in it, since the coefficients
\((i+1),i,\ldots,1\) are nonzero. Bracketing this
with the whole algebra gives
\(e_{-1},e_0,\ldots,e_{p-3}\). Finally
\([e_1,e_{p-3}]=(p-4)e_{p-2}\) supplies the last
vector. This proves simplicity.

Its Killing form is nevertheless zero.
Invariance under \(e_0\) says that \(B(e_i,e_j)\)
vanishes unless \(i+j=0\) modulo \(p\).
If the integer sum is \(p\), the product
\(\operatorname{ad}e_i\operatorname{ad}e_j\)
has zero diagonal: its final index has increased
by \(p\), beyond the basis range. The remaining
possibilities have integer sum zero. Here
\[
 B(e_0,e_0)=\sum_{a\in\mathbb F_p}a^2=0
 \qquad(p>3),
\]
and invariance with \([e_{-1},e_1]=2e_0\) gives
\(B(e_{-1},e_1)=-2B(e_0,e_0)=0\).
These exhaust the possibilities. Thus this simple
nonabelian algebra has zero Killing form.

For further positive-characteristic theory the appropriate
additional datum is a restricted operation \(x\mapsto x^{[p]}\):
it satisfies
\(\operatorname{ad}(x^{[p]})=(\operatorname{ad}x)^p\),
\((a x)^{[p]}=a^p x^{[p]}\), and the additive correction
given by the \(p\)-th-power commutator expansion.
For a derivation algebra the operation is the \(p\)-fold
composition, which is again a derivation by the binomial
formula in characteristic \(p\).
The restricted enveloping algebra is the quotient by
\(x^p-x^{[p]}\); these are additional relations, not the
ordinary characteristic-zero enveloping algebra.
The classification of restricted or modular simple
algebras, their modules, and the small-characteristic
exceptions belongs to the separate positive-characteristic
module in the scope map. The concrete counterexamples
above establish the boundaries used in this volume;
compare \cite[Section 15.4]{etingof},
\cite[Section 1]{premet}, and \cite[Sections A.1--A.2]{jantzen}.
\end{proof}


\backmatter
\begin{thebibliography}{99}
\addcontentsline{toc}{chapter}{Bibliography}
\bibitem{humphreys}
James E. Humphreys.
\emph{Introduction to Lie Algebras and Representation Theory}.
Graduate Texts in Mathematics 9, Springer, 1972, archived first-edition
reprint with the 1994 afterword.
\href{https://doi.org/10.1007/978-1-4612-6398-2}{DOI}.
\bibitem{etingof}
Pavel Etingof.
\emph{Lie Groups and Lie Algebras I and II}.
Combined MIT 18.745/18.755 lecture notes, archived OCW version,
Spring 2024 course resource.
\href{https://ocw.mit.edu/courses/18-755-lie-groups-and-lie-algebras-ii-spring-2024/resources/mit18_755_s24_lec_full_pdf/}{MIT OpenCourseWare}.
\bibitem{weibel}
Charles A. Weibel.
\emph{An Introduction to Homological Algebra}.
Cambridge Studies in Advanced Mathematics 38, Cambridge University
Press, archived 1995 paperback.
\bibitem{hall}
Brian C. Hall.
\emph{Lie Groups, Lie Algebras, and Representations: An Elementary
Introduction}. Second edition, corrected printing, Springer, 2015.
\href{https://doi.org/10.1007/978-3-319-13467-3}{DOI}.
\bibitem{knapp}
Anthony W. Knapp.
\emph{Lie Groups Beyond an Introduction}.
Progress in Mathematics 140, second edition, Birkh\"auser, 2002.
\bibitem{fialowski}
Alice Fialowski and Dmitry Fuchs.
\emph{Construction of Miniversal Deformations of Lie Algebras}.
Archived author version; sections on infinitesimal deformations and
obstructions.
\bibitem{jantzen}
Jens Carsten Jantzen.
\emph{Representations of Lie Algebras in Positive Characteristic}.
Archived author lecture notes.
\bibitem{premet}
Alexander Premet and Helmut Strade.
\emph{Classification of Finite Dimensional Simple Lie Algebras in
Prime Characteristics}. arXiv:math/0601380v2, 2006.
\bibitem{p103}
MathGenome.
\emph{Algebra Core} (P1-03), accepted notes and formalization,
MathGenome v0.21.0. Chapters 2, 4 and 10--11 provide module,
integral-extension and homological prerequisites.
\bibitem{multivariable}
MathGenome.
\emph{Multivariable Analysis} (D04), accepted notes and formalization,
MathGenome v0.21.0. Finite-dimensional inverse and implicit function
theorems.
\bibitem{homological}
MathGenome.
\emph{Advanced Homological Algebra}.
November 2026 notes, especially Chapter 6, results A06--A08,
the PBW paragraph and the coefficient long exact sequence.
\bibitem{finite-representations}
MathGenome.
\emph{Advanced Finite Group Representation Theory}.
November 2026 notes. Chapter 4 proves Schur--Weyl duality and
the classification of homogeneous polynomial representations.
\end{thebibliography}
\end{document}
